\documentclass[11pt,a4paper]{article}
\usepackage[margin=23mm,headheight=16pt]{geometry}
\usepackage{fontspec}
\setmainfont{NotoSerifCJK-Regular.ttc}[Path=/usr/share/fonts/opentype/noto/,FontIndex=2,BoldFont=NotoSerifCJK-Bold.ttc,ItalicFont=NotoSerifCJK-Regular.ttc,ItalicFeatures={FakeSlant=0.2}]
\setsansfont{NotoSansCJK-Regular.ttc}[Path=/usr/share/fonts/opentype/noto/,FontIndex=2,BoldFont=NotoSansCJK-Bold.ttc]
\XeTeXlinebreaklocale "zh"
\XeTeXlinebreakskip = 0pt plus 1pt
\usepackage{amsmath,amssymb,booktabs,longtable,array,graphicx,xcolor,tikz,fancyhdr}
\usetikzlibrary{arrows.meta,positioning,calc,shapes.geometric,fit,backgrounds,decorations.pathreplacing}
\usepackage{hyperref}
\definecolor{inkblue}{HTML}{173B58}
\hypersetup{colorlinks=true,linkcolor=inkblue,urlcolor=inkblue,pdftitle={移动与传感器网络课程讲义},pdfauthor={}}
\setlength{\parindent}{2em}
\setlength{\parskip}{5pt plus 1pt}
\setlength{\emergencystretch}{3em}
\linespread{1.18}
\renewcommand{\contentsname}{目录}
\renewcommand{\figurename}{图}
\renewcommand{\tablename}{表}
\pagestyle{fancy}\fancyhf{}\fancyhead[L]{\small\sffamily 移动与传感器网络}\fancyhead[R]{\small\sffamily 课程讲义}\fancyfoot[C]{\thepage}
\begin{document}
\begin{titlepage}
\vspace*{24mm}
{\sffamily\color{inkblue}\Huge 移动与传感器网络\par}
\vspace{6mm}{\Large Mobile and Sensor Networks\par}
\vspace{12mm}{\LARGE 课程讲义\par}
\vspace{16mm}
{\large 从无线信号到介质访问，再到多跳与低功耗网络\par}
\vfill
\noindent 本讲义依据2026年10月8日至16日已提供的19段课程记录及配套课件编写，沿授课顺序建立概念、模型与协议流程。正文采用中文叙述，术语首次出现附英文；示意图重新绘制，公式逐一说明使用条件。\par
\vspace{5mm}
\noindent 课程讲授中的简化模型用于解释基本机制；涉及具体标准时，保留必要的适用条件。课堂计算题与补充算例在题目处区分。材料定位列于文末，不插入连续讲解。\par
\vspace{8mm}{\small 2026年10月19日\par}
\end{titlepage}
\section*{课程结构与考试范围}
\addcontentsline{toc}{section}{课程结构与考试范围}
无线网络的性能由信号传输、共享信道的规则以及网络组织方式共同决定。本讲义依次讨论这三个层次：首先建立无线信号、信道与可靠传输的基础；随后说明Wi-Fi中多个节点如何有序发送；最后研究多跳网络和电池供电传感器网络如何在连通性、时延与能耗之间取得平衡。

课程考试说明规定笔试时长为2小时，满分100分，60分及格。题型由基础Wi-Fi信道接入计算题（30分）、传感器网络异步信道接入计算题（30分）以及两道课堂内容问答题（每题20分）组成。考试提供答题纸，需准备笔或铅笔、橡皮和简单计算器。

两类接入计算题是本讲义重点展开的部分。前者要求在帧交换过程中同时跟踪物理侦听、虚拟侦听、帧间隔与退避计数；后者要求区分随机等待、信道检查、接入失败和数据重传。问答题覆盖课堂内容，不能据题型比例将其他章节视为可忽略内容。

上述题型与分值依据课程\textit{Information about the final exam}说明页；并非依据讲义作者预测。正文中的补充例题用于说明机制，不冒充教师原题。\clearpage
\setcounter{tocdepth}{2}
\tableofcontents\clearpage
% Part 1: course order, 8 October through 9 October 07:59.
\section{从无线网络到一条无线链路}
\subsection{先看应用需要什么，再看网络怎样组成}
移动与传感器网络（Mobile and Sensor Networks）的起点不是某个通信公式，而是设备为什么需要连接。环境传感器周期性报告少量测量值，视频设备连续上传大量数据，车辆则需要在移动过程中交换信息。它们都使用无线网络（Wireless Network），却不要求相同的服务。视频重视吞吐量（Throughput），也就是单位时间内实际交付的数据量；及时控制重视时延（Latency），也就是信息从产生或发送到到达所经历的时间；电池供电的传感器还重视能耗（Energy Consumption）。可靠性（Reliability）关注信息能否按要求成功交付，例如用一定条件下成功交付的比例来衡量。设备数量、移动速度和覆盖范围，也会改变网络的设计。

因此，“通信速率越高，网络就越好”并不是普遍成立的判断。一台每天只上报少量数据的传感器，可能更需要几年不用换电池；一台高速移动的终端，需要在离开原来覆盖区以后继续保持业务。课程先介绍这些不同需求，是为了说明后面的调制、传播、错误控制和接入机制，都是解决具体约束的工具。

一个无线主机（Wireless Host）可以是手机、笔记本，也可以是没有屏幕的传感器。主机通过无线链路（Wireless Link）与另一个节点交换数据。基础设施网络（Infrastructure Network）中，终端通常连接到接入点（Access Point, AP）或基站（Base Station, BS），再由它们连接其他网络。接入点是无线局域网中常见的连接节点，基站是蜂窝网络中的常见连接节点，两者在这里承担相似的接入作用，但不意味着它们的协议和内部结构相同。

自组织网络（Ad Hoc Network）则允许设备直接通信，不要求每次交换都经过固定接入点。混合架构（Hybrid Architecture）把直接通信与基础设施接入结合起来，例如车辆之间直接交换消息，同时通过路侧设备进入更大的网络。网状网络（Mesh Network）进一步让节点相互转发数据：不是每个节点都有独立互联网出口，数据可以经过其他节点，到达具有出口的网关（Gateway）。这种架构把“能与邻居通信”和“能直接连接互联网”分成了两个问题。

\begin{figure}[htbp]
\centering
\begin{tikzpicture}[>=stealth, every node/.style={font=\small}]
\node[draw,rounded corners] (a) at (0,0) {终端};
\node[draw,rounded corners] (b) at (2.4,0) {接入点};
\node[draw,rounded corners] (c) at (5,0) {其他网络};
\draw[<->] (a)--(b); \draw[<->] (b)--(c);
\node at (2.5,-0.65) {基础设施接入};
\node[draw,circle] (d) at (0,-2) {甲};
\node[draw,circle] (e) at (2,-2) {乙};
\node[draw,circle] (f) at (4,-2) {丙};
\node[draw,rounded corners] (g) at (6,-2) {网关};
\draw[<->] (d)--(e); \draw[<->] (e)--(f); \draw[<->] (f)--(g);
\draw[<->] (e) to[bend left=28] (g);
\node at (3,-2.75) {网状转发：先到邻居，再寻找出口};
\end{tikzpicture}
\caption{连接方式的原创示意。线表示可用连接，不表示所有连接同时发送；网状网络中的一条路可能包含多次无线传输。}
\end{figure}

从覆盖尺度看，蓝牙（Bluetooth）常用于个人设备附近的短距离连接，Wi-Fi 常覆盖房间、楼宇或局部室外区域，蜂窝网络（Cellular Network）以基站服务较大的区域，卫星网络（Satellite Network）可以提供广域覆盖。这里比较的是课程中的典型尺度，不给每项技术设定固定不变的距离；天线、功率、环境和具体部署都会改变实际范围。游牧式移动（Nomadic Mobility）强调用户换了位置以后重新接入，例如从教室到宿舍后再使用网络；连续移动中的业务保持，则需要网络在移动期间支持接入变化。

覆盖范围与移动性（Mobility）是两个不同维度。较大的覆盖区通常让移动终端在同一个接入节点下停留更久，但移动支持还依赖网络是否提供合适的管理机制。切换（Handover/Handoff）是终端从一个接入节点转向另一个接入节点的过程。课程比较 Wi-Fi 与蜂窝网络时，既考虑覆盖区大小，也考虑切换管理，而不是单独由信号传播距离推断全部移动能力。至此，网络已经被拆成节点、链路与连接结构；接下来把目光缩到一条链路，研究信息究竟如何穿过空气。

\subsection{从比特到电磁波的完整路线}
比特（Bit）是取值为 0 或 1 的信息单位。空气中传播的不是打印出来的“0”和“1”，而是随时间变化的电磁波（Electromagnetic Wave）。发送设备必须把比特转换成可发射的信号，接收设备再根据信号判断发送端的选择。物理层（Physical Layer, PHY）负责这类实际信号的生成、传播接口和恢复工作。

调制（Modulation）就是让信号的某些特征随信息发生可辨认的变化。数字调制（Digital Modulation）的“数字”说明输入信息是离散比特，不说明实际无线波形只具有两个数值。即使输入只有 0 和 1，天线上产生的信号仍然是连续变化的物理量。接收端进行解调（Demodulation），依据事先约定的信号集合恢复比特。

课程用两个处理步骤解释发射过程。第一步把比特映射成承载信息的波形；第二步把该波形放到选定的无线工作频段。第二步在课件中称作模拟调制（Analog Modulation），其重点是载波上的频率搬移。这里不应把“模拟调制”误读为“用户的信息一定是模拟信息”。同一条发送链路完全可以先处理数字比特，再通过连续波形和射频电路发射出去。

\begin{figure}[htbp]
\centering
\begin{tikzpicture}[>=stealth,every node/.style={font=\small},x=1cm,y=1cm]
\node[draw,align=center] (a) at (0,0) {信息比特\\0101\ldots};
\node[draw,align=center] (b) at (2.3,0) {数字调制\\选择信息波形};
\node[draw,align=center] (c) at (5.1,0) {频率搬移\\进入工作频段};
\node[draw,align=center] (d) at (7.7,0) {发射天线};
\draw[->] (a)--(b);\draw[->] (b)--(c);\draw[->] (c)--(d);
\node[draw,align=center] (e) at (7.7,-1.7) {接收天线};
\node[draw,align=center] (f) at (5.1,-1.7) {搬回基带};
\node[draw,align=center] (g) at (2.3,-1.7) {数字解调\\判断发送信号};
\node[draw,align=center] (h) at (0,-1.7) {恢复比特};
\draw[->] (d)--node[right,align=left] {无线\\传播} (e);
\draw[->] (e)--(f);\draw[->] (f)--(g);\draw[->] (g)--(h);
\end{tikzpicture}
\caption{发送和接收是方向相反的处理链。图暂未加入纠错编码；后文将在比特与调制之间加入它。传播可能损坏信号，所以“恢复比特”并不保证每个比特都正确。}
\end{figure}

\section{频率、频谱与调制：信息怎样进入波形}
\subsection{三个可以改变的波形特征}
先用一个未承载复杂信息的正弦信号（Sinusoidal Signal）描述基本量：
\begin{equation}
s(t)=A\cos(2\pi f t+\phi).
\end{equation}
这里，$t$ 是时间，单位为秒；$A$ 是幅度（Amplitude），表示波形偏离零值的大小；$f$ 是频率（Frequency），单位为赫兹（Hertz, Hz），表示每秒振荡多少次；$\phi$ 是相位（Phase），单位为弧度，表示振荡在一个周期内的起始位置。$s(t)$ 可以理解为与发射信号对应的电压或电场分量，具体单位随所选物理量而定。这个表达式用于解释波形，不把幅度直接当作功率。

周期（Period）$T$ 是完成一次振荡所需的时间，因此
\begin{equation}
T=\frac{1}{f}.
\end{equation}
例如，频率越高，同一小段时间内能看到的振荡次数越多；相位改变则相当于沿时间轴移动波形。两个频率和幅度相同、相位相差 $\pi$ 弧度，即 $180^\circ$ 的正弦波，在任意时刻的取值互为相反数。

发送端可以分别改变这三个量。幅移键控（Amplitude Shift Keying, ASK）以不同幅度表示不同信息；频移键控（Frequency Shift Keying, FSK）以不同频率表示信息；相移键控（Phase Shift Keying, PSK）以不同相位表示信息。在最简单的二进制例子中，FSK 约定频率 $f_0$ 表示比特 0，频率 $f_1$ 表示比特 1。接收端观察一个约定好的时间区间，判断这一段波形更接近哪种频率，再输出对应比特。

关键是双方共享同一张映射表。某种波形本身没有天然的“0”或“1”含义；发送端和接收端必须对波形、观察时间以及映射规则达成一致。调制的过程不是把数字写在波上，而是在允许的信号集合中作出选择。

\subsection{载波频率不等于带宽，也不等于数据速率}
载波（Carrier）是用于承载并搬移信息的高频振荡，其频率称为载波频率（Carrier Frequency）$f_c$。基带（Baseband）是频率搬移之前、围绕较低频率表示信息的信号范围。把基带信号搬到 $f_c$ 附近后，它占用的是一段频率范围，而不仅是一个没有宽度的频率点。

频谱（Spectrum）描述信号包含哪些频率成分，以及这些成分的相对强弱。频带（Frequency Band）指一段频率区间。带宽（Bandwidth）描述这段区间的宽度；在用上下边界近似描述信号占用频带时，
\begin{equation}
B=f_H-f_L.
\end{equation}
其中，$f_L$ 为下边界频率，$f_H$ 为上边界频率，$B$ 为带宽，三者均以 Hz 为单位。实际信号的频带边界需要结合具体占用标准确定；这里采用理想化区间，是为了分清位置和宽度。

假设一个教学信道（Channel）从 $2.400$ GHz 延伸至 $2.420$ GHz，它的带宽是 $0.020$ GHz，即 $20$ MHz，而中间位置是 $2.410$ GHz。这个数值例子只说明单位换算和区间含义，不指定某个实际 Wi-Fi 信道的完整参数。把整个区间移到更高频率，若上下边界一起移动相同距离，带宽仍然不变。

\begin{figure}[htbp]
\centering
\begin{tikzpicture}[>=stealth,every node/.style={font=\small}]
\draw[->] (0,0)--(10,0) node[right] {频率};
\fill[blue!12] (2,0) rectangle (7,1.2);
\draw[thick] (2,0)--(2,1.2)--(7,1.2)--(7,0);
\draw[dashed] (4.5,0)--(4.5,1.6);
\node[below] at (2,0) {$f_L$};\node[below] at (7,0) {$f_H$};
\node[above] at (4.5,1.6) {$f_c$：频带所在位置};
\draw[<->] (2,-0.8)--node[below] {$B=f_H-f_L$：频带宽度} (7,-0.8);
\end{tikzpicture}
\caption{中心频率与带宽的区别。方框是用于理解占用区间的示意，不表示所有真实信号的频谱均为矩形。}
\end{figure}

比特率（Bit Rate）$R_b$ 则表示每秒发送多少比特，单位是 bit/s。GHz 与 Mbit/s 度量的对象不同：前者是振荡或频谱位置，后者是信息传送速度。不能看到“工作在 5 GHz”就推断“速率为 5 Gbit/s”。调制方式、可用带宽、编码开销、信道条件及接入等待共同影响最终的数据传输表现。

频谱资源需要协调使用。授权频谱（Licensed Spectrum）允许特定使用者在规定条件下使用某些频带；免许可频谱（Unlicensed Spectrum）也受使用规则约束，并不意味着没有发射功率限制或没有干扰。课程借频谱许可的成本强调：频率资源有限，网络设计不仅要让一对设备通信，还要让许多用户有效共享资源。

\subsection{从一个比特扩展到一个符号}
符号（Symbol）是在一个约定的符号时间内选择并发送的一种信号状态。若每个符号只在两个信号之间选择，它承载一个比特。若每次能可靠区分八种信号，就可以把三个比特作为一组，因为三个比特具有 $2^3=8$ 种组合。

一般地，对于具有 $M=2^m$ 个可区分状态、采用固定长度映射的调制，
\begin{equation}
m=\log_2 M,\qquad R_b=mR_s=\frac{m}{T_s}.
\end{equation}
这里，$M$ 是信号状态数；$m$ 是每符号承载的比特数；$R_s$ 是符号率（Symbol Rate），单位为符号每秒，也称波特（Baud）；$T_s=1/R_s$ 是符号持续时间，单位为秒；$R_b$ 是映射到这些符号上的比特率。这个式子尚未扣除编码冗余和协议开销，因此不能直接把它当作应用有效吞吐量。

例如，同样每秒发送一百万个符号，二进制调制每个符号带一个比特，对应 $1$ Mbit/s；16 个状态的调制每个符号带四个比特，对应 $4$ Mbit/s。提高速率的原因不是每个波形振荡了四倍次数，而是每次选择能够表达的可能性增加了。与之相伴的困难是：接收端必须更精细地区分这些可能性。

信号星座（Signal Constellation）把允许发送的信号状态画成若干点。点之间距离越近，受扰动后越容易混淆。课程先比较相位相差 $180^\circ$ 的两个状态，再把相位状态扩展到八个。若所有点都放在相同半径的圆上，状态数增加就意味着相邻相位间隔减小。

\subsection{QAM 为何更高效，也更依赖信道质量}
正交幅度调制（Quadrature Amplitude Modulation, QAM）同时利用幅度和相位的不同组合。常用星座图采用同相分量（In-phase Component, I）与正交分量（Quadrature Component, Q）作为两个坐标；这两个分量对应相差 $90^\circ$ 的参考振荡。点到原点的距离体现幅度，点相对坐标轴的角度体现相位。

16QAM 有 16 个允许状态，因此每符号承载 4 个比特；256QAM 有 256 个允许状态，因此每符号承载 8 个比特。这里的 16 和 256 是信号集合的大小，不是频率数，也不是发射天线数。

\begin{figure}[htbp]
\centering
\begin{tikzpicture}[>=stealth,scale=0.72,every node/.style={font=\small}]
\draw[->] (-4.3,0)--(4.6,0) node[right] {I};
\draw[->] (0,-4)--(0,4.4) node[above] {Q};
\foreach \x in {-3,-1,1,3}{\foreach \y in {-3,-1,1,3}{\fill[blue!70!black] (\x,\y) circle (0.09);}}
\draw[dashed,gray] (2,-3.6)--(2,3.6);
\draw[->,thick,red!70!black] (1,1)--(2.5,1.25);
\fill[red!70!black] (2.5,1.25) circle(0.08);
\node[align=left,anchor=west] at (4.4,2) {蓝点：允许发送的状态\\红点：一次受扰动的观测\\虚线：两列之间的示意边界};
\end{tikzpicture}
\caption{原创 16QAM 示意，坐标为相对值。发送点 $(1,1)$ 受到足够大的扰动后，观测点可能落到相邻状态一侧；图不指定具体比特标记或接收算法。}
\end{figure}

接收信号会受到衰减（Attenuation）、噪声（Noise）和干扰（Interference）的影响。衰减让信号变弱；噪声来自不可避免的随机扰动；干扰来自其他发射端的信号。接收端实际观察到的点不一定正好落在发送星座点上，只能根据观测结果判断最可能发送了哪个状态。如果在相同平均发射功率下把更多点放入有限的星座范围，点间间隔往往缩小，错误判断的风险随之上升。

所以，更大的星座不是无条件的改进。信道好时，可以使用更大的星座，提高每符号的信息量；信道差时，较少的状态更容易区分。所谓稳健性（Robustness）与效率的权衡，就是在当前信号质量下决定“每次选择表达多少信息”，而不是孤立地追求最大的调制阶数（Modulation Order）。

\section{无线传播：为什么收到的信号会变弱、变形}
\subsection{频率、波长与天线尺寸}
波长（Wavelength）$\lambda$ 是同一时刻波形沿传播方向重复一次的空间距离。对于传播速度为 $c$ 的电磁波，
\begin{equation}
\lambda=\frac{c}{f}.
\end{equation}
这里，$\lambda$ 的单位为米，$f$ 的单位为 Hz。在真空中 $c$ 约为 $3\times10^8$ m/s；对本节空气中传播的估算，使用这个近似值。频率越高，波长越短。例如 $2.4$ GHz 对应约 $0.125$ m，即 $12.5$ cm；这不是说所有 $2.4$ GHz 天线都必须长 $12.5$ cm，而是说明天线设计的尺寸尺度与波长相关。

天线（Antenna）把电路中的信号与空间中的电磁波联系起来。很多天线的特征尺寸是波长的某个比例，因此高频、短波长有利于在小型设备中布置天线，甚至布置多个天线。不过，高频传播也可能带来更大的损耗。课程在这一处强调的是两个方向相反的因素：小天线有利于移动设备，而更容易衰减的信号又不利于远距离覆盖。

\subsection{功率、能量与分贝要分别理解}
功率（Power）表示每单位时间传递的能量，单位为瓦特（Watt, W）；能量（Energy）的单位为焦耳（Joule, J）。如果一段持续时间 $\Delta t$ 内功率近似恒定为 $P$，则
\begin{equation}
E=P\Delta t.
\end{equation}
这里 $E$ 是该段信号的能量，$P$ 以 W 为单位，$\Delta t$ 以 s 为单位。若比特率为 $R_b$，每比特时间 $T_b=1/R_b$，并在相同平均功率意义下记账，则每比特能量（Energy per Bit）为
\begin{equation}
E_b=PT_b=\frac{P}{R_b}.
\end{equation}
因此，在功率固定时，提高每秒发送的比特数，会减少分摊到每个比特的能量。使用这个式子时，还必须说清 $R_b$ 指编码后的发送比特率还是用户信息比特率；两者不同，所定义的每比特能量也不同。本节先以进入调制器的比特率为准。

分贝（Decibel, dB）用于表示比值。两个正功率 $P_2$ 与 $P_1$ 的比值写成
\begin{equation}
L_{\mathrm{dB}}=10\log_{10}\!\left(\frac{P_2}{P_1}\right).
\end{equation}
如果 $P_2$ 是 $P_1$ 的十倍，结果为 $10$ dB；如果只有十分之一，结果为 $-10$ dB。这个量只表达两个功率的相对关系，不能脱离参考值充当绝对功率。表示路径损耗时通常将大功率放在分子，使损耗为正；表示接收相对发射的变化时则可能得到负数。计算前先确定比值的方向，比机械记忆符号更重要。

以 $1$ mW 为参考的分贝毫瓦（Decibel-milliwatt, dBm）用于表示绝对功率：
\begin{equation}
P_{\mathrm{dBm}}=10\log_{10}\!\left(\frac{P}{1\,\mathrm{mW}}\right).
\end{equation}
例如 $1$ mW 是 $0$ dBm，$10$ mW 是 $10$ dBm，$0.1$ mW 是 $-10$ dBm。负的 dBm 不表示负能量，只说明功率小于参考功率。还要注意两个 dBm 功率不能直接相加作为总功率；应先换回线性功率相加，再按需要转回 dBm。

\subsection{自由空间模型说明了什么}
自由空间传播（Free-space Propagation）先排除建筑、树木和反射路径，假设发射与接收之间存在无遮挡传播。各向同性天线（Isotropic Antenna）是假想的参考天线，它向所有空间方向均匀辐射。定向天线（Directional Antenna）则把更多能量集中到某些方向。天线增益（Antenna Gain）比较某方向与参考天线的辐射或接收能力，不是凭空产生能量。

在自由空间、远场（Far Field）、天线方向与极化匹配且其他损耗忽略的条件下，课件的接收功率关系可写成
\begin{equation}
P_r=P_tG_tG_r\left(\frac{\lambda}{4\pi d}\right)^2.
\end{equation}
其中，$P_t$ 是送入发射天线的功率，$P_r$ 是接收到的功率，二者用相同功率单位；$G_t$ 和 $G_r$ 分别为收发天线在相应方向上的线性增益，不能直接代入 dB 数值；$d$ 是传播距离，单位为米；$\lambda$ 是波长，单位也为米。远场条件要求不能把天线极近距离区域直接套入这个模型。

把不含天线增益的自由空间路径损耗（Free-space Path Loss）记为 $L_{\mathrm{fs}}$，有
\begin{equation}
L_{\mathrm{fs}}=\left(\frac{4\pi d}{\lambda}\right)^2
=\left(\frac{4\pi df}{c}\right)^2,
\qquad P_r=\frac{P_tG_tG_r}{L_{\mathrm{fs}}}.
\end{equation}
$L_{\mathrm{fs}}$ 是无量纲线性比值。固定频率、增益和发射功率时，距离增加为原来的两倍，接收功率降至四分之一；固定距离、增益和发射功率时，频率增加为两倍，在这一模型下接收功率也降至四分之一。两种变化都相当于约 $6$ dB 的额外路径损耗。

“衰减与频率平方成正比”必须连同这些固定条件一起理解。若换频率的同时也改变了天线增益、天线尺寸或传播环境，不能仅凭 $f^2$ 推断实际接收功率。这里的公式提供一个清晰基准，真实链路还要面对其他传播因素。

\subsection{从覆盖范围到空间复用}
高频带来的较小覆盖区，对单条远距离链路可能不利，对一个密集网络却可能有另一面。空间复用（Spatial Reuse）是让空间上相互影响足够小的连接重复使用相同频率资源。网络密集化（Network Densification）通过增加接入节点、缩小单节点服务区域，为更多局部连接创造复用机会。

课程用多个基站重复使用同一组频带来说明这一点。设可用频带数为 $L$，每个频带在给定条件下可服务 $U$ 个用户，那么一个基站的示意容量是 $LU$。如果有 $K$ 个基站能够在互不造成不可接受干扰的条件下重复使用这些频带，理想化总服务量可达到 $KLU$。这里 $L$、$U$、$K$ 均为计数，$LU$ 和 $KLU$ 只是课堂容量示意，不是实际网络的保证值。

增加基站并不自动把容量乘上基站数。相邻同频发射会相互干扰，覆盖、发射功率和资源使用都需要协调。这解释了为什么课程在讨论“高频损耗较大”之后，紧接着讨论“高频有利于更小覆盖和更多复用”：前者评价单条链路，后者评价整个区域内的总服务能力，两者的评价对象不同。

\subsection{障碍物、多径与衰落}
真实传播环境会吸收、散射或改变电磁波的传播方向。大气吸收（Atmospheric Absorption）与雨雾等因素能增加损耗；折射（Refraction）改变传播方向；反射（Reflection）使信号在表面改变行进方向；衍射（Diffraction）使障碍物边缘附近也可能出现传播分量。课程用建筑、地面、山体和树叶说明：传播条件不仅与两端距离有关，还与路径上的环境有关。

大气造成的附加损耗也与频率有关，并不在所有频率上平滑、等量地增加。课件中的雨、雾与水蒸气曲线用于说明，不同环境因素具有不同的频率响应，某些频率附近可能出现明显吸收。因此，选择工作频率时，除了自由空间损耗，还必须考虑实际路径所穿过的介质。课程举树叶季节变化的例子，也是在说明：两端设备不动，不代表传播环境恒定。

地球曲率还会限制视距传播（Line-of-Sight Propagation）。在课件采用的等效地球半径近似中，一根高度为 $h$ 的天线到无线电地平线（Radio Horizon）的距离近似为
\begin{equation}
d_{\max}\simeq\sqrt{2KR_T h}.
\end{equation}
这里，$R_T$ 是地球半径，课件取约 $6370$ km；$K$ 是描述大气折射影响的无量纲等效地球半径系数；$h$ 是天线相对当地地面的高度；$d_{\max}$ 为单端天线到地平线的近似距离。$R_T$ 与 $h$ 必须采用相同长度单位，结果才具有同一单位。此近似使用 $h\ll KR_T$，并不包含山体遮挡或接收灵敏度的限制。它描述几何与折射意义上的视距，不保证在该距离内功率一定足够通信。两端天线都有明显高度时，还需要考虑两端各自的地平线范围，不能把单端式直接当作所有链路的最大距离。

多径传播（Multipath Propagation）意味着接收端收到了同一个发射信号沿不同路径到达的多个副本。不同路径具有不同长度，因此副本到达的时间、幅度和相位可能不同。它们在接收端叠加，既可能增强，也可能削弱总信号。

\begin{figure}[htbp]
\centering
\begin{tikzpicture}[>=stealth,every node/.style={font=\small}]
\node[draw,rounded corners] (tx) at (0,0) {发送端};
\node[draw,rounded corners] (rx) at (8,0) {接收端};
\draw[->,thick] (tx)--node[below] {直接路径 $d_0$} (rx);
\draw[thick] (2.7,2)--(5.3,2);
\node[above] at (4,2) {反射表面};
\draw[->,blue!70!black] (tx)--(4,2)--node[right] {较长路径 $d_1$} (rx);
\node[align=center] at (4,-1.2) {同一信号的副本先后到达\\接收端把它们叠加，而不是分别“挑选”一条};
\end{tikzpicture}
\caption{原创两径示意。真实环境可存在更多路径；图中反射只用于解释不同路径长度与接收叠加。}
\end{figure}

若只考虑路径长度造成的差异，路径差 $\Delta d=d_1-d_0$ 对应时延差 $\Delta\tau=\Delta d/c$，以及载波相位差
\begin{equation}
\Delta\phi=2\pi f\Delta\tau=\frac{2\pi\Delta d}{\lambda}.
\end{equation}
这里 $d_0,d_1,\Delta d,\lambda$ 均以米为单位，$\Delta\tau$ 以秒为单位，$\Delta\phi$ 以弧度为单位；相位按 $2\pi$ 周期重复。真实反射还可能引入额外相位变化，因此该式只提取“多走了一段路”的影响，不能把路径差本身直接当成以弧度计的相位差。

在两个副本幅度相同、频率相同的简化情形下，同相到达会相长叠加（Constructive Interference）；相位相差 $\pi$ 会相消叠加（Destructive Interference）。只有幅度也相同的理想情形才会完全抵消。一般情形只是部分增强或部分减弱。接收位置稍有改变，路径差就可能改变，所以“离接入点更近”不保证每一个局部位置的信号都更强。

衰落（Fading）描述接收信号随位置或时间变化的起伏。阴影衰落（Shadowing）与建筑、地形等遮挡有关，常表现为较缓慢的变化；多径引起的快速起伏在课程中称为快衰落（Fast Fading）。慢衰落（Slow Fading）与快衰落是相对于观察或通信过程的时间尺度而言的。静止设备也可能因周围人或障碍物移动而经历变化，移动设备则更容易不断进入不同的传播组合。课程此处建立物理直觉，不要求进一步计算详细的统计衰落模型。

\subsection{用 SNR 和 SINR 把传播联系到解调}
接收功率大并不必然意味着容易解调：如果噪声和干扰也很强，有用信号仍然难以分辨。信噪比（Signal-to-Noise Ratio, SNR）比较有用接收功率与噪声功率：
\begin{equation}
\mathrm{SNR}=\frac{P_r}{N}=\frac{P_r}{N_0B}.
\end{equation}
这里 $P_r$ 是经历传播损耗后到达接收端的有用功率；$N$ 是接收带宽内总噪声功率；$N_0$ 是噪声功率谱密度（Noise Power Spectral Density），单位为 W/Hz；$B$ 是所考虑的接收带宽，单位为 Hz。式子 $N=N_0B$ 假设带内噪声谱密度近似平坦，并采用课件中与该带宽相配的谱密度约定。所有功率都应以线性单位代入。

如果同时存在其他发射端，就使用信干噪比（Signal-to-Interference-plus-Noise Ratio, SINR）：
\begin{equation}
\mathrm{SINR}=\frac{P_r}{I+N_0B}.
\end{equation}
其中 $I$ 是接收端看到的总干扰功率，也就是各干扰信号经历其各自传播损耗之后的功率总和，不是邻居设备标称的发射功率总和。分母里的噪声与干扰必须使用相同线性功率单位相加。

例如，在统一的任意线性功率单位下，$P_r=100$、$N=1$ 时，SNR 为 $100$，即 $20$ dB；若新增干扰 $I=9$，SINR 变为 $100/(9+1)=10$，即 $10$ dB。有用信号并未变弱，接收质量却显著下降。这正是商场中多个接入点以及共用频段的不同无线技术可能互相影响的原因。

误比特率（Bit Error Rate, BER）描述接收比特判错的比例或相应概率。对于确定的调制、编码和信道条件，较差的 SNR 或 SINR 通常意味着更高的误比特率。但不存在适用于所有调制方式的单一“正确接收阈值”；高阶星座需要更好的区分条件。到这里，传播与调制形成闭环：传播决定观测信号有多可靠，调制决定在该可靠程度下能承载多少信息。

\section{错误控制：增加冗余，还是重新发送}
\subsection{检错与纠错解决不同的问题}
错误检测（Error Detection）判断数据中是否存在可检测的错误；错误纠正（Error Correction）进一步恢复原来的数据。发现“答案不对”不等于知道“应该改哪一位”。当接收端只有受损信息本身时，通常缺少辨别依据，因此发送端需要加入冗余（Redundancy）。冗余不是独立的新用户信息，而是由原有信息按规则生成的附加约束。

前向纠错（Forward Error Correction, FEC）通过信道编码（Channel Coding）预先加入冗余，使接收端在能力范围内直接修复错误。自动重传请求（Automatic Repeat reQuest, ARQ）则主要用检错结果决定是否重新传送。两种思路都付出代价：FEC 要传额外比特，ARQ 可能需要额外发送和等待。实际机制可以组合使用，但课程在这里先分别理解。

\subsection{分组码与码率：冗余究竟占多少}
分组码（Block Code）把 $k$ 个信息比特编码成 $n$ 个发送比特，形成码字（Codeword）。在课件采用的信息位加冗余位的表示中，其中 $n-k$ 个是冗余比特。码率（Code Rate）定义为
\begin{equation}
r_c=\frac{k}{n},\qquad 0<r_c\leq 1.
\end{equation}
这里 $k$ 和 $n$ 是每个块的比特数，$r_c$ 是无量纲比例。码率越低，单位发送比特中用于承载新信息的比例越小，但更多冗余可为保护信息提供空间。它不意味着“码率低就一定在所有条件下更好”：实际检测和纠正能力还取决于码的具体结构。

若编码后的发送比特率为 $R_{\mathrm{coded}}$，在忽略其他开销时，信息比特率为
\begin{equation}
R_{\mathrm{info}}=r_cR_{\mathrm{coded}}=r_cmR_s.
\end{equation}
$R_{\mathrm{info}}$ 和 $R_{\mathrm{coded}}$ 均以 bit/s 为单位；$m$ 是前文定义的每调制符号比特数，$R_s$ 为符号率。这个式子把调制与编码接到一起：大星座提高每符号可发送的比特数，编码冗余则决定其中有多少属于用户信息。协议首部、确认与重传仍未计入，所以应用真正得到的吞吐量还可能更低。

\subsection{重复码：最直观的纠错过程}
重复码（Repetition Code）把同一信息发送多次。三重复码把信息位 0 编为 000，把信息位 1 编为 111，因此 $k=1$、$n=3$、$r_c=1/3$。接收端对这三个位置作多数判决（Majority Decision）：至少两个位置为 0 就判为 0，至少两个位置为 1 就判为 1。

例如，发送 111，接收成 101，虽然中间一位出错，仍有两票支持 1，因此能够恢复原信息。若接收成 001，则已有两位出错，多数判决会错误地输出 0。因此，这个例子每组三位能保证纠正一个翻转错误，不能保证纠正两个错误。“重复三次”提升可靠性的代价，是用三个发送比特承载一个信息比特。

重复码的价值主要在于说明编码如何利用附加结构。接收端知道合法发送结果只能是 000 或 111，才有依据把 101 解释成 111 的受损版本；若任意三位串都合法，就无法据此纠正。更复杂的编码同样依靠合法信号集合受到约束这一原理，并不要求简单地逐位重复。

\subsection{奇偶校验：为什么能发现，却未必能修复}
奇偶校验（Parity Check）给一组信息比特增加一个校验位（Parity Bit）。偶校验要求包含校验位在内的整个码字中，1 的数量为偶数。比如信息串 1011 含有三个 1，所以附加校验位 1，得到 10111，整个码字有四个 1。

如果传输中恰好一位翻转，1 的数量的奇偶性一定变化，接收端就能发现不满足约束。但它并不知道哪一位翻转，因为每个位置的单比特错误都会造成相同的“奇偶性不对”。因此，单个奇偶位可以检测任意奇数个翻转错误，却不能仅靠这个结果定位并纠正一个错误；偶数个翻转错误可能保持原来的奇偶性，从而漏检。

校验位本身同样是通过信道发送的，也可能出错。课堂中为了展示比较步骤，有时先假设收到的校验位正确；完整的偶校验规则则是对收到的整个码字检查奇偶性。无论翻转发生在信息位还是校验位，一个翻转都会破坏整个码字的偶校验约束。

二维奇偶校验（Two-dimensional Parity Check）把信息位排列成矩阵，分别给每行和每列增加校验。若只有一个信息位翻转，就会恰好出现一行不通过、一列不通过，两者交点定位出错位置。

\begin{center}
\begin{tabular}{c|ccc|c}
 & 第1列 & 第2列 & 第3列 & 行偶校验位\\\hline
第1行 & 1 & 0 & 1 & 0\\
第2行 & 0 & 1 & 1 & 0\\
第3行 & 1 & 1 & 0 & 0\\\hline
列偶校验位 & 0 & 0 & 0 & \\
\end{tabular}
\end{center}
在这个原创例子中，假设第 2 行第 2 列由 1 翻成 0，且校验位都正确，那么第 2 行与第 2 列的 1 的数量都变为奇数，其余行列仍通过。接收端找到交点并翻转该位，完成单个信息位的纠正。这个简化示例未在右下角加入校验位的交叉保护，因此不据此声称能处理任意校验位错误，也不声称能纠正任意多个错误。

卷积码（Convolutional Code）是课件提到的另一类结构。信息比特依次进入移位寄存器（Shift Register），输出由当前输入及保留的历史状态共同决定。课件以每次输入 $k$ 位、输出 $n$ 位描述其码率 $k/n$。此处掌握它与独立块处理的区别即可；本段课堂没有展开具体译码算法，因此不引入额外计算程序。

\subsection{从比特串到数据包与 CRC}
数据包（Packet）把一组比特作为一个完整传送单位。首部（Header）携带地址等控制信息，有效载荷（Payload）携带需要交付的用户内容。信件的信封与正文可以帮助理解这种分工：地址使网络知道如何处理这份数据，而正文才是用户主要想交付的内容。

循环冗余校验（Cyclic Redundancy Check, CRC）根据受保护的数据计算检错值。课堂把它描述为有效载荷的简短“摘要”：发送端根据原载荷生成 CRC，并把它与数据一起发出；接收端根据实际收到的载荷重新计算 CRC，然后比较结果。这里的摘要只是功能类比，不表示它能像正文一样概括语义。

课程为解释流程，把 CRC 放在首部，并假定首部能够正确接收。这个安排是教学模型，不意味着所有协议的 CRC 都位于首部，也不意味着真实首部永远不会出错。理解这一前提后，接收过程就可以按次序读清楚：先提取发送端提供的 CRC，再接收载荷，再用相同规则计算本地 CRC，最后比较两个值。

若两个值不同，说明数据没有通过检查，不能按成功接收处理。若两个值相同，接收端按协议视为检查通过。但有限长度的检错值不能为任意长度数据提供绝对无差错证明，仍可能存在未被检测的错误模式。因此，教材将课堂口语中的“正确收到”精确写为“通过检错，按成功处理”。课程没有进行 CRC 多项式除法，本节也不加入这类运算。

\subsection{停止等待 ARQ：发送、确认、超时、重传}
停止等待（Stop-and-Wait）是本段课程实际讲解的 ARQ 方法。发送端一次发送一个数据包，然后等待接收端的确认（Acknowledgment, ACK）。ACK 是一个控制包，表达前一个包已通过接收端检查。如果等待期间得到了有效 ACK，发送端才继续发送下一个新包；如果等到超时（Timeout）仍没有有效 ACK，就重传（Retransmission）原来的包。

完整过程可以按两个参与者的动作展开。发送端先保留待发送数据包的副本，为载荷生成 CRC，发送包并启动计时。接收端解调后计算本地 CRC；在课程采用的流程中，校验通过便返回 ACK，校验失败便不返回 ACK。发送端在计时期限内收到 ACK，结束当前包的等待并处理下一个包；如果计时到期，则再次发送同一个包，而不是直接跳过它发送新包。

\begin{figure}[htbp]
\centering
\begin{tikzpicture}[>=stealth,every node/.style={font=\small}]
\node at (0,0.4) {发送端};\node at (7,0.4) {接收端};
\draw[->] (0,0)--(0,-6.5);\draw[->] (7,0)--(7,-6.5);
\draw[->,thick] (0,-0.4)--node[above,sloped] {包 A：首次发送} (7,-1.2);
\node[right,align=left] at (7,-1.3) {CRC 失败\\不发 ACK};
\draw[decorate,decoration={brace,amplitude=5pt}] (-0.3,-2.5)--(-0.3,-0.4);
\node[left,align=center] at (-0.65,-1.5) {等待\\超时};
\draw[->,thick] (0,-2.7)--node[above,sloped] {包 A：重传同一包} (7,-3.5);
\node[right] at (7,-3.5) {CRC 通过};
\draw[->,blue!70!black] (7,-3.8)--node[above,sloped] {ACK(A)} (0,-4.6);
\draw[->,thick] (0,-5.0)--node[above,sloped] {包 B：收到确认后才发送} (7,-5.8);
\node[left] at (0,-6.4) {时间向下};
\end{tikzpicture}
\caption{原创停止等待时序图。斜线表示传送需要时间，不按真实比例绘制。第一次失败后，发送端因没有收到确认而超时重传。}
\end{figure}

“自动重传请求”不要求接收端一定发出一个写着“请重传”的消息。在这里，不返回 ACK 与发送端的超时机制结合，就触发了重传。另一方面，没有收到 ACK 也不能严格证明数据包一定没到达：确认本身可能丢失。因此，课件还指出包需要序列号（Sequence Number），以区分新包和重复包。序列号让接收端有依据识别重传，而不把重复收到的同一份内容当成两份新内容。课堂没有展开其他 ARQ 协议，本节也停留在停止等待机制。

发送端通常设有最大重试次数。到达上限仍未完成确认，就丢弃当前包或向上层报告失败，不会无穷次占用链路。这也说明 ARQ 不是“保证所有包必定到达”的魔法：它是在时间、能量和可靠性之间作取舍。链路条件很差时，重传次数增加，不仅消耗发射能量，也拉长交付时间。

停止等待名称中的“等待”揭示了它的主要效率限制：在前一个包得到确认之前，发送端不发送下一个新包。一个成功周期包括数据包发送、信号传播、接收处理以及 ACK 返回等过程。即使物理层瞬间能达到较高比特率，应用仍然可能因为等待和重传而获得较低吞吐量。本段课程重点是正确读懂这个过程，不要求推导额外的停止等待效率公式。

\section{从通信基础走向介质访问控制}
经过前面各节，一条无线传输链路已经具有完整含义：用户数据先形成比特，必要时加入冗余，经调制选出波形，再搬到工作频段，由天线发射；信号经历传播损耗、多径、噪声和干扰后，接收端解调、检错或纠错，必要时通过确认和重传完成可靠交付。

然而，这还没有回答一个网络问题：如果许多设备都想在同一无线信道上发送，谁应当先发？介质访问控制（Medium Access Control, MAC）负责协调对共享通信介质的使用。物理层解决“怎样把数据变成可发送、可接收的信号”，MAC 层进一步解决“什么时候允许一个设备发这个包”。即使某个包已经准备好，设备也可能需要等待合适的发送时机。

课程在 10 月 9 日由此转入 Wi-Fi 与 IEEE 802.11。无线局域网（Wireless Local Area Network, WLAN）提供局部区域内的无线连接；IEEE 802.11 是课程用来解释 Wi-Fi 物理层和 MAC 层工作的标准体系。应用（Application）产生数据之后，数据沿协议栈（Protocol Stack）向下传递，到达 MAC 层，最后交给物理层进行实际发送。接入点收到后，再把数据转入后续网络。

这里要保留两条相互关联、但不能混同的主线。第一条是信号质量：衰减和干扰影响判决，调制与编码决定如何适应这种质量。第二条是共享资源：多个设备竞争同一个介质，接入规则决定谁在何时使用它。后续 Wi-Fi 章节将沿第二条主线展开，但每次讨论碰撞、等待、确认与重传时，都要回到第一条主线，理解它们最终如何影响真实的无线发送。

\section{Wi-Fi 网络的组织与介质访问问题}
\label{ch:wifi-organization}

无线链路首先要把比特变成可以传播的信号，但一个可用的无线网络还必须解决另外两件事：设备如何找到并加入网络，以及已经加入的设备如何共享信道。本章先建立 Wi-Fi 的网络结构与收发过程，再把注意力收拢到介质访问控制。后续的退避、确认和信道预留，都是在这个共同模型上逐步增加的机制。

\subsection{从一条无线链路到一张网络}

\subsubsection{物理层与 MAC 层各自负责什么}

IEEE 802.11 主要规定物理层（Physical Layer，PHY）与介质访问控制层（Medium Access Control，MAC）。物理层负责调制、无线发射、接收与解调；MAC 层则组织帧、决定何时允许发送，并在本章讨论的单播机制中检查接收结果和实施确认。逻辑链路控制层（Logical Link Control，LLC）以及更高的网络层、传输层和应用层，不应全部归入 Wi-Fi 标准本身。

以手机上传照片为例，应用产生的数据经过上层处理后到达 MAC 层。此时有数据可发，并不意味着天线立即发射。MAC 必须先取得一次使用共享介质的机会，才将帧交给物理层。接收端从无线信号恢复比特，再由 MAC 检查帧是否正确；正确的数据向上交付，适用时还要生成确认帧，经本机物理层发回发送端。这一往返过程把前面学过的调制、差错检测与停止等待连接起来。

家用设备常把接入点和路由器装在同一个盒子里，但两者是不同的逻辑功能。接入点（Access Point，AP）提供无线接入；路由器负责网络层的分组转发。手机与 AP 之间的一跳采用 Wi-Fi，并不表示从 AP 到互联网服务器的整条路径都采用 Wi-Fi。数据进入后方网络后，可以经有线链路和多台路由器继续前进。

\subsubsection{BSS、IBSS 与 ESS}

站点（Station，STA）是参与 802.11 通信的设备。基本服务集（Basic Service Set，BSS）表示采用同一协调机制共享无线介质的一组站点；基本服务区（Basic Service Area，BSA）则强调其无线覆盖的空间范围。集合与覆盖区域相关，但不是同一个概念。

在基础设施模式（Infrastructure Mode）下，AP 是 BSS 的中心接入节点。在本课程采用的基本模型中，站点的数据通过 AP 转发；即便两个终端能够听到彼此的无线信号，也不因此直接交换普通用户数据。\textbf{能听到一个节点、与它关联、把它作为当前帧的接收方，是三个不同的关系。}这一区分将在理解旁听、隐藏终端与 NAV 时反复用到。

独立基本服务集（Independent BSS，IBSS）对应不依赖 AP 的自组织模式（Ad Hoc Mode）。站点直接通信，适用于临时组网或缺少固定设施的环境。这里的独立 BSS 并不自动等于具备完整多跳路由能力的网络；直接无线接入与多跳路径组织仍是不同问题。

多个基础设施 BSS 可以通过分布系统（Distribution System，DS）互联，形成扩展服务集（Extended Service Set，ESS）。校园里分散的 AP 因此可以共同提供较大的服务范围。DS 是连接 BSS 的逻辑系统，不必把它限定成某一种特定电缆。用户从一个 BSS 移动到另一个 BSS 时，需要改变关联关系；后方系统也必须能够把后续数据送到新的接入位置。

\begin{figure}[htbp]
\centering
\begin{tikzpicture}[x=1cm,y=1cm, every node/.style={font=\small},>=stealth]
\draw[dashed,rounded corners] (-0.8,-0.7) rectangle (4.1,2.1);
\draw[dashed,rounded corners] (5.4,-0.7) rectangle (10.3,2.1);
\node at (1.6,1.8) {BSS 1}; \node at (7.8,1.8) {BSS 2};
\node[draw,rounded corners] (s1) at (0,0.7) {STA A};
\node[draw,rounded corners] (s2) at (0,-0.1) {STA B};
\node[draw,fill=blue!10] (a1) at (3,0.3) {AP 1};
\node[draw,fill=blue!10] (a2) at (6.5,0.3) {AP 2};
\node[draw,rounded corners] (s3) at (9.2,0.7) {STA C};
\node[draw,rounded corners] (s4) at (9.2,-0.1) {STA D};
\draw[<->] (s1)--(a1); \draw[<->] (s2)--(a1);
\draw[<->] (a2)--(s3); \draw[<->] (a2)--(s4);
\draw[<->,thick] (a1)--node[above]{DS}(a2);
\node at (4.75,-1.2) {多个 BSS 经 DS 互联，构成 ESS};
\end{tikzpicture}
\caption{基础设施网络的逻辑结构。连线表示通信或互联关系，并不表示所有无线节点都能彼此侦听。}
\label{fig:wifi-ess}
\end{figure}

\subsection{加入网络：发现、认证与关联}

服务集标识符（Service Set Identifier，SSID）是用户通常看到的网络名称；基本服务集标识符（BSS Identifier，BSSID）标识具体的 BSS。在常见基础设施网络中，BSSID通常与对应 AP 无线接口的 MAC 标识相关。同一网络名称可以由多个 BSS 提供，因此手机列表里的一个名称不必只对应一台 AP。这两个标识用于区分“想加入哪张网络”和“当前选择哪个接入实例”。

站点使用一个 BSS 前，必须先知道它在哪里、使用哪个信道，并完成相应的加入过程。课程将这一过程分成扫描、认证和关联三步。这个划分说明的是功能次序；实际安全接入中，认证与密钥建立的细节会随所采用的安全机制而不同，不能把所有认证都简化成同一种用户名和密码交换。

\subsubsection{扫描为何需要逐信道进行}

同一 BSS 的基本通信发生在选定的无线信道上。尚未加入的站点不知道附近 AP 使用哪些信道，因此必须收集候选网络的信息。被动扫描（Passive Scanning）让站点依次驻留在各信道，等待 AP 周期发送的信标帧（Beacon）。信标是管理帧，用来宣告网络的存在、传递参数，并提供时间同步参考。接收信标时，站点还可以测量信号质量，形成候选 AP 列表。

主动扫描（Active Scanning）由站点发送探测请求（Probe Request），再收集探测响应（Probe Response）。它不必总等待下一次周期信标，因此可以缩短发现过程，但会引入额外报文。站点记住经常使用的网络，并优先探测这些网络，是减少重复发现开销的一种思路。不过，若探测请求显式携带常用网络标识，旁听者可能推断设备曾经连接哪些网络。这里的教学重点是：加快发现与减少信息泄露之间存在取舍，不能只评价接入速度。

\subsubsection{关联把“发现 AP”变成“由该 AP 服务”}

认证（Authentication）用于建立相应的身份或接入资格关系；关联（Association）则让站点与 AP 建立服务关系。关联请求和关联响应交换设备能力、支持速率等信息。对于移动站点，重关联（Reassociation）还涉及新的服务位置以及原有连接的信息。

BSS 的基本速率集合（Basic Rate Set）表示参与该网络所需支持的一组速率要求。它不等价于“所有数据帧都只能使用一个固定的最低速率”。数据帧可以依据链路条件采用不同的可用速率；控制和管理帧则必须满足相应的可解码性与速率选择规则。理解这一点，可以避免把广播覆盖要求、关联能力检查和单播速率自适应混成一件事。

\begin{figure}[htbp]
\centering
\begin{tikzpicture}[every node/.style={font=\small},>=stealth]
\node[draw,rounded corners,text width=2.4cm,align=center] (a) at (0,0) {扫描\\发现信道与 AP};
\node[draw,rounded corners,text width=2.4cm,align=center] (b) at (3.8,0) {认证\\建立接入资格};
\node[draw,rounded corners,text width=2.4cm,align=center] (c) at (7.6,0) {关联\\交换能力与参数};
\node[draw,rounded corners,text width=2.4cm,align=center] (d) at (7.6,-1.8) {数据通信\\每次发送仍需接入};
\draw[->] (a)--(b); \draw[->] (b)--(c); \draw[->] (c)--(d);
\end{tikzpicture}
\caption{加入网络与使用信道是两个层次。关联成功不会使站点永久获得无线信道的独占权。}
\end{figure}

\subsection{物理层回顾：频率资源与可用速率}

\subsubsection{频段、信道与相邻网络的干扰}

频段（Frequency Band）是一整段可用频率范围，信道（Channel）是其中用于一次无线通信的一部分资源。Wi-Fi 从早期主要使用 2.4 GHz，逐步扩展到 5 GHz 和 6 GHz，使部署者能够在更多可用频率资源之间选择。这里的 GHz 表示载波所在的频率数量级，不是“每秒能发送多少吉比特”。频率、带宽与数据速率相关，却是不同物理量。

课程用免许可频谱（Unlicensed Spectrum）解释 Wi-Fi 的易部署性：使用者通常无需像蜂窝运营商那样取得一块专属频谱许可，但仍须遵守所在地区的频段、功率及设备要求。免许可并不意味着频谱完全没有管理，也不意味着同频设备互不干扰。

以传统 2.4 GHz 信道规划为例，相邻信道中心可以相隔 5 MHz，而早期 DSSS 信号的名义信道宽度约为 22 MHz。中心间隔小于信道宽度，所以“信道编号不同”不等于频谱不重叠。课堂采用 1、6、11 作为常见的三个互不重叠信道配置：它们的中心间隔为 25 MHz，可以把相邻小区分开安排。这里的“三个”指常见规划下可并列部署的不重叠信道数，不能误写为整个 2.4 GHz 频段总共只有三个合法编号。具体可用编号及配置受地区与物理层方式影响。

如果相邻 AP 全部使用同一信道，它们的站点可能需要共享介质，或者在互相听不到时产生干扰。给相邻 BSS 安排不同且充分分离的信道，可以减少这种相互影响。课堂的不同教室类比在这里尤其合适：同一教室里要协调发言，分到隔音良好的不同教室后，可以同时开展不同对话。但空间部署扩大后仍需复用信道，不能认为多增加几个 AP 就自动按比例增加容量。

\begin{figure}[htbp]
\centering
\begin{tikzpicture}[x=1cm,y=1cm,every node/.style={font=\small},>=stealth]
\draw[->](0,0)--(12,0)node[right]{频率};
\draw[fill=blue!15](.7,.2)rectangle(3.7,.9);
\draw[fill=orange!20](1.35,1.1)rectangle(4.35,1.8);
\node at(2.2,.55){信道 1};\node at(2.85,1.45){相邻编号信道};
\draw[fill=green!15](4.2,.2)rectangle(7.2,.9);\node at(5.7,.55){信道 6};
\draw[fill=blue!15](7.7,.2)rectangle(10.7,.9);\node at(9.2,.55){信道 11};
\draw[dashed](1.35,0)--(1.35,1.9);\draw[dashed](3.7,0)--(3.7,1.9);
\node[align=center]at(6,-.8){相邻编号的频谱可能大量重叠；\\1、6、11 说明常见的分离配置，并非精确频谱形状};
\end{tikzpicture}
\caption{信道中心间隔与信道占用宽度的区别。矩形表示概念上的占用区间，不表示真实发射谱为理想矩形。}
\end{figure}

\subsubsection{速率自适应：在速度与可靠性之间选择}

即使不改变所用信道，一个站点仍可以采用不同的 PHY 速率。较高阶的调制让每个符号携带更多比特，但接收端必须更精细地区分可能的信号状态，因此通常要求更好的信噪比。信道较差时，较稳健的调制与编码可以牺牲速率、提高正确接收的机会。

以 802.11b 为例，课程列出的可用速率包括 1、2、5.5 和 11 Mb/s。最大 11 Mb/s 只表示设备在适用条件下支持的最高 PHY 速率，不承诺任意距离和位置都能达到这个值。终端远离 AP、穿过墙体或遇到干扰时，可能选择较低速率；靠近 AP 且链路质量良好时，可能选择较高速率。距离是影响传播的一个因素，并不是通过距离便能唯一查出速率的刻度尺。课堂的室内外距离图用于说明趋势，不是所有建筑和设备共同遵守的覆盖承诺。

速率自适应（Rate Adaptation）把这一取舍变成运行时决策。课程介绍两种反馈思路。第一种利用接收端对信噪比（Signal-to-Noise Ratio，SNR）的估计；发送端据此选择更高效或更稳健的速率。第二种利用近期发送的成功与失败历史：成功较稳定时试探更高速率，连续失败时退回较稳健的选择。后一种方法可以利用 ACK 记录，但仍面临“失败来自碰撞还是传播”的歧义。

标准规定可使用的物理层模式，并不等于给每一家设备规定同一个完整的自适应算法。不同实现可以选择不同的观察窗口、试探策略和回退规则。评价算法时不能只看名义速率，还要看成功率和重传时间：更高的发送速率若带来很多失败，最终交付数据可能反而更慢。

\subsubsection{OFDM：在一个信道内并行传输}

正交频分复用（Orthogonal Frequency Division Multiplexing，OFDM）把一个较宽频率信道分成多条窄子载波（Subcarrier），在同一时间并行传输多个符号。这里先区分两个层次：前一小节的“信道 1、6、11”用于不同网络的频率安排；OFDM 子载波则是在一个所选信道内部组织信号。不能把每条子载波直接当成一台 AP 独占的独立 WLAN 信道。

课堂以具有中心峰和旁瓣的频谱形状解释子载波：不同子载波的频谱可以在频率轴上重叠，但按合适的间隔安排时，理想情况下它们在接收处理所用的观察区间上彼此正交。直观地说，某条子载波被提取时，其他子载波不形成同样的干扰分量。因此，正交并不等于所有信号的频谱完全不接触。

发送端把待发数据分配到多个并行分支，分别调制后组合成一个发送信号；接收端对整体信号进行对应处理，恢复各分支的信息，再组合成数据流。若第 $k$ 条有效数据子载波提供的净贡献速率为 $r_k$，共有 $N_d$ 条数据子载波，则并行部分的理想总贡献为
\begin{equation}
R_{\mathrm{parallel}}=\sum_{k=1}^{N_d}r_k.
\label{eq:ofdm-intuition}
\end{equation}
这里 $R_{\mathrm{parallel}}$ 为各数据分支贡献之和，$N_d$ 不包括不承载用户数据的保护或导频位置，$r_k$ 是在一致时间口径下定义的分支贡献。此式仅将课堂的“并行提高总速率”思想形式化，不是用总带宽直接计算标准峰值的完整公式；物理保护时间、编码、导频与 MAC 开销仍需另外正确计入。

课程在这一阶段只要求理解并行传输的基本思路，并未要求从傅里叶分析推导 OFDM 实现。其重点是看懂提高速率的路线：除了让每个符号承载更多比特，也可以同时利用更多合适的信号分支。

\subsubsection{MIMO：进一步利用空间维度}

多输入多输出（Multiple Input Multiple Output，MIMO）在发送端和接收端使用多根天线，使通信不仅能够利用时间和频率，也能够利用空间传播结构。课程以多对天线承载多路数据流说明空间复用（Spatial Multiplexing）：在接收端能够区分各路空间信号的条件下，多条数据流可以同时经过同一频率资源。

多天线本身不保证数据速率按天线数成倍增加。接收端必须有足够的观测能力，传播路径也必须提供可区分的空间信息；若多个通道高度相似，就不一定能可靠分离所有流。把几根天线简单贴在一起，也不能自动得到几条独立的信道。课堂提及天线间距，正是为说明空间结构的重要性；后续 Wi-Fi 增强章节再展开实际可支持的流数及信道条件。

OFDM 与 MIMO 可以组合：OFDM 在频率分支上组织并行信号，MIMO 在空间分支上增加可用的数据流。两者也都不意味着本站可以在同一信道上边发边收。多条同时发送的数据流，与半双工／全双工讨论的发送和接收方向，是不同维度的问题。

从标准演进看，802.11b、802.11a 与 802.11g 展示了早期多速率与 OFDM 路线，802.11n 把多天线等增强带入更高速 WLAN，后续版本继续增加频宽、空间与多用户能力。标准字母名与 Wi-Fi 代际名提供检索线索，但本课程更重要的主线是机制如何增加容量，以及为何更快的 PHY 仍无法省去 MAC 协调。下面转向所有这些增强都绕不开的共享接入问题。


\subsection{为何需要介质访问控制}

无线信号会传播到一定范围内的多个接收机。若两个发射同时到达同一接收端，另一个信号通常成为干扰，可能使有用帧无法解码。课堂以多人同时讲话说明这个问题：每个人单独讲话时都可听清，重叠讲话则会使听者失去可理解的信息。MAC 的任务，就是在多个发送需求之间建立可以执行的协调规则。

需要强调，碰撞（Collision）是围绕接收结果讨论的事件。两台设备都在发射，并不能单凭这一点断定两帧必然失败；还要考虑频率、空间位置、功率与接收条件。本章时序分析采用课程的简化碰撞模型：同一接收端处互相干扰的帧若重叠，就视为不能正确接收。不讨论捕获效应等物理层例外。

\subsubsection{半双工决定了失败信息的取得方式}

课程采用半双工（Half-Duplex）收发模型：同一无线收发链路在某一时刻发送或接收，不能一边发送自己的帧，一边可靠接收其他节点的弱信号来检测碰撞。因此，Wi-Fi 的基本 MAC 使用载波侦听多址接入／碰撞避免（Carrier Sense Multiple Access with Collision Avoidance，CSMA/CA）。

有线以太网历史上的载波侦听多址接入／碰撞检测（CSMA/CD）可用来比较设计思路，但不能移植成“Wi-Fi 发送时发现碰撞就立即停止”。而且 CSMA/CD 也不是现代交换式全双工以太网的工作机制。对这里的 Wi-Fi 单播，发送端依据确认是否到达来判断交换是否成功；没有确认只表示未得到成功证据，并不直接揭示失败原因。

\subsubsection{MAC 可以看成一个带服务台的队列}

高层到达的分组先进入 MAC 队列。先入先出（First In First Out，FIFO）是课程采用的基本服务模型：队头帧获得服务时，需要经历信道接入、数据发送和确认等待；失败时还可能经历重新竞争和重传。服务时间因而不是固定的物理发送时长。

为了明确统计边界，可写成
\begin{equation}
D_{\mathrm{MAC}}=W_{\mathrm{q}}+S_{\mathrm{MAC}}.
\label{eq:mac-delay}
\end{equation}
式中，$D_{\mathrm{MAC}}$ 为一个帧从进入 MAC 队列到服务结束的总时间；$W_{\mathrm{q}}$ 为等待前方帧完成的排队时间；$S_{\mathrm{MAC}}$ 为该帧成为服务对象后经历的 MAC 服务时间。这里假设一个串行服务台，并把该帧的接入等待、交换以及必要重试都计入服务时间。此式是时间分解，不能据此认为排队时间与服务时间固定或相互独立。

上行（Uplink）表示站点向 AP 发送，下行（Downlink）表示 AP 向站点发送。AP 在 DCF 中同样要争用无线介质；它不是只接收而不发送的特殊旁观者。若多条下行业务集中在 AP 队列中，还要区分“每个竞争实体的接入机会”与“每条用户业务获得的吞吐量”。

\section{DCF：从侦听到一次受保护的交换}
\label{ch:dcf-exchange}

\subsection{分布式协调的基本原则}

分布式协调功能（Distributed Coordination Function，DCF）由各站点依据共同规则自行决定接入时机，不依赖 AP 为每个普通数据帧逐次发放发送许可。集中式的点协调功能（Point Coordination Function，PCF）则由协调者组织发送。这两种思路在课程开头同时出现，但本章先完成 DCF，因为理解它才能看清 PCF 为何需要另一套优先级与调度安排。

DCF 的第一条原则是先听再发。物理层通过空闲信道评估（Clear Channel Assessment，CCA）判断介质活动；课程用接收能量与门限比较说明其直观含义。实际 CCA 还可以利用可识别信号的信息，因此“能量检测”应理解为教学模型，而不是所有实现的唯一判据。

单次观察为空闲仍然不够。一方面，其他节点可能恰好也作出相同决定；另一方面，空闲可能只是一个正在进行的交换中两个帧之间的短暂停顿。于是需要等待间隔、确认优先权、虚拟侦听和随机退避。下面先暂时固定“当前站点已经具备接入资格”，把一次数据交换的内部结构讲清楚，再补足资格如何获得。

\subsection{时隙与帧间间隔构成时间规则}

时隙（Slot Time）是随机退避计数的基本单位，记为 $\sigma$。它与载波侦听、传播以及收发切换等时间要求有关，具体数值取决于所用物理层。这里的离散时隙不是为每个用户永久分配的独占时段，也不能把一帧数据的长度限制成一个时隙。

帧间间隔（Interframe Space，IFS）是一类有规定用途的等待时间。最短帧间间隔（Short IFS，SIFS）用于已开始交换中的即时响应；DCF 帧间间隔（DCF IFS，DIFS）用于普通分布式竞争；PCF 帧间间隔（PCF IFS，PIFS）处在二者之间，使集中式协调能够先于普通 DCF 竞争取得介质。基本关系为
\begin{align}
T_{\mathrm{PIFS}}&=T_{\mathrm{SIFS}}+\sigma,\\
T_{\mathrm{DIFS}}&=T_{\mathrm{SIFS}}+2\sigma.
\label{eq:ifs-relations}
\end{align}
这里 $T_{\mathrm{SIFS}}$、$T_{\mathrm{PIFS}}$ 和 $T_{\mathrm{DIFS}}$ 分别表示三种间隔的持续时间，$\sigma>0$ 为当前物理层的时隙长度。因而 $T_{\mathrm{SIFS}}<T_{\mathrm{PIFS}}<T_{\mathrm{DIFS}}$。这些关系说明优先权来自等待时间差，而不是某个节点可以随时打断已经开始的传输。

例如，数据帧结束后，正确接收它的节点只需等待 SIFS 就发送确认。其他准备发送普通数据的站点必须等待 DIFS；在它们等够之前，确认已开始，使它们重新看到信道忙。短间隔因此保护了现有交换的连续性。优先级属于特定帧和特定协议阶段，不能理解为同一台设备所有帧永远拥有相同优先权。

\subsection{单播的 DATA--ACK 交换}

单播（Unicast）面向一个接收方。发送端发出一个 MAC 协议数据单元（MAC Protocol Data Unit，MPDU），接收端检查循环冗余校验（Cyclic Redundancy Check，CRC）后，仅在正确接收时返回确认（Acknowledgment，ACK）。ACK 的等待间隔是 SIFS，它不重新参加一次普通数据竞争。

\begin{figure}[htbp]
\centering
\begin{tikzpicture}[x=0.95cm,y=0.95cm,every node/.style={font=\small},>=stealth]
\node[anchor=east] at (0,1.5) {发送端};
\node[anchor=east] at (0,0) {接收端};
\node[anchor=east] at (0,-1.5) {旁听站点};
\foreach \y in {1.5,0,-1.5}{\draw[->] (0,\y)--(11.6,\y);}
\draw[fill=gray!15] (0.4,1.25) rectangle (2,1.75);\node at (1.2,1.5){DIFS};
\draw[fill=blue!15] (2,1.25) rectangle (6.5,1.75);\node at (4.25,1.5){DATA};
\draw[dashed] (6.5,1.1)--(6.5,-1.7);
\draw[<->] (6.5,0.6)--node[above]{SIFS}(7.8,0.6);
\draw[fill=green!20] (7.8,-0.25) rectangle (9.5,0.25);\node at (8.65,0){ACK};
\draw[dashed] (9.5,0.35)--(9.5,-1.7);
\draw[fill=orange!15] (3,-1.75) rectangle (9.5,-1.25);
\node at (6.25,-1.5){根据首部预留到 ACK 结束};
\node[align=center] at (5.4,2.35){发送端已满足本次接入条件\\此处先省略随机退避};
\end{tikzpicture}
\caption{基本单播交换。忽略传播时延，矩形仅表达先后关系，不按真实时长比例绘制。}
\label{fig:data-ack}
\end{figure}

发送端在规定的等待窗口内没有收到预期 ACK，就不能把本次交换记为成功。可能是数据帧碰撞，也可能是数据帧受噪声破坏，还可能是数据已经到达而 ACK 丢失。停止等待自动重传请求（Stop-and-Wait ARQ）要求发送端在相应重试限制内重新尝试；但重新尝试必须重新取得信道，不能在超时瞬间直接重发。

课堂将发送完成后用于等待 ACK 的失败等待时间记为 $T_F$。这一符号用于描述发送端的超时判断。$T_F$ 的设置必须覆盖预期响应的时间要求，具体实现还涉及物理层参数；它不是一个任意的 DIFS，也不是后面讨论的 EIFS。\textbf{“发送端等自己的 ACK”与“旁听者为了保护别人的 ACK 而推迟接入”是不同的状态。}

\subsection{物理侦听之外，还要有虚拟侦听}

\subsubsection{Duration 与 NAV 保护的是一段交换}

先区分 MAC 的三类帧：数据帧承载要传输的用户数据；控制帧帮助组织交换，例如 ACK、RTS 与 CTS；管理帧建立和维护网络关系，例如信标、探测与关联报文。三类帧都要经过物理层发射，但它们承担的功能、是否向上交付以及所适用的响应规则并不相同。不能把“控制信息”理解成无需占用无线时间。

无线信道在 DATA 和 ACK 之间的 SIFS 内可能没有信号，但这不代表一个全新的竞争周期已经开始。帧首部中的持续时间字段（Duration）告知能够正确读取它的邻居：当前交换还需要占用介质多久。邻居据此维护网络分配向量（Network Allocation Vector，NAV），在预留结束前按介质仍忙处理。这称为虚拟载波侦听（Virtual Carrier Sensing）。

NAV 可以理解为一个随时间递减的逻辑计时器。它不是站点的退避计数器：NAV 表示需要尊重的预约时间，退避计数器表示本站尚未消耗完的竞争等待。NAV 会随真实时间流逝而走向零；退避计数器却只能在合格的空闲时隙中递减。

对一个没有后续分片的基本 DATA--ACK 交换，若以 DATA 接收结束时刻作为参考点，后续预留时间是
\begin{equation}
D_{\mathrm{DATA}}=T_{\mathrm{SIFS}}+T_{\mathrm{ACK}}.
\label{eq:duration-data}
\end{equation}
式中，$D_{\mathrm{DATA}}$ 表示 DATA 结束之后需要保护的剩余交换时长，$T_{\mathrm{ACK}}$ 是 ACK 的完整空中发送时间，包含其物理层开销。假定传播时间在图解中忽略，且本次交换在 ACK 后结束。若从 DATA 开始或从首部被读取的时刻画整段保护区，还必须把尚未结束的 DATA 部分计入。\textbf{不能把“从哪个时刻画 NAV”与“Duration 字段表示的剩余时间”混用。}

\subsubsection{两种侦听必须同时允许}

DCF 将物理与虚拟侦听共同作为判断依据。可用逻辑关系表示为
\begin{equation}
\text{可认为介质空闲}
\quad\Longleftrightarrow\quad
\text{物理侦听空闲}\ \land\ \mathrm{NAV}=0.
\label{eq:two-sensing}
\end{equation}
式中，$\land$ 表示逻辑“并且”；$\mathrm{NAV}=0$ 表示本地预留计时已结束。该式只是介质空闲的判据，不等于已经满足全部发送条件：站点还须完成适用的 IFS 和退避要求。

物理空闲而 NAV 为正并不矛盾。例如，某个旁听站点听不到接收端的 ACK，或者当前处于帧间间隔，它的物理层都可能看不到信号，但仍应尊重已经了解到的交换。在 DATA 失败、接收端不发 ACK 时，已经设置 NAV 的邻居也不会凭自己的猜测立即取消预留。局部观察必须服从既有的预约信息。

\subsection{隐藏终端：局部空闲不代表接收端安全}

隐藏终端（Hidden Terminal）问题说明物理侦听的空间局限。设 A 与 B 都能和 AP 通信，但 A 与 B 互相听不到。A 先发送给 AP，B 随后拿到一个待发帧。B 的本地侦听可能持续为空闲，于是也向 AP 发送；两帧在 AP 处重叠。在本章的碰撞模型下，AP 无法正确接收，自然也不能给出成功确认。

\begin{figure}[htbp]
\centering
\begin{tikzpicture}[every node/.style={font=\small},>=stealth]
\node[draw,circle,minimum size=0.85cm,fill=blue!10] (a) at (0,0) {A};
\node[draw,rectangle,minimum size=1cm,fill=green!15] (ap) at (4,0) {AP};
\node[draw,circle,minimum size=0.85cm,fill=blue!10] (b) at (8,0) {B};
\draw[<->,thick] (a)--node[above]{可通信}(ap);
\draw[<->,thick] (ap)--node[above]{可通信}(b);
\draw[dashed] (a) to[bend right=25] node[below]{A 与 B 不能相互侦听} (b);
\node[align=center] at (4,-2){A 发送时，B 可能误判本地信道为空闲；\\真正需要保护的是 AP 的接收过程};
\end{tikzpicture}
\caption{隐藏终端的最小场景。无线可听关系与接收端的干扰关系必须分别分析。}
\end{figure}

问题不在于 B 故意违反规则，而在于 B 没有取得足够的信息。即使它认真等待 DIFS，也可能始终听不到 A。因此，单纯增加一个很短的固定侦听间隔，无法解决所有隐藏终端问题。需要让接收端也发出可被其邻居听到的预留信息。

\subsection{RTS/CTS：由收发双方共同宣布预留}

\subsubsection{四帧交换的完整次序}

请求发送（Request To Send，RTS）与允许发送（Clear To Send，CTS）构成数据传输前的短握手。发送端先按普通竞争规则获得机会，发 RTS；预期接收端收到 RTS 后等待 SIFS，发 CTS；发送端正确收到 CTS 后等待 SIFS，发 DATA；接收端正确收到 DATA 后等待 SIFS，发 ACK。因此完整顺序是
\begin{center}
RTS $\rightarrow$ SIFS $\rightarrow$ CTS $\rightarrow$ SIFS $\rightarrow$ DATA $\rightarrow$ SIFS $\rightarrow$ ACK。
\end{center}
DIFS 和必要的退避发生在 RTS 之前，不应机械插到每两个帧之间。CTS、DATA 和 ACK 都属于已经发起并被预留保护的交换。

\begin{figure}[htbp]
\centering
\begin{tikzpicture}[x=0.91cm,y=0.9cm,every node/.style={font=\small},>=stealth]
\node[anchor=east] at (0,2){发送端};\node[anchor=east] at (0,0.8){接收端};
\node[anchor=east] at (0,-0.5){听到 RTS};\node[anchor=east] at (0,-1.7){仅听到 CTS};
\foreach \y in {2,0.8,-0.5,-1.7}{\draw[->](0,\y)--(12.7,\y);}
\draw[fill=blue!15] (0.8,1.75) rectangle (2.3,2.25);\node at (1.55,2){RTS};
\draw[fill=green!20] (3.2,0.55) rectangle (4.5,1.05);\node at (3.85,0.8){CTS};
\draw[fill=blue!15] (5.4,1.75) rectangle (9.3,2.25);\node at (7.35,2){DATA};
\draw[fill=green!20] (10.2,0.55) rectangle (11.5,1.05);\node at (10.85,0.8){ACK};
\foreach \a/\b in {2.3/3.2,4.5/5.4,9.3/10.2}{\draw[<->] (\a,2.65)--(\b,2.65);\node[font=\scriptsize,above] at ({(\a+\b)/2},2.65){SIFS};}
\draw[fill=orange!15] (2.3,-0.75) rectangle (11.5,-0.25);\node at (6.9,-0.5){RTS 给出的预留};
\draw[fill=orange!15] (4.5,-1.95) rectangle (11.5,-1.45);\node at (8,-1.7){CTS 给出的剩余预留};
\node[align=center] at (6,-2.65){两类邻居获知预留的时刻不同，但都保护到同一交换终点};
\end{tikzpicture}
\caption{RTS/CTS 的作用不仅是取得接收端许可，更是把预约信息传播到收发双方的邻域。图中省略初始 DIFS 与退避，并忽略传播时延。}
\label{fig:rts-cts-nav}
\end{figure}

以各帧结束时刻为参考，且忽略传播时延，剩余时长可写为
\begin{align}
D_{\mathrm{RTS}}&=3T_{\mathrm{SIFS}}+T_{\mathrm{CTS}}+T_{\mathrm{DATA}}+T_{\mathrm{ACK}},\\
D_{\mathrm{CTS}}&=2T_{\mathrm{SIFS}}+T_{\mathrm{DATA}}+T_{\mathrm{ACK}},\\
D_{\mathrm{DATA}}&=T_{\mathrm{SIFS}}+T_{\mathrm{ACK}}.
\label{eq:rts-durations}
\end{align}
其中 $D_{\mathrm{RTS}}$、$D_{\mathrm{CTS}}$、$D_{\mathrm{DATA}}$ 分别表示对应帧结束后的剩余预留时间；$T_{\mathrm{RTS}}$ 等带帧名的 $T$ 表示该帧的完整发送时长。$T_{\mathrm{RTS}}$ 没有出现在第一式右侧，正因为参考点已经是 RTS 结束。上式适用于这里的一次 RTS--CTS--DATA--ACK 交换，不是所有分片或聚合场景的通用 Duration 公式。

\subsubsection{分别站在每个旁听者的位置分析}

听到 RTS 的邻居知道发送端准备占用信道，于是设置 NAV。听不到 RTS、却能听到 CTS 的隐藏邻居，也能从 CTS 获知接收端接下来要接收数据，因而保持静默。两个帧都听到的节点则依据后续合法预留信息维护 NAV。分析时应先画出谁能听到谁，再逐条标明哪个帧触发了哪个节点的预留，不能把所有“其他节点”都画成同一时刻开始倒计时。

课程还使用相邻同频 AP 的例子。AP 1 向某个终端发送时，AP 2 虽然属于另一个 BSS，但若能正确听到相关控制帧，也可能需要推迟发送。网络的管理归属不会使电磁干扰自动消失。相反，一个已关联于 AP 1 的终端，也不一定听得到 AP 1 覆盖区里的每个其他终端。

\subsubsection{预留减少风险，但没有消灭碰撞}

RTS 本身仍通过竞争发送，多个 RTS 可以碰撞。邻居还可能由于信道差错漏收 CTS；若没有得到预留信息，它仍可能在后续 DATA 期间发射。因此，RTS/CTS 是降低隐藏终端风险与减少长帧碰撞损失的工具，不是“有了握手就绝不会碰撞”的保证。

如果发送端没有收到 CTS，它不能认定数据发送已经获准，应在相应超时后重新竞争。没有 CTS 也不能唯一归因于接收端不在线：RTS 丢失、CTS 丢失或碰撞都可能导致同一表象。

握手的代价是额外的控制帧及两段额外 SIFS。设忽略竞争变化、传播与差错，比较同一 DATA--ACK 交换是否增加握手，则额外时间为
\begin{equation}
\Delta T=T_{\mathrm{RTS}}+T_{\mathrm{CTS}}+2T_{\mathrm{SIFS}}.
\label{eq:rts-overhead}
\end{equation}
$\Delta T$ 表示纯粹由 RTS/CTS 引入的交换时间增量；三个 $T$ 分别是 RTS、CTS 发送时长与 SIFS 时长。这只是成功交换的增量，未计入握手减少重传所节省的时间，因此不能仅凭它断言握手降低总体性能。小载荷时额外开销占比大，长载荷发生碰撞时浪费更大，这解释了为何可依据 RTS 门限（RTS Threshold）决定是否启用握手。门限是取舍参数，并非在任何拓扑、负载和误码条件下都最优的固定常数。

\section{随机退避、竞争公平性与协议效率}
\label{ch:backoff-fairness}

\subsection{共同等待结束，并不等于可以共同发送}

设 A 正在与 AP 交换数据，B 和 C 都已有待发帧。B、C 观察到同一忙期结束，继而同时等待相同的 DIFS。如果等待完就直接发送，两者会再次同时占用介质。在这个简化场景里，固定等待把两个发送者同步到了同一个时刻，反而不能把它们分开。

随机退避（Random Backoff）在共同的等待规则之后加入独立的随机计数。各站点不再全部选择同一个接入时刻，从而以概率方式打破同步。碰撞避免的“避免”应理解为降低风险，并非保证每次只有一个节点发送。

\subsection{退避抽样与计数：时间只在空闲时前进}

\subsubsection{从竞争窗口抽取整数}

竞争窗口（Contention Window，CW）给出随机整数的上界。若当前窗口为 $CW$，站点抽取
\begin{equation}
K\sim\operatorname{Uniform}\{0,1,\ldots,CW\},\qquad B=K\sigma.
\label{eq:backoff-draw}
\end{equation}
式中，$K$ 是以时隙为单位的退避计数初值；$CW$ 是非负整数上界；$\sigma$ 是一个时隙的时长；$B$ 是需要累计消耗的空闲退避时间。离散均匀分布表示每个整数以相同概率被选择。因为零也在范围内，一共有 $CW+1$ 个选择，而不是 $CW$ 个。

$B$ 不是一个从抽样时刻开始连续运行、到点必发的墙钟定时器。若中途信道变忙，倒计数会冻结；实际经过的日历时间可能远大于 $B$。把计数器画成在他人传输期间仍然递减，是理解 DCF 时最常见的错误之一。

\subsubsection{冻结保留进度，恢复须重新等待 DIFS}

执行退避时，站点先等介质连续空闲 DIFS，然后才开始消耗空闲时隙。每完整经过一个合格的空闲时隙，$K$ 减一；若信道变忙，保留当前剩余值。忙期结束后还须再次连续空闲 DIFS，才能继续递减。\textbf{冻结不是抽取新的随机数，也不是把计数器重置为最初的值。}

例如抽到 $K=3$，第一个时隙空闲后剩余为 2；第二个时隙尚未完整结束就出现他人的传输，因此这个未完成的空闲时隙不应被扣除。整个忙期内剩余值仍为 2。再次等够 DIFS，再消耗两个完整空闲时隙，计数器归零，站点才获得本次发送机会。

\begin{figure}[htbp]
\centering
\begin{tikzpicture}[x=0.9cm,y=1cm,every node/.style={font=\small},>=stealth]
\draw[->](0,0)--(13,0);
\draw[fill=gray!15](0,0)rectangle(1.7,0.65);\node at(.85,.325){DIFS};
\draw[fill=blue!15](1.7,0)rectangle(2.7,.65);\node at(2.2,.325){1 槽};
\draw[fill=orange!20](2.7,0)rectangle(6.8,.65);\node at(4.75,.325){介质忙：冻结};
\draw[fill=gray!15](6.8,0)rectangle(8.5,.65);\node at(7.65,.325){DIFS};
\draw[fill=blue!15](8.5,0)rectangle(9.5,.65);\node at(9,.325){1 槽};
\draw[fill=blue!15](9.5,0)rectangle(10.5,.65);\node at(10,.325){1 槽};
\draw[fill=green!20](10.5,0)rectangle(12.5,.65);\node at(11.5,.325){发送};
\node at(1.7,-.45){$K=3$};\node at(2.7,-.45){$2$};
\node at(4.75,-.45){始终为 $2$};\node at(8.5,-.45){$2$};
\node at(9.5,-.45){$1$};\node at(10.5,-.45){$0$};
\node[align=center]at(6.2,1.5){只累计三个完整空闲时隙；\\忙期和恢复前的 DIFS 不消耗计数};
\end{tikzpicture}
\caption{退避冻结与恢复的例子。时间轴按逻辑片段示意，并非按标准时长比例绘制。}
\label{fig:backoff-freeze}
\end{figure}

\subsubsection{一次竞争的流程}

\begin{figure}[htbp]
\centering
\begin{tikzpicture}[every node/.style={font=\small},>=stealth]
\node[draw,rounded corners,text width=5.2cm,align=center](a)at(0,0){确定当前 CW，抽取初始计数 $K$};
\node[draw,rounded corners,text width=5.2cm,align=center](b)at(0,-1.2){等待物理及虚拟介质连续空闲 DIFS};
\node[draw,rounded corners,text width=5.2cm,align=center](c)at(0,-2.4){检查剩余计数：$K=0$？};
\node[draw,rounded corners,text width=5.2cm,align=center](d)at(0,-3.8){侦听下一个时隙：完整空闲？};
\node[draw,rounded corners,text width=5.2cm,align=center](e)at(0,-5.2){$K\leftarrow K-1$，继续计数};
\node[draw,rounded corners,text width=2.2cm,align=center](f)at(4.6,-2.4){发送 DATA\\或先发送 RTS};
\draw[->](a)--(b);\draw[->](b)--(c);
\draw[->](c)--node[above]{是}(f);
\draw[->](c)--node[right]{否}(d);
\draw[->](d)--node[right]{是}(e);
\draw[->](e.west)--(-3.3,-5.2)--(-3.3,-2.4)--(c.west);
\draw[->](d.east)--(3.2,-3.8)--node[right,align=left]{否：冻结\\保留 $K$}(3.2,-1.2)--(b.east);
\end{tikzpicture}
\caption{退避核心流程。失败后的窗口更新与成功后的发送后退避，在流程之外决定下一次进入时采用的 CW。}
\end{figure}

\subsection{三类需要退避的情形}

第一类是\textbf{新帧遇到忙信道}。站点原本没有待发数据，新帧到达时发现介质忙，不能等大家共同的 DIFS 一结束就直接发，而要进入随机退避。

第二类是\textbf{失败后重新竞争}。发送端没有取得预期确认，在超时后需要重试。它必须重新执行退避，而不是仅等待一个固定间隔。对同一帧连续失败，窗口还要按规则扩大。

第三类是\textbf{成功后的发送后退避}。即便已经成功完成一个交换，站点也不能因队列中还有下一帧，就获得不经过随机竞争的连续发送特权。此时以最小窗口开始新的退避周期，下一节将解释其公平性目的。

“空闲一个 DIFS 后直接发送”的简化叙述只适用于没有尚未完成的退避义务、没有相关错误恢复要求的空闲接入情形。例如，站点此前的发送后退避已完成，队列随后才收到新帧，且介质满足空闲条件。它不能用作所有帧的通用接入规则，更不能覆盖已经被要求退避的帧。

\subsection{二进制指数退避：扩大的是选择空间}

最小竞争窗口记为 $CW_{\min}$，最大竞争窗口记为 $CW_{\max}$。为与课堂推导一致，令 $i=1$ 表示为当前帧执行的第一轮退避；后续 $i$ 随该帧失败后需要再次抽样而增加。窗口序列为
\begin{align}
CW_1&=CW_{\min},\\
CW_i&=\min\{2(CW_{i-1}+1)-1,\ CW_{\max}\},\quad i\geq2.
\label{eq:cw-recursion}
\end{align}
式中，$CW_i$ 为第 $i$ 轮退避采用的整数上界；$\min\{\cdot\}$ 表示不得超过最大窗口。该式采用课程的连续重试模型。这里的 $i$ 不应不加说明地称为“空中发送次数”：某帧的首次接入可能处于简化的无需初始退避情形，课堂索引针对的是退避轮次。

在尚未达到上限时，也可用退避指数（Backoff Exponent）$m_i$ 表示：
\begin{equation}
CW_i=2^{m_i}-1,\qquad m_i=m_1+i-1.
\label{eq:backoff-exp}
\end{equation}
$m_1$ 为初始指数，$m_i$ 为第 $i$ 轮指数；每轮增加一，使可选整数个数 $CW_i+1=2^{m_i}$ 精确翻倍。\textbf{翻倍的是包含零在内的选择数，并非数值上严格把 CW 乘二。}所以从 31 变成的是 63，而不是 62。

以课程的早期 802.11b 参数为例，$m_1=5$，因而 $CW_{\min}=31$。采用课件给定的 $CW_{\max}=1023$ 与 $\sigma=20\,\mu\mathrm{s}$，得到表\ref{tab:cw-growth}。这些数值是指定物理层示例，不能当作所有 Wi-Fi 版本的统一常数。

\begin{table}[htbp]
\centering
\caption{早期 802.11b 示例中的窗口增长与纯退避均值}
\label{tab:cw-growth}
\begin{tabular}{cccr}
\toprule
退避轮次 $i$ & $CW_i$ & 可选整数个数 & 平均空闲退避时间\\
\midrule
1 & 31 & 32 & $310\,\mu\mathrm{s}$\\
2 & 63 & 64 & $630\,\mu\mathrm{s}$\\
3 & 127 & 128 & $1270\,\mu\mathrm{s}$\\
4 & 255 & 256 & $2550\,\mu\mathrm{s}$\\
5 & 511 & 512 & $5110\,\mu\mathrm{s}$\\
6 及以后 & 1023 & 1024 & $10230\,\mu\mathrm{s}$\\
\bottomrule
\end{tabular}
\end{table}

表中的均值来自离散均匀分布：
\begin{equation}
\mathbb{E}[K]=\frac{CW}{2},\qquad
\mathbb{E}[B]=\frac{CW}{2}\sigma.
\label{eq:mean-backoff}
\end{equation}
$\mathbb{E}[\cdot]$ 为数学期望；$K$、$B$、$CW$、$\sigma$ 的含义与式\eqref{eq:backoff-draw}相同。这里仅统计要消耗的空闲退避时长，不含 DIFS、不含冻结忙期，也不含重传帧本身的发送时间。因此真实平均接入时延一般比该数值更大。

扩大窗口并不保证某一次抽样结果变大。窗口从 31 扩到 63 后，站点仍然可能抽到零；增大的是可选范围和平均值。达到上限后继续失败，也不再无限增加窗口。当前帧处理结束、转向新帧时，从最小窗口重新开始，以免一次历史失败长期惩罚后续业务。

\subsection{两节点碰撞概率的完整推导}

考虑课程中的两个竞争节点 B、C。设它们在同一竞争起点开始倒计数，使用相同窗口 $CW$，独立且均匀抽取初始整数 $K_B$ 与 $K_C$；没有隐藏终端、没有其他竞争者，也不考虑传播造成的额外碰撞窗口。若两者抽到同一整数，就在同一时刻发送，形成该模型中的碰撞；抽到不同整数时，较大的计数会在对方开始发送后冻结。

先固定 B 的抽样值为 $n$。由于 C 的抽样独立，条件概率为
\begin{equation}
\Pr(K_C=n\mid K_B=n)=\frac{1}{CW+1},
\qquad n\in\{0,1,\ldots,CW\}.
\label{eq:collision-conditional}
\end{equation}
$\Pr(\cdot)$ 表示概率，竖线表示已知条件；$n$ 为某个合法的整数退避值。独立性意味着知道 B 抽到了什么，并不会改变 C 各个选择的概率。

再遍历 B 的全部可能取值，使用全概率公式：
\begin{align}
p_{\mathrm{col}}
&=\Pr(K_B=K_C)\\
&=\sum_{n=0}^{CW}\Pr(K_C=n\mid K_B=n)\Pr(K_B=n)\\
&=\sum_{n=0}^{CW}\frac{1}{CW+1}\frac{1}{CW+1}
=\frac{1}{CW+1}.
\label{eq:two-node-collision}
\end{align}
$p_{\mathrm{col}}$ 是这两个节点在本轮抽样和竞争假设下发生同计数碰撞的概率；求和包含 $CW+1$ 项。最后一式说明，$CW=31$ 时概率为 $1/32=3.125\%$，扩大到 63 后变为 $1/64=1.5625\%$。

这个结果回答的是限定模型中的问题，不是一个任意繁忙 Wi-Fi 网络的通用丢包率。多个竞争者、不同重试阶段、不同残余计数、隐藏终端和无线差错都会改变实际行为。尤其不能把 $1/(CW+1)$ 当成任意一个帧的独立成功／失败概率，再直接套入所有吞吐量计算。

推导同时揭示了性能取舍：窗口越大，两个独立节点抽中相同值的概率越低，但平均空闲等待越长。若失败确实来自拥挤，扩大窗口有帮助；若失败主要来自衰落或噪声，扩大窗口并不能改善接收信号质量，反而增加时延。发送端只观察到缺失 ACK，未必能够区分这两类失败，这也是指数退避的局限。

\subsection{发送后退避为何关乎公平}

发送后退避（Post-Backoff）要求成功发送者再次使用 $[0,CW_{\min}]$ 范围参与等待。它不是对“成功”进行处罚，而是防止刚刚成功的站点凭借连续队列获得系统性的抢先机会。

设 A 已完成 DATA--ACK，B 与 C 在 A 发送期间已获得待发数据，正在等候竞争。如果取消 A 的发送后退避，A 队列中的下一帧可能只等待 DIFS 就发出，而 B、C 仍须等待 DIFS 加各自的随机剩余时隙。只要它们剩余计数不为零，A 就能先于它们再次占用信道，迫使它们冻结。若 A 连续有数据，这种优势可以反复出现。

不能以“A 成功时队列恰好为空”为理由免除发送后退避。新帧可能在成功后几微秒到达；若此时便允许它跳过退避，同样会获得相对于其他等待者的优势。所以发送后退避可以在队列空时运行。若新帧到来之前它已完成，且满足介质空闲等条件，才可能落入前述无需额外初始退避的情形。

课堂借移除某个机制的思路说明其必要性：先观察完整协议，再假设删除发送后退避，追踪谁获益、谁的等待不断延长。这种消融式推理比背诵“它用于公平”更有解释力。这里的公平是接入机会层面的目标，不意味着每个节点在任意短时间窗口内发送次数完全相同，也不意味着不同 PHY 速率的业务必然具有相同吞吐量。

\subsection{广播和多播为何不用同样的 ACK 过程}

广播（Broadcast）面向广播范围内的所有接收者；多播（Multicast）面向一个接收者组。若所有成功接收的站点都在同一 SIFS 后立即发 ACK，确认帧会互相干扰。若再设计一套确认调度，又会带来更多状态和开销。

因此，本章讨论的基本广播／多播机制不使用单播式逐帧 ACK，也不依据 ACK 超时进行 MAC 层重传。发送者仍须遵守介质访问规则，遇到忙信道仍须退避；只是发送结束后，没有一个“等待所有接收端确认”的阶段。\textbf{无需 ACK 不等于无需争用，也不等于每个接收者一定收到。}更高层若需要可靠性，仍须采用适当机制；后续标准的扩展不能反向套入这里的基本模型。

\subsection{PHY 速率为何不等于有效吞吐量}

\subsubsection{先固定计算的时间边界}

物理层速率 $R$ 表示实际发送相应比特时的速率。有效吞吐量（Goodput）则把成功交付的有效载荷除以整个观察时间；等待、帧首部、前导码、确认、碰撞与重传都会占用时间。为便于与课堂比较，这里用一个无碰撞、无差错、单链路连续发送的模型，只计算一次成功服务周期。

设有效载荷为 $L$ 字节，MAC 层开销为 $H$ 字节，数据速率为 $R$ 比特每秒，物理前导码与物理头开销为 $T_{\mathrm{PHY}}$，则
\begin{equation}
T_{\mathrm{DATA}}=T_{\mathrm{PHY}}+\frac{8(H+L)}{R}.
\label{eq:data-airtime}
\end{equation}
$T_{\mathrm{DATA}}$ 为 DATA 完整发送时间。因 $H$、$L$ 的单位为字节，乘 8 才得到比特；除以 $R$ 得到秒。此式假定被合并进 $H+L$ 的字段都以所列数据速率发送。不能把物理前导码的低速开销重复算入这部分。

一个简化的平均成功周期为
\begin{equation}
\overline{T}_{\mathrm{cyc}}=T_{\mathrm{DIFS}}+\mathbb{E}[B]+T_{\mathrm{DATA}}
+T_{\mathrm{SIFS}}+T_{\mathrm{ACK}},\qquad
G=\frac{8L}{\overline{T}_{\mathrm{cyc}}}.
\label{eq:goodput-cycle}
\end{equation}
$\overline{T}_{\mathrm{cyc}}$ 是平均周期时长，$G$ 是以比特每秒计的有效吞吐量。退避可归在本次发送之前，也可在连续流量的周期记账中归为上一次的发送后退避；只要边界一致，每周期只计一次，不能前后重复加入。此式不包含排队、其他节点导致的冻结或失败重传。

\subsubsection{11 Mb/s 链路的数值例子}

采用课件的 802.11b 长前导码示例：$L=1500$ 字节，$H=34$ 字节，$R=11\,\mathrm{Mb/s}$，$T_{\mathrm{PHY}}=192\,\mu\mathrm{s}$，$T_{\mathrm{SIFS}}=10\,\mu\mathrm{s}$，$T_{\mathrm{DIFS}}=50\,\mu\mathrm{s}$，$\sigma=20\,\mu\mathrm{s}$，$CW_{\min}=31$。这里的 34 字节是该计算所选用的合并 MAC 开销参数，并非宣称所有 MAC 帧的首部长度都是 34 字节。

ACK 采用本例明确指定的 $1\,\mathrm{Mb/s}$ 速率，长度为 14 字节，并计入同样的 $192\,\mu\mathrm{s}$ 物理开销。于是
\begin{align*}
T_{\mathrm{DATA}}&=192+\frac{8(1500+34)}{11}
\approx1307.6\,\mu\mathrm{s},\\
\mathbb{E}[B]&=\frac{31}{2}\times20=310\,\mu\mathrm{s},\\
T_{\mathrm{ACK}}&=192+\frac{14\times8}{1}=304\,\mu\mathrm{s},\\
\overline{T}_{\mathrm{cyc}}&=50+310+1307.6+10+304
=1981.6\,\mu\mathrm{s},\\
G&=\frac{12000}{1981.6}\approx6.06\,\mathrm{Mb/s}.
\end{align*}
上述代入利用了“比特除以 Mb/s 的数值等于微秒”这一单位换算。$12000$ 是 1500 字节有效载荷对应的比特数；因此约 $6.06\,\mathrm{Mb/s}$ 不是无线芯片的 PHY 速率下降了，而是有用数据只占整个服务周期的一部分。

若定义效率 $\eta=G/R$，其中 $\eta$ 为无量纲的有效吞吐量与数据 PHY 速率之比，则本例 $\eta\approx55\%$。该结果依赖所列参数与理想化假设，不能当作所有 11 Mb/s 设备的实测速率。碰撞、低质量链路和多用户竞争都会进一步改变结果。

对固定开销而言，较长载荷能提高摊销效率；但长帧一旦失败，损失的空中时间也更多。这与 RTS/CTS 的选择形成同一条逻辑：设计不能只追求开销小，也不能只追求保护多，而应比较成功收益与失败代价。

\subsubsection{地址字段反映不同的转发角色}

课堂在讨论开销时还区分了原始源、最终目的、当前无线发送者与当前无线接收者。普通简单链路中，部分角色可以重合；跨分布系统转发时，它们可能不同。802.11 的不同帧格式因此可能需要不同数量的地址字段，特定无线分布转发格式可以使用四地址。

这不意味着 ESS 中的每一个帧都具有四个地址。必须先标明当前这一跳是谁发给谁，再区分原始源与最终目的，最后选择对应格式。地址开销随格式变化，也解释了为何吞吐量例子必须明确自己的 MAC 开销假设。

\subsection{EIFS：错误接收者需要额外的保护时间}

扩展帧间间隔（Extended IFS，EIFS）在本段课程结束时引入，用于处理站点听到了一个帧、但没有正确解码的情况。这样的站点可能不知道首部所载的剩余预留时间，也可能听不到接收端接下来发出的 ACK。如果它仅按通常 DIFS 恢复竞争，便有机会破坏实际上仍在继续的交换。

EIFS 的思路是让错误接收的站点在介质重新空闲后更保守地等待，为潜在的 SIFS 和 ACK 留出时间。为解释这一保护目的，可将其基本时间构成写为
\begin{equation}
T_{\mathrm{EIFS}}\approx T_{\mathrm{SIFS}}+T_{\mathrm{ACK,ref}}+T_{\mathrm{DIFS}}.
\label{eq:eifs-intro}
\end{equation}
其中 $T_{\mathrm{EIFS}}$ 为扩展等待时间，$T_{\mathrm{ACK,ref}}$ 为按适用物理层规定的参考 ACK 发送时间计算的保护量，其余符号与前文一致。此式为帮助衔接后续内容而补充的时间构成说明，并非本节录音已经完成的推导；这里不为所有 PHY 指定一个共同数值。具体 ACK 参考速率与计时细节应服从对应规范。也不能把任意数据速率下实测的 ACK 时长直接代入所有场景。

再次对照两种状态：ACK 超时由\textbf{发送端}使用，用于判定自己的单播交换未获确认；EIFS 由\textbf{错误接收了某个帧的站点}使用，用于避免干扰别人的后续交换。它们可能都涉及 ACK 所需时间，却不是同一个定时器。本段课程尚未完整展开 EIFS 的恢复时序，因此这里保留其作用与构成，后续章节再处理具体恢复条件。

\subsection{通向 PCF：从分布式竞争到集中式协调}

DCF 通过侦听、间隔、预留与随机化组织竞争。它能适应终端自行产生的发送需求，但访问时间具有随机性：队列中的帧可能经历不同的退避、冻结与重试。对于希望得到更可控服务的业务，课程由此引出另一条路线，即由一个协调者决定各站点何时发送。

点协调功能 PCF 的基本思想是由 AP 中的点协调者（Point Coordinator，PC）组织无竞争的轮询访问，而不是让每个被服务站点都重新随机争用。因为需要这个中心角色，它依赖基础设施模式。前面介绍的 PIFS 比 DIFS 短，使协调者能够在适用的阶段先于普通 DCF 站点取得介质；它仍然不会抢断已经开始的帧，也不会破坏 SIFS 响应。

截至本部分所覆盖的课堂，详细推导的主体是 DCF，PCF 只提供概念入口。轮询次序、无竞争周期与竞争周期的交替，以及实际时序的异常处理应在后续 PCF 讲解中展开。这里应先带走一条完整因果链：\textbf{无线信道共享引出接入控制，半双工引出确认式失败判断，局部侦听引出 NAV 与握手，共同等待引出随机退避，连续成功又引出发送后退避的公平性要求。}

\section{从集中轮询到无线服务质量}
\label{chap:qos}

传统分布式协调功能（Distributed Coordination Function，DCF）回答的是多个站点如何共享信道，却没有直接回答语音分组应当等待多久、视频业务应当获得多少连续发送时间。一个可靠送达但已经错过播放时刻的语音分组，应用价值可能接近于零；后台文件同步则通常能够容忍等待。因而，服务质量（Quality of Service，QoS）的核心不是让所有分组同样快，而是使接入行为与业务需求相匹配。本章从点协调功能的尝试出发，逐步建立差异化竞争、时间授权和集中调度之间的联系。

\subsection{点协调功能：把随机争用改为逐站授权}

点协调功能（Point Coordination Function，PCF）利用接入点（Access Point，AP）中的集中协调者决定下一次由谁发送。与DCF相比，PCF的变化不在于无线链路突然能够避免接收错误，而在于同一受控区域内的站点不再自行争夺每次发送机会。站点必须先得到轮询（Polling）授权，再使用信道。因此，集中控制可以消除被协调站点之间的竞争碰撞，但不能消除噪声、衰落或外部网络造成的干扰。

\subsubsection{超帧与两类接入的共存}

课堂用超帧（Superframe）描述信标（Beacon）之间的宏观时间结构。在这一教学模型中，一个周期先包含无竞争期（Contention-Free Period，CFP），随后进入竞争期（Contention Period，CP）。前者采用PCF，后者采用DCF。两套接入方式在同一无线信道上按时间分工，而不是分别占用两条频率不同的信道。

\begin{figure}[htbp]
\centering
\begin{tikzpicture}[x=0.95cm,y=0.8cm,>=stealth,font=\small]
\draw[->] (0,0)--(13,0) node[right]{时间};
\draw[fill=gray!35] (0.4,0.2) rectangle (0.8,1);
\draw[fill=blue!15] (0.8,0.2) rectangle (5.2,1);
\draw[fill=green!15] (5.2,0.2) rectangle (11.7,1);
\draw[fill=gray!35] (11.7,0.2) rectangle (12.1,1);
\node[above] at (0.6,1){信标};
\node at (3,0.6){CFP：按轮询发送};
\node at (8.5,0.6){CP：按DCF竞争};
\node[above] at (11.9,1){下一信标};
\draw[<->] (0.4,-0.6)--(11.7,-0.6) node[midway,below]{一个示意超帧};
\node[below] at (5.2,0.1){CF-End};
\end{tikzpicture}
\caption{PCF与DCF的时分共存。此图表示机制关系，不把实际信标间隔假定为永远没有延迟的精确周期。}
\end{figure}

参加PCF的意愿在关联阶段提供给AP，AP据此维护轮询列表。进入CFP后，AP按列表发出无竞争轮询（Contention-Free Poll，CF-Poll）。收到轮询的站点如果有数据，便在规定间隔后发送；如果没有数据，仍需用空帧（Null Frame）表明没有待发送内容，使AP能够继续轮询下一站。无竞争结束帧（Contention-Free End，CF-End）结束这一受控阶段，使普通竞争恢复。

一次轮询不必只携带一个孤立的控制指令。若AP正好有下行数据发给被轮询站点，可以在同一帧中组合下行数据与轮询；对上一帧的确认也可以捎带（Piggyback）在后续适当帧中。这样的组合减少单独控制帧的发送次数。它不是省掉确认语义，而是让一份物理传输同时承担数据、确认或授权等多个逻辑功能。

\subsubsection{短间隔维持授权链，PIFS帮助取得控制权}

短帧间间隔（Short Interframe Space，SIFS）使已经开始的对话能够继续。轮询及被轮询站点的响应构成受控交换，相邻适当帧之间使用短间隔，而不是每次重新执行DCF竞争。站点观察到信道刚刚空闲时，并不能立刻加入，因为既有对话有更短的响应等待时间。

AP开始无竞争服务时需要比普通竞争站点更早取得信道。点协调帧间间隔（PCF Interframe Space，PIFS）正是为此设置。设时隙（Slot）长度为$\sigma$，在本章所用传统时序模型中，
\begin{equation}
T_{\mathrm{PIFS}}=T_{\mathrm{SIFS}}+\sigma,
\qquad
T_{\mathrm{DIFS}}=T_{\mathrm{SIFS}}+2\sigma.
\end{equation}
这里$T_{\mathrm{DIFS}}$是分布式帧间间隔（Distributed Interframe Space，DIFS）。因此，$T_{\mathrm{SIFS}}<T_{\mathrm{PIFS}}<T_{\mathrm{DIFS}}$。这个次序同时保留了两种优先关系：正在进行的响应优先于集中协调者的新接入，集中协调者的新接入又优先于普通DCF竞争。

目标信标发送时间（Target Beacon Transmission Time，TBTT）是计划时刻，不是允许AP打断任意帧的命令。若某站点在TBTT前取得信道，其数据与确认交换可能跨过TBTT。AP必须等正在进行的交换结束，再在符合PIFS规则时发送信标。由此产生的信标延迟会推迟CFP的开始。高优先级接入意味着在下一个合法机会先接入，绝不意味着能够抢占已经在空中传播的帧。

\subsubsection{PCF的三项局限来自三个不同的信息缺口}

第一项局限是静态名单与实时业务状态之间的差距。愿意参与PCF的站点不等于此刻确实有数据的站点。如果AP轮询一个空队列，轮询和空响应仍然占用信道，其他有实际数据的站点只能等待。名单解决了“谁可以参加”，却没有解决“谁现在需要服务”。

第二项局限是信标可能被既有DCF交换推迟。即使AP内部已经安排好轮询顺序，受控服务的起始时刻仍可能变化。设理想服务开始时刻为$t_b$，信标推迟量为$\Delta_b$，则实际起点为$t_b+\Delta_b$。当$\Delta_b$不固定时，后续站点的服务开始时间也整体移动。周期性产生的语音分组由此可能遭遇不稳定的接入等待。

第三项局限是轮询后的占用时间不容易事先精确掌握。若站点$i$发送有效载荷$L_i$比特、MAC开销$H_i$比特，数据部分速率为$R_i$，物理层固定开销为$P_i$秒，则一个简化的数据帧时长为
\begin{equation}
T_i=P_i+\frac{H_i+L_i}{R_i}.
\end{equation}
即使$L_i$相同，速率$R_i$降低也会拉长服务时间。若第$j$站要等前$j-1$站服务完毕，其等待量含有$\sum_{i=1}^{j-1}T_i$以及轮询、响应等间隔。AP只有名单而没有充分的队列与时长信息时，很难据此做精细的时延规划。

这三项局限共同说明：集中轮询本身并不自动等于确定的QoS。要改善服务，必须进一步获取业务状态、约束每次占用的时间，并让受控服务的出现时机更灵活。

\subsection{混合协调功能与传输机会}

IEEE 802.11e引入混合协调功能（Hybrid Coordination Function，HCF）。其中，增强型分布式信道接入（Enhanced Distributed Channel Access，EDCA）保留竞争方式，受控信道接入（HCF Controlled Channel Access，HCCA）提供集中轮询方式。两者共享一个关键抽象：传输机会（Transmission Opportunity，TXOP）。支持QoS的站点称QoS站点（QoS Station，QSTA），相应基本服务集称QoS基本服务集（QoS Basic Service Set，QBSS）。HCCA的集中控制者称混合协调器（Hybrid Coordinator，HC），本章取其位于AP中的常见架构。

\subsubsection{从“一次接入一帧”到“一次接入一段时间”}

在前面讨论的简化DCF模型中，一次成功竞争通常用来完成一帧数据及其确认交换。TXOP则赋予站点在一个有限时间窗口内继续发送的权利。它由开始时刻与持续时间上限确定，后者称TXOP上限（TXOP Limit）。在上限允许且规则满足的条件下，站点可以连续发送多个帧，而不必在每帧之前重新竞争。

时间授权有两个直接作用。对发送者而言，多帧可以分摊一次竞争的成本；对其他站点而言，一次赢家能够连续占用的时间受到限制。站点必须检查下一次交换能否在剩余授权内完成，而不能仅因为数据帧开始时仍在TXOP内，就无条件把其余交换拖到授权结束之后。

TXOP上限是允许使用的最大值，不代表每次必然用满。队列为空、目的站不可用或剩余时间不足以完成下一次交换，都可能使发送者提前结束。因而，“相同TXOP上限”与“实际每次占用完全相同”是两个不同命题；后面的公平性推导会明确采用用满时间的理想化假设。

\subsection{EDCA：用四个虚拟站点表达四类业务}

\subsubsection{业务分类先于接入竞争}

一个传统单队列MAC把不同业务放进同一个先入先出队列（First-In, First-Out，FIFO）。队首若排着大批文件数据，后来到达的语音分组便可能被挡在后面。EDCA把流量划分到四个接入类别（Access Category，AC），每个类别分别维护队列和竞争状态。语音类别$\mathrm{AC}_{\mathrm{VO}}$通常具有最高相对优先级，视频类别$\mathrm{AC}_{\mathrm{VI}}$次之，尽力而为类别（Best Effort）$\mathrm{AC}_{\mathrm{BE}}$适合普通交互数据，后台类别（Background）$\mathrm{AC}_{\mathrm{BK}}$适合更能容忍等待的业务。

这种分类来自需求差异。语音分组通常较短，但到达太迟就无法用于实时播放；视频流需要较高吞吐量，播放缓冲能吸收一部分到达时延，但缓冲耗尽仍会造成卡顿；网页需要合理的响应时间；后台同步往往可以让出即时机会。四类并不是对应用重要性的永久评价，而是一组使业务需求能够进入MAC调度的接口。

\begin{figure}[htbp]
\centering
\begin{tikzpicture}[x=1cm,y=0.85cm,>=stealth,font=\small]
\node[draw,rounded corners,minimum width=2.1cm,minimum height=0.7cm] at (5,4.1){上层业务分类};
\foreach \x/\n in {0.8/VO,3.6/VI,6.4/BE,9.2/BK}{
\draw[->] (5,3.7)--(\x,3.1);
\node[draw,fill=blue!10,minimum width=2.2cm,minimum height=0.7cm] at (\x,2.7){AC\_\n 队列};
\node[draw,align=center,minimum width=2.2cm,minimum height=1cm] at (\x,1.3){独立AIFS\\CW与退避计数};
\draw[->] (\x,2.3)--(\x,1.9);
\draw[->] (\x,0.7)--(5,-0.1);
}
\node[draw,fill=green!15,minimum width=5cm,minimum height=0.7cm] at (5,-0.5){内部仲裁：选择一个类别发送};
\draw[->] (5,-0.9)--(5,-1.5) node[below]{同一个无线收发器};
\end{tikzpicture}
\caption{四个“虚拟站点”是四套竞争状态，不是四部可同时发送的无线电。}
\end{figure}

\subsubsection{AIFS决定何时有资格开始倒计数}

EDCA仍使用载波侦听多路访问及碰撞避免（Carrier Sense Multiple Access with Collision Avoidance，CSMA/CA）的基本逻辑。优先级不是通过取消低优先级站点的竞争权实现，而是通过不同参数改变获胜概率。

仲裁帧间间隔（Arbitration Interframe Space，AIFS）取代统一的DIFS。对类别$a$，
\begin{equation}
T_{\mathrm{AIFS},a}=T_{\mathrm{SIFS}}+\mathrm{AIFSN}_a\sigma,
\end{equation}
其中$\mathrm{AIFSN}_a$为该类别的仲裁帧间间隔时隙数（Arbitration Interframe Space Number）。信道从忙变为空闲后，AIFS较短的类别较早获得倒计数资格。若较高优先级类别已经在较低优先级类别尚未完成AIFS时开始发送，后者只能继续等待。

例如，取课堂时序示意中的$T_{\mathrm{SIFS}}=16\,\mu\mathrm{s}$、$\sigma=9\,\mu\mathrm{s}$。若两个类别分别使用$\mathrm{AIFSN}=2$和3，则其AIFS分别为34与43微秒。这个9微秒差值不是给高优先级队列增加发送速率，而是在每次空闲竞争起点提供一个时隙的先行机会。

\subsubsection{竞争窗口决定随机等待的分布}

每个类别另有自己的最小竞争窗口（Minimum Contention Window）与最大竞争窗口（Maximum Contention Window），分别记为$CW_{\min,a}$与$CW_{\max,a}$。为避免窗口“长度”与“上界”混淆，本节把$CW_a$定义为可取退避整数的最大值，即
\begin{equation}
B_a\sim\operatorname{Unif}\{0,1,\ldots,CW_a\}.
\end{equation}
因此，可能取值共有$CW_a+1$个，而平均退避整数为
\begin{equation}
\mathbb{E}[B_a]=\frac{0+1+\cdots+CW_a}{CW_a+1}=\frac{CW_a}{2}.
\end{equation}
在没有其他发送打断、没有碰撞且已经有分组等待的理想条件下，一次开始竞争后的平均等待可写成
\begin{equation}
\mathbb{E}[W_a]=T_{\mathrm{AIFS},a}+\frac{CW_a}{2}\sigma.
\end{equation}
若$CW_a=7$、$\sigma=9\,\mu\mathrm{s}$，退避部分平均为31.5微秒；若$CW_a=15$，则为67.5微秒。真实等待还包括计数冻结、重新满足空闲间隔、碰撞重试和队列排队，上式不是完整网络时延模型。

失败后窗口按照该类别的范围增大。用上述上界记法，一次典型二进制增长写为
\begin{equation}
CW'_a=\min\{2(CW_a+1)-1,\ CW_{\max,a}\}.
\end{equation}
加一再翻倍是因为倍增的是可选整数的数量。成功后，窗口按协议回到该类别的初始设置。高优先级类别通常使用较小窗口，因而统计上更快接入，但仍可能抽到较大随机数，也可能与其他站点碰撞。EDCA提供的是统计差异化，不能仅凭较短AIFS和较小窗口就宣称每个语音分组都有严格时延上界。

\subsubsection{内部碰撞与外部碰撞}

同一物理站点中的两个AC可能在同一时刻把计数器减到零。这称内部碰撞或虚拟碰撞（Internal/Virtual Collision）。由于站点只有一套要使用的无线发送资源，它不会让两个队列真的同时在空中碰撞。内部仲裁把发送机会给较高优先级类别，未获机会的类别按相应失败退避规则处理。外部站点只能看到胜出的那个类别实际发送。

因此，一个站点内部需要维护四套逻辑竞争过程，但外部信道依然只有一段实际发送。内部仲裁解决的是本机多个候选者同时就绪的问题，不能解决不同物理站点在同一时隙到零造成的外部碰撞。

\subsubsection{TXOP上限应与业务形态匹配}

第三组参数是每个类别的TXOP上限。较长TXOP让一次竞争换来更多有效数据，却也让其他类别和其他站点等待更久。最高优先级并不意味着必须具有最长TXOP：语音急需尽早接入，但一批语音数据可能很小；视频持续产生较多数据，一次获准后可能更需要连续发送时间。

这样，三个参数分别作用于不同问题：AIFS改变什么时候开始竞争，CW改变开始竞争后的随机等待分布，TXOP上限改变获胜后允许占用多久。理解其分工比记住某张参数表更重要。具体参数随物理层、接入点配置及所用协议条件而变化，不能把一组课堂示例数值视为所有Wi-Fi设备不变的常量。

QoS数据帧中增加QoS控制字段（QoS Control）。其中的业务标识（Traffic Identifier，TID）及确认策略等信息使接收端能够识别业务处理方式；在相应帧格式和方向下，字段也可表达队列或授权相关信息。这些位的解释依上下文而定，不能把同一组比特同时当成独立、并存的全部信息字段。课堂的重点是：MAC不仅运送数据，还能向协调者提供安排后续服务所需的状态。

\subsection{突发发送与块确认}

\subsubsection{一次竞争如何支撑多个数据帧}

突发发送（Burst Transmission）在一个TXOP内发送多个帧，减少重复争用。若仍逐帧确认，发送者需要完成“数据、SIFS、确认”的多次交换。块确认（Block Acknowledgment，Block ACK）则把多个帧的接收结果集中报告，减少确认开销。课堂讨论的即时块确认流程为：发送多帧数据，发块确认请求（Block ACK Request，BAR），接收方经SIFS后返回块确认。

\begin{figure}[htbp]
\centering
\begin{tikzpicture}[x=0.95cm,y=0.8cm,>=stealth,font=\small]
\node[left] at (0,1.4){发送方};\node[left] at (0,0){接收方};
\draw[->] (0,1.4)--(12.5,1.4);\draw[->] (0,0)--(12.5,0);
\foreach \l/\r/\s in {0.4/2.4/数据1,2.9/4.9/数据2,5.4/7.4/数据3,7.9/9.2/BAR}{
\draw[fill=blue!15] (\l,1.15) rectangle (\r,1.95);\node at ({(\l+\r)/2},1.55){\s};}
\draw[fill=green!20] (9.8,-0.25) rectangle (12,0.55);\node at (10.9,0.15){Block ACK};
\draw[->] (9.2,1.15)--(9.8,0.55);\node[above] at (9.65,1.65){SIFS};
\draw[<->] (0.4,-0.85)--(12,-0.85) node[midway,below]{一个TXOP；图中各必要短间隔不按比例绘制};
\end{tikzpicture}
\caption{课堂中的多帧突发与即时块确认流程。该流程是一个具体模型，并非所有版本的块确认都必须采用完全相同的请求时机。}
\end{figure}

块确认中的位图（Bitmap）依据约定的起始序号与窗口表示多个帧的状态。假设连续四帧的状态为$1,0,1,1$，则发送者知道第二帧未被确认，而第一、第三、第四帧已被确认；后续可以针对未成功的部分安排重传。位图位的位置必须与序号对应，不能脱离窗口把一个比特解释成任意分组的确认。课堂“一帧一位”的图式对应常用的简化理解，不应据此推断所有块确认变体都有相同位图格式。

\subsubsection{减少开销的收益与代价}

设突发包含$m$个等长数据帧，每帧空中时长为$T_D$，逐帧确认时长为$T_A$，BAR与块确认的时长分别为$T_R$与$T_B$。令$S$为SIFS，暂不计TXOP外的首次竞争。若逐帧交换之间也用SIFS衔接，则总时长近似为
\begin{equation}
T_{\mathrm{individual}}=mT_D+mT_A+(2m-1)S.
\end{equation}
课堂BAR模型中的$m$帧数据、BAR和块确认共$m+2$个帧，相邻之间有$m+1$个短间隔，因此
\begin{equation}
T_{\mathrm{block}}=mT_D+T_R+T_B+(m+1)S.
\end{equation}
两式相减得到
\begin{equation}
\Delta T=mT_A-T_R-T_B+(m-2)S.
\end{equation}
当$\Delta T>0$时，这个简化模型下块确认减少了空中时间。由于BAR和块确认也有自身开销，不能仅凭“只回一次确认”就断言任意小批量都一定更省时。

另一方面，较长连续突发可能增大其他站点的等待与时延抖动（Delay Jitter），接收端还需保存多个接收状态，发送端需维护未确认窗口。提高效率往往来自一次处理更多数据，但实时系统还要约束数据在组批和等待期间增加的时延。

\subsection{分组公平、时间公平与吞吐量}

\subsubsection{课堂明确提示可能考查：先固定比较对象与假设}

教师在10月13日08:03场次的01:01:14明确提示这一吞吐量比较“可能出现在测试中”。复习重点是从成功交付比特数推导公式，说明相等条件，而不是把两条式子脱离假设背诵。

公平性（Fairness）必须先说明公平的对象：是成功接入次数、帧数、比特数，还是占用时间。本节比较两个饱和站点（Saturated Station）A与B，即两者始终有数据等待。设它们竞争参数相同，在足够长观察期$T$内各成功获得$N$次机会；忽略有限样本偏差，并把不同错误率造成的不对称排除在外。

在一次成功机会只发送一帧的DCF简化模型中，A和B各送出$N$帧。令每帧有效载荷分别为$L_A$、$L_B$比特，吞吐量（Throughput）定义为观察期内成功交付有效比特数除以观察期长度，故
\begin{equation}
G_A=\frac{NL_A}{T},\qquad G_B=\frac{NL_B}{T},\qquad
\frac{G_A}{G_B}=\frac{L_A}{L_B}.
\end{equation}
所以，相同成功接入次数不自动产生相同吞吐量；只有此处的有效载荷相同时，两者吞吐量才相同。若载荷相同而速率不同，慢站发送同一帧会占用更长时间。快站也要等待慢站结束，整个系统的机会循环被拉长。

\subsubsection{相同时间授权改变了公平性的基准}

现在令两站使用同一个AC，获得相同可用发送时间$\tau$，并假设每次都用满。若先忽略开销，A的发送速率为$R_A$比特每秒，则一次机会发送$\tau R_A$比特。于是
\begin{equation}
G_A\simeq\frac{N\tau R_A}{T},\qquad
G_B\simeq\frac{N\tau R_B}{T},\qquad
\frac{G_A}{G_B}\simeq\frac{R_A}{R_B}.
\end{equation}
两站拥有相同时间，但更高的速率可以把这段时间转化为更多比特。若$R_A=2R_B$，则理想吞吐量也约为两倍。这不违背时间公平（Temporal Fairness），因为其平等对象是时间而非字节数。

\begin{figure}[htbp]
\centering
\begin{tikzpicture}[x=0.95cm,y=0.8cm,font=\small,>=stealth]
\node[left,align=right] at (0,1.5){每次一帧\\相同载荷};
\node[left,align=right] at (0,-0.2){每次等时\\相同TXOP};
\draw[->] (0,1.1)--(12.4,1.1);\draw[->] (0,-0.6)--(12.4,-0.6);
\foreach \l/\r/\s in {0.3/1.8/A,2.2/5.2/B,5.7/7.2/A,7.6/10.6/B}{\draw[fill=blue!15] (\l,1.2) rectangle (\r,1.9);\node at ({(\l+\r)/2},1.55){\s};}
\foreach \l/\r/\s in {0.3/2.6/A,3/5.3/B,5.7/8/A,8.4/10.7/B}{\draw[fill=green!15] (\l,-0.5) rectangle (\r,0.2);\node at ({(\l+\r)/2},-0.15){\s};}
\node[below,align=center] at (6,-1){上排：慢站B占用更久；下排：等时窗口内A可发送更多数据};
\end{tikzpicture}
\caption{令A速率高于B时，两种公平基准的差异。图中交替顺序仅为示意，随机接入并不保证严格轮流。}
\end{figure}

作为数值比较，设观察期$T=1$ s，两站各成功接入$N=100$次。DCF情形若$L_A=1000$字节、$L_B=1500$字节，则$G_A=100\times8000/1=0.8$ Mb/s，$G_B=100\times12000/1=1.2$ Mb/s。它们发送帧数相同，但B的每帧载荷更大，吞吐量也更高。

另看相同TXOP的独立理想场景：令每次可用时间$\tau=1$ ms，$R_A=12$ Mb/s、$R_B=6$ Mb/s。每次A可发送12000 bit，B可发送6000 bit；在同样$N=100$次、$T=1$ s的观察条件下，$G_A=1.2$ Mb/s、$G_B=0.6$ Mb/s。两站各累计使用0.1 s的授权时间，吞吐量却不相等。前一情形比较载荷大小，后一情形比较时间内的发送能力；两个场景不要求具有相同的实际竞争过程。

若考虑开销，可把$R_i$替换为窗口内平均有效载荷速率$r_i^{\mathrm{eff}}$，从而写成$G_i=N\tau r_i^{\mathrm{eff}}/T$。此外，不同帧长度可能使最后一帧无法装进剩余时间，队列也可能提前为空。因此，课上的$N\tau R_i/T$是解释公平基准的理想模型，而不是直接使用物理层标称速率就能精确预测的实测吞吐量公式。

\subsection{HCCA：让轮询掌握需求、时长与时机}

\subsubsection{动态队列信息与有限时间授权}

HCCA针对PCF的第一项问题，引入更充分的业务状态。站点可以报告队列需求，使HC知道谁有数据以及大致需要多少服务。于是轮询不必机械遍历一张长期不变的参与名单。信息并非永远即时准确：一次报告之后仍可能到达新分组，因而具体调度必须考虑状态陈旧和新增业务。

针对第二类信息不足，即发送占用时长未知，HC在QoS轮询中授予有限TXOP。它不只是说“轮到你”，还限制“你可以使用多久”。令站点$i$获准时间为$\tau_i$，轮询和必要转换开销为$c_i$。若一个调度周期中依次服务$m$个站点，其已计划时长为
\begin{equation}
T_{\mathrm{plan}}=\sum_{i=1}^{m}(c_i+\tau_i).
\end{equation}
这使调度者能进行时间预算。它并不证明所有分组一定按时交付，因为重传、外部干扰与突发到达仍可能破坏预期，但相比不知道每站要发多久，至少有了可管理的资源单位。

\subsubsection{受控接入阶段可以出现在竞争期内}

PCF的服务主要受固定CFP结构约束。HCCA进一步允许HC在竞争期中，于信道空闲达到PIFS后发起受控接入阶段（Controlled Access Phase，CAP）。因此，一个CP内部既可以出现EDCA竞争获得的TXOP，也可以出现HC轮询授予的TXOP。

\begin{figure}[htbp]
\centering
\begin{tikzpicture}[x=0.95cm,y=0.9cm,font=\small,>=stealth]
\draw[->] (0,0)--(13,0) node[right]{时间};
\draw[fill=gray!30] (0.3,0.2) rectangle (0.7,1);
\draw[fill=blue!15] (0.7,0.2) rectangle (3.6,1);
\draw[fill=green!15] (3.6,0.2) rectangle (6.3,1);
\draw[fill=blue!15] (6.8,0.2) rectangle (9.4,1);
\draw[fill=green!15] (9.4,0.2) rectangle (12.2,1);
\node[above] at (0.5,1){信标};\node at (2.1,0.6){CFP轮询};
\node at (4.95,0.6){EDCA};\node at (8.1,0.6){CAP轮询};\node at (10.8,0.6){EDCA};
\draw[<->] (3.6,-0.55)--(12.2,-0.55) node[midway,below]{竞争期CP并不排斥受控服务};
\node[above,align=center] at (6.55,1.1){PIFS后\\HC发起};
\end{tikzpicture}
\caption{HCCA可在CP内插入CAP。协调者先等已有交换结束；“插入”不表示打断正在发送的帧。}
\end{figure}

这使得信标延迟不再意味着所有受控服务只能一并等到下一个固定窗口。HC有机会在后续空闲点补充服务，尤其适合处理紧迫需求。不过，若信道一直被合法交换或外部干扰占用，PIFS条件仍然无法满足。HCCA增加了控制手段，而没有取消无线共享介质的物理约束。

\subsubsection{流量协商与接纳控制}

只报告队列长度是状态通知；协商业务所需吞吐量、时延和容许丢失则是资源需求声明。流量规范（Traffic Specification，TSPEC）提供后一类描述。课堂介绍通过添加流量流（Add Traffic Stream，ADDTS）的请求与响应来协商服务，并用删除流量流（Delete Traffic Stream，DELTS）结束相应状态。

接纳控制（Admission Control）的任务是判断是否应接受一个新业务。即使每个站点都使用高优先级类别，如果所有业务需求的总和超过可用空中时间，系统仍然无法同时满足所有要求。因此，QoS不只是把高优先级标签附在数据上，还包括必要时拒绝新的过量负载。

作为预算思想，设业务$i$每隔$P_i$秒至少需要$\tau_i$秒有效服务及$c_i$秒相关开销，则其时间利用率为$u_i=(\tau_i+c_i)/P_i$。若保留一部分比例$\rho$给其他业务及不确定性，一个必要的粗略条件是
\begin{equation}
\sum_i u_i\leq 1-\rho.
\end{equation}
所有量都用同一时间单位，$u_i$与$\rho$均无量纲。这个不等式只是总预算检查，不能代替具有具体截止期、重传和突发约束的可调度性证明。标准提供信令与授权机制，如何具体决定轮询顺序和服务频率仍需要调度算法。

\subsection{本章小结}

PCF表明，集中发号施令并不足以保证服务质量：调度者还需要实时业务信息、可预算的占用时间和合适的服务时机。EDCA用四类队列以及AIFS、CW和TXOP上限提供统计优先级；HCCA则结合状态、协商和有限时间授权进行集中服务。TXOP把接入资源从“一次一帧”提升为“一次一段时间”，块确认帮助摊薄确认成本。最终，无论竞争还是集中控制，都必须同时考虑吞吐量、时延、公平性和额外控制开销。

\section{高速物理层为何仍需要帧聚合}
\label{chap:aggregation}

\subsection{从物理层能力到有效吞吐量}

课堂在10月13日08:03场次的01:11:16和01:26:25将额外物理层细节作为了解内容。这里保留机制解释，重点是它们如何影响MAC效率，不把具体速率组合表当成需要机械记忆的重点。

物理层（Physical Layer，PHY）数据速率描述特定发送模式每秒能够承载多少比特，应用吞吐量则衡量经过竞争、头部、确认和重传后，每秒真正交付多少有效数据。提高前者并不保证后者等比例提高。理解这种差异，是理解IEEE 802.11n介质访问控制层（Medium Access Control，MAC）增强的起点。

\subsubsection{速率提升的几个来源}

信道绑定（Channel Bonding）把相邻信道组合成更宽的工作带宽。例如，从20 MHz扩展到40 MHz可以提供更多并行频域资源，但也占用更多频谱，并要求存在适合的信道条件。更宽带宽不是免费资源，邻近网络的共存和干扰同样需要考虑。

调制编码方案（Modulation and Coding Scheme，MCS）组合调制阶数与编码率。若一个调制符号可取$M$种状态，则每个符号可标记$\log_2 M$比特；编码率$r_c$表示传输比特中属于信息的比例。提高调制阶数或编码率会提升净速率，同时通常要求更好的接收条件。正交频分复用（Orthogonal Frequency Division Multiplexing，OFDM）把数据分布到多个正交子载波；多输入多输出（Multiple Input Multiple Output，MIMO）则可利用空间维度同时传输多个空间流（Spatial Stream）。

用一个仅用于解释因素的表达式，可以写成
\begin{equation}
R\approx\frac{N_{\mathrm{SS}}N_{\mathrm{SD}}\log_2(M)r_c}{T_{\mathrm{sym}}}.
\end{equation}
其中$N_{\mathrm{SS}}$为空间流数量，$N_{\mathrm{SD}}$为每个流用于承载数据的子载波数，$T_{\mathrm{sym}}$为包含保护间隔的符号周期。分子为每个符号周期承载的信息比特数，除以秒得到比特每秒。实际速率表还受到编码、可用模式及物理层格式限制，不能仅用此式任意组合参数。

波束成形（Beamforming）通过多个天线信号的相对相位和幅度控制，使接收方向上的信号更有利地叠加。课堂用“朝目标方向集中表达”帮助理解其作用。它主要改善接收信号条件，不等于每增加一副天线就必然多出一条独立数据流，也不等于发射功率可以无限提高。

802.11n的600 Mb/s是特定多空间流、带宽、保护间隔与调制编码组合下的最高物理层速率。速率表中的6.5 Mb/s与600 Mb/s是不同配置下的例子，不能把前者一概当成所有802.11n前导码、控制帧和头部都统一采用的“基本速率”。各字段的发送方式应由具体物理层格式决定。

\subsubsection{数据帧内的各部分不一定使用相同速率}

物理层协议数据单元（PHY Protocol Data Unit，PPDU）包含前导码（Preamble）、物理层头部及后续承载MAC内容的部分。MAC协议数据单元（MAC Protocol Data Unit，MPDU）又包含MAC头部、有效载荷和帧校验序列（Frame Check Sequence，FCS）。前导码用于接收机检测、同步及其他物理层处理；头部使接收机知道如何解释后面的比特。即使应用只关心载荷，也不能省去这些必要结构。

设前导码和物理头部共占固定时间$P$，MAC头部及校验等共$H$比特，应用有效载荷为$L$比特，承载MAC部分的速率为$R$，则
\begin{equation}
T_D=P+\frac{H+L}{R}.
\end{equation}
这里$P$直接以秒计量，是为了避免把特殊编码的前导部分错误地当成普通数据比特除以$R$。实际OFDM帧还可能需要服务位、尾比特、填充与完整符号取整；本式用于本课的机制分析，而非替代完整PHY时长计算器。

若一次成功发送还需平均竞争开销$C$、一个SIFS及确认时长$T_A$，则
\begin{equation}
T_{\mathrm{cycle}}=C+P+\frac{H+L}{R}+T_{\mathrm{SIFS}}+T_A,
\qquad G=\frac{L}{T_{\mathrm{cycle}}}.
\end{equation}
$G$是每个周期成功交付的有效比特数除以周期时间，故单位为bit/s；相对物理速率的效率定义为$\eta=G/R$。

\subsection{一个11 Mb/s计算与高速率极限}

课件的传统11 Mb/s算例给定：有效载荷1500字节，MAC开销34字节，物理前导及头部时长192微秒；DIFS为50微秒，平均退避为310微秒，SIFS为10微秒；ACK为14字节，在1 Mb/s下发送并具有相同的192微秒物理开头。这是一个明确指定速率和格式的示例，不是所有确认帧的通用参数。

先把字节换成比特，得到$L=1500\times8=12000$ bit、$H=34\times8=272$ bit。由于1 Mb/s等于每微秒1比特，使用bit与Mb/s相除可直接得到微秒。因此
\begin{align}
T_D&=192+\frac{12000+272}{11}=1307.64\,\mu\mathrm{s},\\
T_A&=192+\frac{14\times8}{1}=304\,\mu\mathrm{s},\\
T_{\mathrm{cycle}}&=50+310+1307.64+10+304
=1981.64\,\mu\mathrm{s},\\
G&=\frac{12000}{1981.64}\approx6.06\,\mathrm{Mb/s},\qquad
\eta\approx55.1\%.
\end{align}
这说明11 Mb/s是承载数据部分时采用的速率，并非应用每秒能连续得到11百万有效比特。一次1500字节发送还支付了侦听、退避、物理头部和确认等成本。

把那些不随$R$缩短的开销合并成$T_0$，则
\begin{equation}
G(R)=\frac{L}{T_0+(H+L)/R}
=\frac{LR}{RT_0+H+L},\qquad
\eta(R)=\frac{L}{RT_0+H+L}.
\end{equation}
在固定$L,H,T_0>0$的模型中，$R$增加会提高$G$，但效率$\eta$却下降；当$R$趋于无穷大时，
\begin{equation}
\lim_{R\to\infty}G(R)=\frac{L}{T_0},\qquad
\lim_{R\to\infty}\eta(R)=0.
\end{equation}
物理传送再快，也无法压缩模型中的固定等待成本。这里的极限不是说高速率无用，而是说只提高速率会出现边际收益递减。要继续提高有效吞吐量，需要增大每次支付固定开销所能运送的数据量，这就引出帧聚合（Frame Aggregation）。

\subsection{两级聚合的边界与可靠性}

\subsubsection{A-MSDU：先合并高层交付的数据}

MAC服务数据单元（MAC Service Data Unit，MSDU）是上层交给MAC处理的数据。聚合MSDU（Aggregated MSDU，A-MSDU）把多个适合共同发送的MSDU放入一个MPDU载荷中，从而共享一个外层MAC头和FCS。内部每个子帧仍需携带必要的地址、长度等子帧头，必要时还需填充（Padding）；聚合并不意味着简单地把纯应用字节无分隔地串起来。

其关键收益是减少重复MAC封装。关键代价则是多个子帧受同一个外层MPDU校验保护：若该MPDU不能正确接收，其内部数据不能像独立MPDU那样分别用块确认逐项恢复。聚合越大，分摊成本越好，但在错误信道上一次失败牵连的数据也越多。

\subsubsection{A-MPDU：保留多个MAC帧的独立边界}

聚合MPDU（Aggregated MPDU，A-MPDU）把多个已经形成的MPDU放进一次物理层发送中，共享前导码及物理层头部。各MPDU仍保留自己的MAC头和FCS；分隔符（Delimiter）与填充帮助接收机识别边界。由于保留独立校验单元，接收方可以使用块确认指出哪些MPDU成功，发送方再选择性重传未成功部分。

因此，A-MSDU与A-MPDU节省的层次不同：前者减少多个MSDU分别封装成MPDU时的外层MAC开销；后者减少多个MPDU分别形成物理发送时的PHY开销。两者都不能任意合并所有目的站和业务。实际聚合需满足接收方、业务标识及已建立确认关系等限制；本章采用同一接收方、同一适用业务流的常规情形。

\begin{figure}[htbp]
\centering
\begin{tikzpicture}[x=0.95cm,y=0.85cm,font=\small,>=stealth]
\node[left] at (0,2.8){A-MSDU};
\draw[fill=gray!25] (0.1,2.4) rectangle (1.7,3.2);\node at (0.9,2.8){MAC头};
\foreach \l/\r/\s in {1.7/4.5/子帧1,4.5/7.3/子帧2,7.3/10.1/子帧3}{\draw[fill=blue!15] (\l,2.4) rectangle (\r,3.2);\node at ({(\l+\r)/2},2.8){\s};}
\draw[fill=gray!25] (10.1,2.4) rectangle (11.2,3.2);\node at (10.65,2.8){FCS};
\node[below] at (5.5,2.35){一个MPDU：整个外层帧共享校验};
\node[left] at (0,0.5){A-MPDU};
\draw[fill=gray!25] (0.1,0.1) rectangle (2,0.9);\node[align=center] at (1.05,0.5){PHY开头};
\foreach \l/\r/\s in {2.3/4.9/MPDU 1,5.2/7.8/MPDU 2,8.1/10.7/MPDU 3}{\draw[fill=green!15] (\l,0.1) rectangle (\r,0.9);\node at ({(\l+\r)/2},0.5){\s};\draw[fill=gray!60] ({\l-0.3},0.1) rectangle (\l,0.9);}
\node[below,align=center] at (6,-0.15){每个MPDU保留MAC头与FCS；窄灰块表示分隔结构};
\draw[->,thick] (5.5,1.6)--(6.5,1.05);
\end{tikzpicture}
\caption{两级聚合的封装边界。一个A-MSDU可以成为下层某个MPDU的载荷；图不按字节长度比例绘制。}
\end{figure}

\subsubsection{两级聚合与组批等待之间的折中}

两级聚合（Two-Level Aggregation）先把若干MSDU组成A-MSDU，再将携带这些A-MSDU的多个MPDU组成A-MPDU。这样既能分摊MAC外层头部，又能分摊PHY开头。但实现仍需要决定何时停止组批并发送，不能为追求最大聚合而无限等候。

用一个不含误码的简化模型说明收益。设每个子数据有效长度为$L$比特，其聚合后仍需保留的平均开销为$h$比特，每次物理发送共享固定成本$T_0$，则聚合$n$份数据的吞吐量为
\begin{equation}
G_n=\frac{nL}{T_0+n(L+h)/R}
=\frac{L}{T_0/n+(L+h)/R}.
\end{equation}
当$n$增大时，每份有效数据承担的固定时间$T_0/n$下降；但$n$很大时，$G_n$仍受到$R L/(L+h)$限制。共享成本可以分摊，逐份保留的子帧结构不能凭空消失。

误码还带来另一种限制。若为解释趋势而假定各比特独立出错，误码率为$p$，一个长度$\ell$比特的校验单元全部正确的概率为
\begin{equation}
P_{\mathrm{ok}}=(1-p)^\ell.
\end{equation}
在$p>0$时，校验单元越长，一次完全正确的概率越低。真实无线误码通常具有相关性，此式不是实测信道的完整模型，但清楚说明了“尽可能大”未必总是最优。A-MPDU保留独立MPDU边界的价值，正是在分摊物理开销的同时，让失败恢复仍有较细粒度。

若数据到达率为$\lambda$份每秒，等待凑足$n$份也可能增加组批时延。在理想匀速到达、凑满即发的批次中，第一个数据要等$(n-1)/\lambda$，最后一个无需等候，平均组批等待为$(n-1)/(2\lambda)$。真实系统常用数量上限与等待计时共同限制聚合，从而在效率和实时性之间折中。

\subsection{机制之间的联系}

TXOP、块确认与聚合不是同一个概念。TXOP规定谁能使用信道以及能使用多久；块确认规定如何报告多个数据单元的接收结果；聚合规定把多个数据单元装进怎样的MAC或PHY结构。一次TXOP内可以安排多次物理发送，单次物理发送也可能包含多个MPDU。把这三个层次分开，才能正确画时序、统计开销，并理解高速物理层为何需要配套的MAC设计。

\section{DCF综合时序：从局部观察推出全局顺序}
\label{chap:dcf-examples}

本章沿用课堂两道综合练习，重点是把“每个站点观察到了什么”转化为“什么时候可以递减计数器”。所有站点彼此可听见，数据都发往AP，忽略传播时延，不考虑捕获效应（Capture Effect），即重叠传输不使其中某个帧意外被成功解码。数据发送时长已包含数据帧的各层开销，不再另加前导码或MAC头部。

\subsection{统一记号与时间线画法}

用$D$表示DIFS，$S$表示SIFS，$E$表示扩展帧间间隔（Extended Interframe Space，EIFS），$\sigma$表示时隙长度，$A_{t}$表示ACK发送时长，$F$表示从自身数据发送结束起计算的确认超时（ACK Timeout）等待量。令$Q=S+A_{t}$，则一个数据帧发送完毕之后，成功交换还需$Q$时间。网络分配向量（Network Allocation Vector，NAV）保护已知交换的剩余占用时间。

画图时首先标出实际发送的区段，再逐站判断物理忙、可否解码以及NAV状态。没有收到ACK只说明发送者知道失败；旁听者收到损坏帧才是另一个局部事件。半双工（Half-Duplex）发送者在自己发送期间不能同时观察碰撞。因此，几个站点面对同一次空中重叠，可能采取不同的后续等待。

EIFS结束并不意味着自动获得发送权。若站点已经处于需要退避的状态，它仍必须完成退避；EIFS在相应错误恢复条件下替代普通空闲等待，而不是再机械叠加一次DIFS。若后来正确接收了帧，按正常载波侦听和恢复规则继续，不应永久保留错误惩罚。

\subsection{例一：等长碰撞与随后到达的新分组}

\subsubsection{已知条件与第一次成功发送}

取$t_0=0$。A与B同时发送，二者数据帧均持续1 ms；C的数据帧持续1.5 ms。A与B重试所选退避整数分别为10与11。C在1.25 ms才从上层得到一个新分组。原题截图给定$D=50\,\mu\mathrm{s}$、$S=10\,\mu\mathrm{s}$、$E=200\,\mu\mathrm{s}$、$\sigma=20\,\mu\mathrm{s}$、$A_t=140\,\mu\mathrm{s}$和$F=150\,\mu\mathrm{s}$，所以$Q=150\,\mu\mathrm{s}$。以下统一以毫秒进行绝对时刻计算。

在$[0,1]$ ms内，A、B在发送，AP和C看到相互重叠的信号，无法正确解码。AP不能发ACK。A和B在自身发送结束后等确认，超时后准备重试。C没有发送，它观察到错误接收，因而从1 ms信道空闲起等待EIFS。于是C等待到$1+0.2=1.2$ ms。A、B在$1+0.15=1.15$ ms确认失败，随后等待0.05 ms的DIFS，从1.20 ms开始递减各自计数器。

C的分组到达时，EIFS已经结束。按课堂采用的初次发送简化规则，C观察一个DIFS后发送，因此其开始与结束时刻分别为
\begin{equation}
c=1.25\,\mathrm{ms}+D,
\qquad d=c+1.5\,\mathrm{ms}+Q.
\end{equation}
$d$是C的ACK结束时刻。A和B原先从$1\,\mathrm{ms}+F+D$开始倒计数。利用课上明确给出的5槽递减，有
\begin{equation}
c-(1\,\mathrm{ms}+F+D)=5\sigma,
\qquad F+5\sigma=0.25\,\mathrm{ms}.
\end{equation}
因此，在C开始发送的瞬间，A剩$10-5=5$槽，B剩$11-5=6$槽。C的单独发送可被正确理解，A、B冻结计数器，并利用Duration信息使虚拟侦听保护其后续交换。

\begin{figure}[htbp]
\centering
\begin{tikzpicture}[x=0.95cm,y=0.8cm,>=stealth,font=\small]
\foreach \y/\n in {3/AP,2/A,1/B,0/C}{\draw[->] (0,\y)--(12.6,\y);\node[left] at (0,\y){\n};}
\foreach \y in {1,2}{\draw[fill=red!15] (0.2,\y-0.2) rectangle (2.4,\y+0.45);\node at (1.3,\y+0.13){碰撞发送};}
\node at (1.3,3.25){无ACK};
\node at (3.4,2.28){退避$10\to5$};\node at (3.4,1.28){退避$11\to6$};
\draw[fill=yellow!20] (2.4,-0.2) rectangle (3.5,0.45);\node at (2.95,0.13){EIFS};
\draw[fill=blue!15] (4.5,-0.2) rectangle (7.5,0.45);\node at (6,0.13){C数据};
\draw[fill=green!20] (7.8,2.8) rectangle (8.5,3.45);\node at (8.15,3.13){ACK};
\foreach \y in {1,2}{\draw[fill=gray!15] (4.5,\y-0.2) rectangle (8.5,\y+0.45);\node at (6.5,\y+0.13){冻结；保护C交换};}
\node at (9.3,2.3){$D+5\sigma$};\node at (9.3,1.3){$6\to1$};
\draw[fill=blue!15] (10.1,1.8) rectangle (12.1,2.45);\node at (11.1,2.13){A数据};
\foreach \x/\v in {0.2/0,2.4/1\mathrm{ms},4.5/c,8.5/d,10.1/3.1\mathrm{ms}}{\draw[dashed] (\x,-0.5)--(\x,3.6);\node[below] at (\x,-0.5){$\v$};}
\end{tikzpicture}
\caption{例一原创事件示意，横轴不按比例。数值采用已核对的原题截图：$F=150$微秒。横轴压缩长发送区段，以保留短间隔与计数变化的可读性。}
\end{figure}

\subsubsection{恢复倒计数并完成发送顺序}

在$d$之后，A和B先等待DIFS，再继续原来的剩余计数，而不是重新抽取10和11。经过5槽，A到零，B剩1槽。A发送的开始时刻因此为
\begin{equation}
a=d+D+5\sigma=3.1\,\mathrm{ms}.
\end{equation}
把$c$、$d$的表达式代入，可得一致性约束
\begin{equation}
2D+Q+5\sigma=0.35\,\mathrm{ms}.
\end{equation}
代入原题数值，$2(0.05)+0.15+5(0.02)=0.35$ ms，两条独立时间链一致。具体地，C在1.30 ms开始发送，到2.80 ms结束；AP等待SIFS后于2.81--2.95 ms发送ACK。A、B在2.95--3.00 ms等待DIFS，再倒计数5槽，A恰在3.10 ms到零。

A的数据持续1 ms，其ACK结束于$a+1\,\mathrm{ms}+Q$。B在A交换期间保持剩余1槽不动，之后等待DIFS并完成这1槽，故B的发送开始时刻为
\begin{equation}
b=a+1\,\mathrm{ms}+Q+D+\sigma.
\end{equation}
代入数值，A数据占3.10--4.10 ms，ACK占4.11--4.25 ms。B在4.25--4.30 ms等待DIFS，再用0.02 ms完成最后1槽，所以从4.32 ms发送到5.32 ms；最后ACK占5.33--5.47 ms。整次过程在5.47 ms结束。成功发送顺序为C、A、B。决定这一顺序的不是“A最早有数据”，而是每次真正可以倒计数的空闲时段以及被冻结的剩余槽数。

\begin{figure}[htbp]
\centering
\begin{tikzpicture}[x=2.05cm,y=0.78cm,>=stealth,font=\small]
\foreach \y/\n in {3/AP,2/A,1/B,0/C}{\draw[->] (0,\y)--(5.65,\y);\node[left] at (0,\y){\n};}
\foreach \y in {1,2}{\draw[fill=red!15] (0,\y-0.18) rectangle (1,\y+0.38);\node at (0.5,\y+0.1){碰撞};}
\draw[fill=blue!15] (1.30,-0.18) rectangle (2.80,0.38);\node at (2.05,0.1){C数据};
\draw[fill=blue!15] (3.10,1.82) rectangle (4.10,2.38);\node at (3.60,2.1){A数据};
\draw[fill=blue!15] (4.32,0.82) rectangle (5.32,1.38);\node at (4.82,1.1){B数据};
\foreach \l/\r in {2.81/2.95,4.11/4.25,5.33/5.47}{\draw[fill=green!30] (\l,2.82) rectangle (\r,3.38);\node[above] at ({(\l+\r)/2},3.42){ACK};}
\foreach \x/\v in {0/0,1/1.00,1.3/1.30,2.8/2.80,3.1/3.10,4.1/4.10,4.32/4.32,5.47/5.47}{\draw[dashed,gray] (\x,-0.45)--(\x,3.4);\node[below,rotate=45,anchor=north east] at (\x,-0.45){\v};}
\node[below] at (2.8,-1.2){绝对时刻，单位ms；细节间隔见前图与正文};
\end{tikzpicture}
\caption{例一完整成功发送顺序与终点。这里按真实时间比例绘制，短ACK和等待间隔因此明显小于数据发送区段。}
\end{figure}

\subsection{例二：较短发送者会听到较长帧的尾部}

\subsubsection{不同结束时刻形成不同局部知识}

A与B仍从0同时开始发送，但A的数据为1 ms，B的数据为1.4 ms；C以后发送的数据为1 ms。C在最初信道忙期间就已经得到待发分组。确认超时为$F=0.15$ ms。此题的时间参数是一组独立条件，尤其EIFS已在课堂说明与前题不同；不能沿用前题的$E=0.2$ ms。由于本题完整参数页尚未核实，下文的$E,D,Q,\sigma$保留为本题自己的符号，给出事件顺序和可检验的时间约束，不把它们偷换成例一数值。课堂后续推导使用A退避19槽、C退避5槽，B按30槽解释时与给出的剩余槽数和成功顺序一致。

在1 ms时A停止发送并转为接收，此时B还要发送0.4 ms。A没听到ACK，却听到了B的帧尾。它错过了B的前导和头部，不能从残余比特可靠地恢复完整帧和Duration，因此在本题错误接收模型下，A也需要从实际忙结束时刻1.4 ms起等待EIFS。其自身ACK超时1.15 ms已经到期，但信道仍忙，不能从1.15 ms开始倒计数。

C从一开始旁听到碰撞，直到1.4 ms才见信道空闲，也从这里起等待EIFS。B则在整个1.4 ms内都处于发送状态，没有在发送期间执行接收；其发送结束后等待$F$，未收到ACK，再按普通DIFS进入重试。因此A与C的倒计数起点相同，而B的起点由另一条时间链决定：
\begin{equation}
x_A=x_C=1.4\,\mathrm{ms}+E,
\qquad x_B=1.4\,\mathrm{ms}+F+D.
\end{equation}
这里没有把$F$与$E$在A上首尾相加，因为A的失败确认等待与继续侦听残余忙信号发生在同一真实时间轴上；必须先判断哪个约束在何时结束。

\subsubsection{用已消耗槽数校验起点偏移}

课上解答指出，C第一次发送之前，A递减5槽，B递减11槽。因此B比A、C提前倒计数6槽，有
\begin{equation}
x_A-x_B=6\sigma,
\qquad E-F-D=6\sigma.
\end{equation}
C的开始时刻为$c_2=1.4\,\mathrm{ms}+E+5\sigma$。到此时，各站剩余计数为
\begin{equation}
B_A=19-5=14,
\qquad B_B=30-11=19,
\qquad B_C=5-5=0.
\end{equation}
C发送时，A、B都能正确接收其头部，建立NAV并冻结。C交换结束于$d_2=c_2+1\,\mathrm{ms}+Q$。

\begin{figure}[htbp]
\centering
\begin{tikzpicture}[x=0.9cm,y=0.85cm,>=stealth,font=\small]
\foreach \y/\n in {2/A,1/B,0/C}{\draw[->] (0,\y)--(13.6,\y);\node[left] at (0,\y){\n};}
\draw[fill=red!15] (0.2,1.8) rectangle (2.4,2.45);\node at (1.3,2.13){发送};
\draw[fill=orange!20] (2.4,1.8) rectangle (3.5,2.45);\node[above,align=center] at (2.95,2.5){听到B尾部};
\draw[fill=red!15] (0.2,0.8) rectangle (3.5,1.45);\node at (1.85,1.13){较长发送};
\draw[fill=gray!15] (0.2,-0.2) rectangle (3.5,0.45);\node at (1.85,0.13){碰撞忙；C已有包};
\foreach \y in {0,2}{\draw[fill=yellow!20] (3.5,\y-0.2) rectangle (6.6,\y+0.45);\node at (5.05,\y+0.13){EIFS};}
\node at (4.45,1.25){$F+D$};\node at (6.2,1.25){先减6槽};
\node at (7.7,2.25){$19\to14$};\node at (7.7,1.25){$24\to19$};\node at (7.7,0.25){$5\to0$};
\draw[fill=blue!15] (8.7,-0.2) rectangle (11.8,0.45);\node at (10.25,0.13){C数据};
\foreach \y in {1,2}{\draw[fill=gray!15] (8.7,\y-0.2) rectangle (12.8,\y+0.45);\node at (10.75,\y+0.13){冻结至C的ACK结束};}
\foreach \x/\v in {0.2/0,2.4/1,3.5/1.4,6.6/x_A,8.7/c_2,12.8/d_2}{\draw[dashed] (\x,-0.55)--(\x,2.5);\node[below] at (\x,-0.55){$\v$};}
\node[below] at (6.8,-1.1){标注1与1.4的单位为ms；其余为事件时刻};
\end{tikzpicture}
\caption{例二原创时序，突出B先开始倒计数而C先完成。转写出现B取13或30的冲突；此处采用与“先减11槽、之后A先发”相容的30，未把它视为已核验原题。}
\end{figure}

C成功之后，A、B同时等DIFS并恢复各自剩余计数。A只需14槽，B还需19槽，所以A先发，B剩5槽。于是
\begin{align}
a_2&=d_2+D+14\sigma,\\
b_2&=a_2+1\,\mathrm{ms}+Q+D+5\sigma.
\end{align}
B最后发送的数据长1.4 ms，因而最终交换结束于$b_2+1.4\,\mathrm{ms}+Q$。这个解清楚地区分了原来的初始退避数、C发送前已经消耗的槽数，以及C和A发送期间无法消耗的冻结时间。

若把B初始值改成13而保留“已减11槽”，C开始时B仅剩2槽，C结束后B会先于尚剩14槽的A发送。因此13与课堂随后给出的A先发顺序不相容。这个检验不依赖缺失的时隙具体数值，是利用协议状态和相对顺序排查转写错误的有效方法。

\subsection{从两个例子提炼通用解题过程}

首先把每个实际发送区段画在同一时间轴上，尤其不能把不同长度的碰撞帧画成同时结束。其次分别记录每个站点的状态：正在发送、正在等待ACK、听到但不能解码、成功读到Duration，或信道确实空闲。接着确定各站首次具备倒计数资格的时刻；错误接收导致的EIFS应从实际信道空闲处计算，发送失败的超时则要明确从哪里开始计量。

然后只在满足规则的完整空闲时隙上递减计数器。某站开始发送时，先记下其他站的剩余值，再冻结到受保护交换结束。恢复时保留剩余值，重新满足必要的空闲间隔，再继续倒计数。最后用“初始值等于各次实际递减之和加剩余值”检查每个站点，以防把冻结时间误算为退避时间。

两题最重要的区别是C的到达时刻和A听到的内容。第一题C在错误保护已经结束后才得到新分组，课堂模型让它按初次接入流程抢在重试者之前；第二题C早已在忙期间等待，因此需要在EIFS后完成退避。第一题A、B同时结束，二者后续对称；第二题A先结束并听到B的帧尾，必须采用与B不同的错误恢复等待。只有把这些局部观察逐项落实，全局发送顺序才会自然得到。

% Part 4: course-ordered continuation, 2026-10-14 to 2026-10-16.
\section{从隐藏节点时序到设备间直接通信}
\label{ch:d2d}

\subsection{隐藏节点条件下的退避推演}

分布式协调功能（Distributed Coordination Function, DCF）的计算题不能只画一条代表全网的“信道忙闲线”。每个站点掌握的信道状态取决于它能够听到哪些发送端。同一时刻，接入点可能正在接收，某个隐藏站点却仍把信道判断为空闲。因此，时序分析必须先确定拓扑，再分别记录各节点的物理侦听、网络分配向量（Network Allocation Vector, NAV）、等待间隔和退避计数。

\subsubsection{不同包长碰撞后的不对称状态}

设两个站点同时发送，但A的数据帧较短，B的数据帧较长。A发送期间不能同时接收；A结束发送后，可能只听到B帧的尾部。虽然检测到了无线能量，A并没有从头接收到可正确解码的帧，不能根据这个尾部恢复可靠的Duration字段。因而，在错误接收所对应的情形下，A待信道重新空闲后使用扩展帧间间隔（Extended Interframe Space, EIFS）。B则需要根据确认超时判断本次发送失败，再进入重传竞争。

这里有两条不同的因果链：A的额外等待来自接收错误或无法正确解码的活动，B的重传来自未收到确认帧（Acknowledgment, ACK）。不能因为两者参加了同一次碰撞，就让两者在完全相同的时刻进入同一种等待状态。若C在碰撞期间收到上层待发分组，还须根据C自身是否成功读取首部，决定它使用NAV还是错误接收后的保护等待。待C先取得发送机会后，能够正确接收C首部的其他站点设置相应NAV，并冻结尚未完成的退避。

\subsubsection{两包连续发送与后退避}

考虑图\ref{fig:p4-hidden}的拓扑。A与B相互隐藏，C能够听到A、B和接入点，接入点也能与所有站点通信。B先收到两个待发分组，第一帧发送成功后，不能直接把第二帧接在后面独占信道，而要执行后退避（Post-backoff）。如果只研究每站一帧的例子，成功发送后的后退避可能不影响所关心的结果；当同一站点还有第二帧时，它就成为决定发送顺序的关键状态。

\begin{figure}[htbp]
\centering
\begin{tikzpicture}[scale=0.95,every node/.style={font=\small}]
\node[draw,circle,minimum size=9mm] (a) at (0,0) {A};
\node[draw,circle,minimum size=9mm] (c) at (3,0) {C};
\node[draw,circle,minimum size=9mm] (b) at (6,0) {B};
\node[draw,rounded corners,minimum width=16mm,minimum height=9mm] (ap) at (3,2) {AP};
\draw (a)--(c)--(b); \draw (a)--(ap)--(b); \draw (c)--(ap);
\draw[dashed] (0,-1)--(6,-1);
\node[fill=white] at (3,-1) {A与B之间无可侦听链路};
\end{tikzpicture}
\caption{隐藏节点例子的原创拓扑示意。实线表示题设中的相互可听关系，不表示节点必定同时发送。}
\label{fig:p4-hidden}
\end{figure}

推演时应先标出上层分组到达、已经给定的发送开始时刻和帧长，再推导其余动作。尤其要区分“收到待发分组”和“已经开始发送”：前者还可能需要等待DCF规定的空闲间隔，后者则是时间轴上的既定事件。接下来分别画出ACK、NAV与退避，最后才比较谁的计数器先归零。

第三个例子以下列已经确定的中间计数状态为起点，展开剩余退避约束，并不重新给定完整的绝对到达时间题干：B在A的数据发送及ACK附近的危险边界前，共经历13个可递减时隙。与此同时，C原来抽取7，在A接入前已递减2，因此在A与AP完成交换后还剩5。令B初始后退避整数为$b$，则在这个重新竞争的阶段，其剩余值为
\begin{equation}
 r_B=b-13,\qquad r_C=7-2=5.
\end{equation}
B要在C之前发送，要求$r_B<5$；B又不能在AP准备发送ACK或仍不能接收的边界上归零，要求$r_B\geq1$。两条约束合并为
\begin{equation}
 1\leq b-13\leq4
 \quad\Longrightarrow\quad b\in\{14,15,16,17\}.
\end{equation}
这不是把ACK持续时间也折算成B必须额外抽取的时隙。B能够听到AP，一旦ACK开始，B的物理侦听就会使剩余退避冻结；ACK结束并经过分布式帧间间隔（DCF Interframe Space, DIFS）后，才继续消耗剩余时隙。

取$b=13$的问题在于计数器过早归零：A的数据发送虽然可能刚结束，AP尚处于准备回复ACK或发送ACK的阶段。半双工（Half-duplex）接入点不能一边发ACK、一边接收B的新数据。取$b=18$则留下5个时隙，使B与C同时归零、同时发送。更大的值可能让C先发送而B随后成功，但不能满足“B的两帧都先于C完成”的限定。

\subsubsection{只限制碰撞次数时，可行区间如何改变}

第四个例子保留同一隐藏节点拓扑，但先给定B开始发送，A稍后开始发送，并使两帧在同一时刻结束。C在重叠发生之前已经读取B的首部，因此即使后续数据碰撞，C仍可按先前读出的Duration维持预期交换的NAV。A、B均因没有收到ACK而进入失败后的竞争。

继续使用例题给出的中间计数状态：A的给定退避整数为1，C的为5；B看不到A的发送，故仍会递减，并在AP的ACK使其冻结前已经消耗16个完整可递减时隙。这里把16作为已给出的中间结果，随后完整求解可行退避集合。于是
\begin{equation}
 r_B=b-16,\qquad r_C=5.
\end{equation}
若B先发送，必须有$1\leq r_B\leq4$，得到$b=17,18,19,20$。若B、C同时归零，则$r_B=5$，对应$b=21$，应当排除。若允许C先发送，$r_B\geq6$也可行；C发送时，B听到C并冻结，C完成交换之后，B继续剩余退避。

因此，在例题限定的抽样范围$0\leq b\leq31$内，可行集合为
\begin{equation}
 \mathcal{B}=\{17,18,19,20\}\cup\{22,23,\ldots,31\}.
\end{equation}
$b=16$不能取，原因仍是AP的收发切换与ACK保护，而非A的数据帧本身尚未结束。两个例子的共同方法是“先保护不可侦听的发送，再比较可共同侦听阶段的剩余计数”；不同答案来自发送顺序约束不同。这里的31是例题采用的候选上限，不应据此取消一般DCF在失败后调整竞争窗口的规则。

\subsection{设备到设备网络的动机与工作方式}

设备到设备通信（Device-to-Device Communication, D2D）允许邻近设备通过直接无线链路交换数据。自组织网络（Ad hoc Network）进一步强调：节点可以按需要组成网络，不必先部署一套固定基础设施。灾害救援中既有网络可能损坏；临时活动中可能需要快速建立局部通信；个人手机、电脑、耳机和物联网标签也常需要近距离互联。直接通信减少了对固定设施的依赖，部署较快，且能够利用短距离链路降低发射功率。

纯自组织网络（Pure Ad hoc Network）的数据完全在参与设备之间传递。混合网络（Hybrid Network）则同时保留直接通信与基础设施连接，某个节点或网关（Gateway）负责把局部网络接入互联网。网关的存在不意味着所有局部数据都必须绕行互联网；同一网络可以让近邻直接交换，同时让跨网络业务经过网关。

在蜂窝网络辅助的D2D场景中，基站仍可能管理资源或提供外部连接，但局部内容分发可由设备协作完成。例如，许多邻近用户需要同一内容时，基站向少量设备传送，随后由这些设备向附近用户转发，就可能降低基站重复发送的负担。这种收益以设备距离近、链路可靠、资源协调适当以及设备愿意转发为前提。

\subsubsection{空间复用的含义}

空间复用（Spatial Reuse）是指在彼此干扰足够小的不同空间区域重复使用同一频率资源。设一个服务区域可使用$K$个信道，每个信道在给定业务负载下支持约$u$个用户，则单个覆盖区域的简单容量估计为$Ku$。把服务区域划分为多个低功率区域，并在满足干扰约束的位置重复使用同一组信道，有机会增加同时服务的总用户数。实际容量不会无条件按区域数线性增长，因为边缘干扰、调度和复用间距都会限制增益。

D2D把这一思想带到设备附近。较短的传输距离通常允许较低的发射功率，较低功率使强干扰的范围缩小，于是远处的另一对设备可能同时使用同一信道。若设备并不邻近，必须提高功率才能完成直接通信，空间复用收益便可能减弱。

\begin{figure}[htbp]
\centering
\begin{tikzpicture}[every node/.style={font=\small}]
\draw[dashed] (0,0) circle (1.25); \draw[dashed] (4.6,0) circle (1.25);
\node[draw,circle] (u1) at (-0.45,0) {$S_1$};
\node[draw,circle] (v1) at (0.45,0) {$D_1$};
\node[draw,circle] (u2) at (4.15,0) {$S_2$};
\node[draw,circle] (v2) at (5.05,0) {$D_2$};
\draw[->,thick] (u1)--(v1); \draw[->,thick] (u2)--(v2);
\node at (0,-1.7) {近距离、低功率链路};
\node at (4.6,-1.7) {可复用同一信道};
\node at (2.3,1.7) {是否能并发取决于接收端的实际干扰};
\end{tikzpicture}
\caption{空间复用的原创概念图。虚线仅表示局部干扰范围的简化示意，不是固定的无线传播边界。}
\label{fig:p4-reuse}
\end{figure}

\subsection{业务、设备与移动性的差异}

D2D中的节点并不一定同质。手机和笔记本通常比微型传感器具有更多计算资源、更大的存储空间和不同的电池容量。设备的逻辑责任也可能不同，例如某种技术安排一个设备协调通信，另一些设备只完成简单数据交换。能力较强并不自动意味着承担所有管理工作，角色仍取决于协议和网络组织。

业务目的决定了交付方式。单播（Unicast）面向一个确定目的端；广播（Broadcast）面向广播范围内的所有节点；多播（Multicast）面向指定接收组。地理广播（Geocast）按地理区域选择接收者，例如让某一道路范围内的车辆接收交通告警；任播（Anycast）则只需把数据交给某个合适的组成员，例如若干可提供同一服务的节点中的一个。任播不要求所有成员都收到。

业务要求至少应同时描述比特率、时延和可靠性。以成功交付比例描述可靠性时，可以写成
\begin{equation}
 \rho=\frac{N_{\mathrm{delivered}}}{N_{\mathrm{generated}}}.
\end{equation}
分母是观察期间产生的目标分组数，分子是按所选成功判据交付的分组数；如果应用另有截止时间，还应明确是否只把及时到达的分组算作成功。低速率温度测量与实时告警的要求不同，不能只凭平均流量判断设计优先级。

移动性（Mobility）既包括移动速度，也包括移动模式。游牧式使用往往是位置改变后停留；车辆沿道路高速移动；家庭设备、机器人或随机移动的人群又具有不同规律。速度越高，邻居关系的有效时间可能越短，网络发现和路由维护的工作也可能越频繁。固定节点同样可能因障碍物、故障或休眠而失去可用链路。

\subsection{拓扑、接入和多跳转发}

\subsubsection{用随时间变化的图表示网络}

网络拓扑（Network Topology）可抽象成图
\begin{equation}
 G(t)=\bigl(V(t),E(t)\bigr),
\end{equation}
其中$V(t)$是当前可参与通信的节点集合，$E(t)$是当前可用链路集合。邻居（Neighbor）指能够直接通信的节点。简单模型常把通信范围内的链路视为可靠，把范围外的链路视为不存在；真实系统中，衰落、干扰和接收门限会使边的可用性更加复杂。

通信范围也不能与干扰范围混为一谈。一个信号可能已经弱到无法正确解调，却仍足以破坏另一个接收机正在接收的数据。因而，“不是邻居”不等于“互不干扰”。半双工限制进一步要求节点分时完成发送和接收。

平面拓扑（Flat Topology）中没有强制的层级管理结构，网状拓扑（Mesh Topology）可提供多条替代路径。对$N$个节点的简单无向图，每对不同节点至多对应一条边，因此完全连接时的最大边数为
\begin{equation}
 |E|_{\max}=\frac{N(N-1)}{2}.
\end{equation}
分子$N(N-1)$把每一对节点计算了两次，除以2得到无向边数。一般的平面网并不必然达到这个上限。

树形拓扑（Tree Topology）在连通且无环时恰有$N-1$条边。两节点之间只有一条简单路径，所以路径中的一条边失效即可使某些节点之间失去连通。网状结构若存在冗余路径，单边失效后可能绕行；但冗余连接也意味着更多潜在干扰和状态维护工作。

分簇拓扑（Cluster-based Topology）把节点组织为若干簇，由簇头（Cluster Head）承担簇内协调或数据转发。桥接节点（Bridge Node）连接相邻簇，使数据能够跨簇逐跳传播。分簇可以减小局部管理范围，却也会使簇头和桥接节点成为重要的负载集中点。

\subsubsection{接入负责一跳，路由负责路径}

信道接入（Channel Access）决定某一跳的多个竞争者何时使用共享无线信道；路由（Routing）决定数据经过哪些中继（Relay）到达目的端。即使每一跳都遵守MAC规则，网络仍需要选出可到达目的端的路径。反过来，即使路径已经选好，每一跳也仍需获得接入机会，路径存在并不意味着能立即发送。

网络层（Network Layer）执行路由功能。源节点选择路径或下一跳，中继收到数据后继续选择下一跳；MAC层完成本跳发送规则，物理层（Physical Layer）执行实际无线信号传输。应用层产生的信息向下交付，成功接收后再由目的端协议栈向上交付。不能把中继仅看成“把听到的信号原样放大”：这里讨论的是节点接收分组后进行下一跳转发。

\begin{figure}[htbp]
\centering
\begin{tikzpicture}[every node/.style={font=\small}]
\node[draw,circle] (a) at (0,0) {A}; \node[draw,circle] (c) at (6,0) {C};
\node[draw,circle] (b) at (3,1.6) {B};
\node[draw,circle] (d) at (2,-1.6) {D}; \node[draw,circle] (e) at (4,-1.6) {E};
\draw[->] (a)--node[above]{一跳直达}(c);
\draw[->,thick] (a)--(b); \draw[->,thick] (b)--node[above right]{两跳}(c);
\draw[->,dashed] (a)--(d); \draw[->,dashed] (d)--(e); \draw[->,dashed] (e)--(c);
\node at (3,-2.3) {三跳候选路径：A--D--E--C};
\end{tikzpicture}
\caption{同一源、目的端的不同候选路径。原创图用于说明路径选择，并不规定某条路径总是最优。}
\label{fig:p4-routes}
\end{figure}

多跳通信（Multi-hop Communication）可把一个较远的传输分解为数个近距离传输。源节点因只需到达邻近中继，可能降低自身发射功率；较小的干扰范围又可能改善空间复用。但每增加一跳，就增加一次信道使用、排队与接入过程，也增加中继的接收和发送工作。因此，源节点节能不等于全网总能耗必然降低，跳数较少也不保证路径在所有指标上最优。

链路失效时，若还存在候选路径，路由可转向备用路径。这样的稳健性以节点保持参与和愿意中继为前提。民用设备的拥有者可能不愿为他人承担电池和资源开销，需要激励机制；任务型网络可能由共同任务规定合作责任。后续分析采用中继愿意合作转发的模型。

\subsection{能量约束与分布式设计}

移动设备通常依赖电池；固定在偏远、恶劣环境中的传感器也可能无法接入电源。电池容量受到体积、重量和维护条件限制。若用能量预算$E_{\mathrm{bat}}$和长期平均功率$\bar P$表示，理想化寿命估计为
\begin{equation}
 T_{\mathrm{life}}\approx\frac{E_{\mathrm{bat}}}{\bar P}.
\end{equation}
这一关系说明，降低长期平均功率与增加电池能量都能延长寿命；但在体积已受限制时，协议设计往往只能从减少无效活动入手。

\subsubsection{功率控制与连通性的平衡}

功率控制（Power Control）通过调整发射功率，在可达性、干扰和能耗之间取舍。采用简化路径损耗模型时，接收功率可写为
\begin{equation}
 P_r=K\frac{P_t}{d^{\alpha}},
\end{equation}
$P_t$是发射功率，$d$是收发距离，$\alpha$是路径损耗指数，$K$汇总模型中的其他固定因素。若需要维持相同的接收门限$P_{\min}$，则
\begin{equation}
 P_t\geq\frac{P_{\min}}{K}d^{\alpha}.
\end{equation}
因此，把距离缩短为原来的某个比例，会以该比例的$\alpha$次方改变所需发射功率。这解释了近邻转发的节能潜力，也说明降低功率不能脱离接收门限随意进行。

高功率通常增加可用邻居和冗余边，但会扩大干扰；低功率减小干扰，却可能使网络变得稀疏甚至分裂。合理目标是以足够低的功率保持任务需要的连通性和可靠性，而非把所有节点始终设为最大功率或最小功率。

\subsubsection{活动调度与休眠深度}

活动调度（Activity Scheduling）决定何时保持设备清醒、何时进入休眠（Sleep）。无线接口即使没有发送，也可能因持续侦听而消耗显著能量。关闭不必要的收发和计算组件可以节能，但睡着的节点不能立即接收数据，还可能暂时从可用拓扑中消失。

浅休眠保留较多状态，唤醒较快，但功耗相对高；深休眠关闭更多组件，功耗更低，但恢复工作通常需要更长时间。选择休眠深度时，应同时考虑下次预期活动时刻、唤醒代价、允许的接收时延和网络对该节点的转发依赖。频繁进入、退出深休眠未必比保持短暂浅休眠更省能。

拓扑、信道接入、多跳路由和能量消耗构成D2D设计中互相影响的四个问题。自治自组织网络通常需要分布式算法（Distributed Algorithm），让节点根据本地观测与有限交换作决定；集中式算法（Centralized Algorithm）则把信息和决策集中于管理者。分布式并不等于没有任何协调，而是不能把所有正确运行条件都建立在一个永远可达的全局控制者之上。自组织Wi-Fi、蜂窝辅助D2D、车联网、蓝牙和体域网（Body Area Network, BAN）体现了不同的角色、移动性和能量约束；传感器网络则把低功耗要求推到了设计中心。

\section{无线传感器网络与低功耗介质访问}
\label{ch:wsn}

\subsection{从物理测量到网络数据}

无线传感器网络（Wireless Sensor Network, WSN）由分布于环境中的传感节点组成，是物联网（Internet of Things, IoT）的重要组成部分。节点测量温度、压力、气体浓度、位置或其他物理量，并把有用信息交给汇聚节点（Sink）。汇聚节点连接局部传感网络与外部系统，承担数据出口，有时也负责向节点下发控制信息。

大范围森林火灾监测说明了其典型工作链：各位置的传感器定期测温，本地处理识别异常，必要时报告告警；远端系统收到告警后再组织响应。偏远或寒冷地区的环境监测同样依赖长期电池供电，维护困难使通信效率和寿命尤为重要。

\subsubsection{传感节点的基本组成}

传感板（Sensing Board）把物理量转换成可处理的测量值，信号处理单元（Signal Processing Unit）进行分析、筛选或事件检测，无线接口（Wireless Interface）负责数据交换。节点还可以控制执行器（Actuator），把计算结果转化成物理动作。例如，超过门限的温度测量既可以触发远端告警，也可以在本地触发某种控制动作。

\begin{figure}[htbp]
\centering
\begin{tikzpicture}[every node/.style={font=\small},box/.style={draw,rounded corners,text width=23mm,align=center,minimum height=11mm}]
\node[box] (sense) at (0,0) {传感板\\物理量采样};
\node[box] (proc) at (3.2,0) {信号处理\\事件检测};
\node[box] (radio) at (6.4,0) {无线接口\\测量与告警};
\node[box] (act) at (3.2,-1.8) {执行器\\本地控制};
\draw[->] (sense)--(proc); \draw[<->] (proc)--(radio); \draw[->] (proc)--(act);
\node at (0,-1.8) {电源支持各模块};
\end{tikzpicture}
\caption{传感节点的信息处理链原创示意。模块可以集成于同一设备，执行器并非每个节点都必须具备。}
\label{fig:p4-sensor}
\end{figure}

物理现象通常随时间和空间连续变化，而网络只能处理有限测量。设传感器位于$\boldsymbol{x}_i$，在时刻$t_k$采样，则样本可记为
\begin{equation}
 y_i[k]=s(\boldsymbol{x}_i,t_k),
\end{equation}
其中$s$是所观察的物理场。时间采样间隔决定变化捕捉的细致程度，空间部署决定不同位置的覆盖程度。本地处理与外部系统可利用这些离散数据重建或估计物理变化；仅有少量采样点并不自动保证重建精确，采样密度应适应现象变化的尺度。

\subsubsection{传输介质与低速率特征}

传感器的无线接口可以采用射频（Radio Frequency, RF）、红外或光通信。红外、光通信对视距（Line of Sight, LOS）和方向通常更敏感，障碍物可能直接遮断链路。射频常更适合存在遮挡的环境，但仍受传播条件影响。微型“智能尘埃”（Smart Dust）体现了把感知、处理和通信集成到微小设备中的思路，不能据此假定所有传感器都使用同一种传输介质。

本课程采用的典型低速率节点尺度包括约5--150米的射频覆盖示例、约1--500 kbit/s的数据速率量级，以及约1 mW的低发射功率示例。这些数值用来说明低功耗短报文设计，而非所有产品或所有IEEE 802.15.4物理层的统一性能。以分贝毫瓦（dBm）表示功率时，
\begin{equation}
 P_{\mathrm{dBm}}=10\log_{10}\left(\frac{P}{1\,\mathrm{mW}}\right),
\end{equation}
所以1 mW对应0 dBm。小型电池、低成本和小尺寸共同限制节点资源，协议也因此不能照搬高吞吐量终端的设计优先级。

\subsection{汇聚业务、寻址与协议栈}

传感节点可以上传原始测量，也可以先在本地处理，只上传压缩后的摘要或检测到的事件。数据常由许多节点流向少数汇聚节点，构成以汇聚为中心的上行（Uplink）业务；配置、查询、时间信息或控制指令则沿下行（Downlink）方向传播。上下行业务量通常不对称，但下行流量较少并不意味着可以忽略，因为它可能决定采样、同步与控制能否正确进行。

当大量节点传输短载荷时，寻址开销十分明显。若应用数据只有$L$字节，地址及其他协议开销共$H$字节，则报文有效比例为
\begin{equation}
 \eta=\frac{L}{L+H}.
\end{equation}
当$L$很小时，固定首部$H$占比就很大。短地址能减轻局部传输负担，但仍须保证所使用范围内的唯一性，并处理与外部寻址体系的衔接。不能只从地址更短就推出整个系统不需要地址管理。

时间同步（Time Synchronization）让不同节点的测量可在时间上比较，也支持协同休眠和时隙调度。自定位（Self-localization）则给测量加上空间含义：温度值相同，但出现位置不同，可能意味着完全不同的事件。感知数据是否有用，不仅取决于它是否到达，还取决于“何时、何地测得”是否可解释。

\subsubsection{应用如何形成设计要求}

环境监测、滑坡预警、动物跟踪、建筑管理、车辆系统、水质和土壤监测，都利用了分散节点在现场收集信息的能力。智能电网（Smart Grid）需要观察发电、负载和设备状态，以支持调节与控制；可再生能源的变化使及时、分布式测量具有价值。土壤污染、水中微生物和农业条件的监测又强调长期部署、覆盖与维护成本。相同的感知网络框架可以服务不同应用，但采样周期、时延容忍度和可靠性目标应随应用改变。

商用传感平台体现了模块化实现：同一平台可配置不同传感板，选择不同无线接口，并提供多种休眠模式。气体、温度、液体、重量和压力等测量能力来自不同传感器模块。接口种类、发射功率和覆盖距离应逐项理解，不应把一个平台不同配置的最佳指标同时当成单一节点的固定能力。

\subsubsection{协议层之间的配合}

物理层负责信号传输；MAC层负责接入共享介质和组织收发活动；网络层负责路由，并适应因移动、故障或休眠而变化的拓扑。较高层还可能提供缓存（Caching）、查询（Query）和数据聚合（Data Aggregation），把多个测量整合成应用需要的信息，再与外部数据库或业务系统连接。

休眠会使“无线链路在地理上可达”和“节点此刻可接收”成为两件事。路由若选择了正在深睡的中继，MAC即使找到空闲信道，也无法立即完成交付。同步、定位和能量管理会跨越多个层次发生影响，分层用于明确职责，而不能消除这些协同需求。

\subsection{MAC设计的主要目标与能量浪费}

无线传感器网络的MAC尤其重视能量效率（Energy Efficiency）、可扩展性（Scalability）和拓扑适应性（Adaptability）。节点增多后，协议不应因控制负担过大而失效；节点休眠或重新加入时，接入机制也应继续运行。信道利用率、时延、吞吐量（Throughput）和公平性（Fairness）同样重要，只是其优先级取决于应用。温度低频采样与视频传输的吞吐量需求不同；突发告警即使平均速率很低，也可能要求很短的交付时延。

分布式接入意味着节点根据自己的观察决定行动。多个节点可能同时发现信道空闲并同时发送，因此“独立作出合理判断”并不等于“全网必然无碰撞”。MAC需要通过竞争分散或预先调度来减少这类同时行动。

\subsubsection{功耗状态与活动时间}

无线接口可处于发送、接收、空闲侦听和休眠状态。低功耗设计中的重要事实是：空闲侦听的功率可能与接收功率处于相近量级，而远高于休眠功率。总能量可按各状态停留时间估计为
\begin{equation}
 E\approx P_{\mathrm{tx}}t_{\mathrm{tx}}+
 P_{\mathrm{rx}}t_{\mathrm{rx}}+
 P_{\mathrm{idle}}t_{\mathrm{idle}}+
 P_{\mathrm{sleep}}t_{\mathrm{sleep}}.
\end{equation}
式中每项都是“该状态的功率乘以停留时间”。如果绝大多数时间都没有有效业务，却保持空闲侦听，那么即使每次发送已经很短，长期电量仍可能被侦听耗尽。CPU、传感板以及模式切换的消耗也应在完整预算中计入，上式仅把无线状态的主线展开。

四类典型浪费分别是：空闲侦听（Idle Listening），即没有有用接收却持续开机；碰撞（Collision），即已花费收发能量但数据不能正确交付；协议开销（Overhead），即控制报文、首部和维护活动消耗的资源；旁听（Overhearing），即接收并处理本不发给自己的报文。节能MAC不能只减少碰撞，还应减少无效清醒和不必要接收。

\subsection{基于调度的接入：LEACH}

基于调度的MAC（Schedule-based MAC）预先安排谁在何时发送。

基于竞争的MAC（Contention-based MAC）则由节点在业务出现时竞争接入。调度可减少碰撞并让节点知道何时休眠，但需要维护日程；竞争更灵活，却可能产生等待、碰撞和重试。

低能耗自适应分簇层次协议（Low-Energy Adaptive Clustering Hierarchy, LEACH）展示了分簇与调度结合的思路。普通节点向本簇簇头报告数据，簇头承担更多收集和转发工作。簇内采用时分多址（Time Division Multiple Access, TDMA），把重复的时间帧划成时隙，每个节点在自己的时隙发送，其他时间则有机会休眠。

\begin{figure}[htbp]
\centering
\begin{tikzpicture}[every node/.style={font=\small},x=1.2cm,y=1cm]
\foreach \x/\s in {0/S_1,1/S_3,2/S_5,3/S_1,4/S_3,5/S_5}{
\draw (\x,0) rectangle (\x+1,0.8); \node at (\x+0.5,0.4) {$\s$};}
\draw[<->] (0,1.15)--node[above]{第一个时间帧}(3,1.15);
\draw[<->] (3,1.15)--node[above]{下一个时间帧}(6,1.15);
\node[align=center] at (3,-0.65) {簇头按表接收；普通节点主要在自己的时隙活动};
\end{tikzpicture}
\caption{LEACH簇内TDMA的原创时隙示意。时隙标签仅用来说明重复调度。}
\label{fig:p4-leach}
\end{figure}

在理想同步与正确调度下，同一簇内不同节点不会占用同一个发送时隙，因此可避免簇内冲突。不同簇可能同时工作，课程采用码分多址（Code Division Multiple Access, CDMA）说明簇间区分：不同码使接收端能够区分目标信号和其他同时存在的信号。其直观意义是选择所关注的信号，而不是让其他信号在物理上消失；可分离性仍取决于码、功率和干扰条件。

簇头需要接收多个节点的数据并向外转发，能耗通常高于普通节点。轮换簇头（Cluster-head Rotation）把较重的职责在节点之间轮流分配，有助于避免固定簇头过早耗尽电池。轮换并不消除协调成本，但能改善负担分布。

静态时隙安排面临可扩展性问题。若时间帧已有三个固定时隙，新节点加入却没有空余时隙，就必须重新调度、延长帧，或降低某些节点的发送频率。设每个时隙长$T_s$，一个簇安排$N$个发送节点且每帧各发送一次，忽略其他开销时，
\begin{equation}
 T_{\mathrm{frame}}=NT_s,\qquad
 f_{\mathrm{report}}=\frac{1}{NT_s}.
\end{equation}
保持时隙时长不变而增加节点，会增加每节点的报告间隔。若坚持固定帧长，就需要其他资源调整，不能凭空增加无碰撞发送机会。

簇数量增加也可能使可用码资源不足；重复使用相同或相互干扰较大的码，需要考虑空间分隔与簇间干扰。LEACH由此说明了调度类方案的基本权衡：低冲突、可休眠的日程，与节点变化、调度维护和资源扩展之间存在张力。

\subsection{IEEE 802.15.4的分层位置、设备与拓扑}

IEEE 802.15.4定义低速率无线个人区域网（Low-Rate Wireless Personal Area Network, LR-WPAN）的物理层和MAC层。Zigbee基于这些底层能力提供更高层的网络与应用组织。两者并不是同一个名称的两种写法：物理层和MAC层可以支撑不同上层协议，上层的路由能力也不应全部归到IEEE 802.15.4本身。

全功能设备（Full-Function Device, FFD）具有较完整的协议功能，可承担协调器等角色；精简功能设备（Reduced-Function Device, RFD）面向较简单的终端任务，降低实现复杂度，通常通过FFD交换数据。具体多跳路由由所采用的上层网络协议完成，不能把“具有全功能MAC能力”直接等同于“IEEE MAC已经规定完整路由算法”。

个人区域网协调器（PAN Coordinator）管理网络成员，记录节点的加入与离开，并可提供同步参考。星形拓扑（Star Topology）以协调器为中心，外围终端可采用较简单的设备。对等拓扑（Peer-to-Peer Topology）允许有能力的节点之间直接通信，扩展连接结构；簇树拓扑（Cluster-tree Topology）则把多个局部星形或簇连接起来，由全功能节点承担中间协调角色，使较大覆盖范围与简单叶节点并存。

这些拓扑影响能量负担与故障后果。星形中心易于集中管理，但外围节点需要能够到达中心；更复杂的连接结构允许逐跳扩展，却增加转发与维护工作。后续MAC分析以传感节点和汇聚或协调节点之间的交换为基本场景，并特别关注节点如何安排清醒时间。

\section{IEEE 802.15.4的信标与竞争接入过程}
\label{ch:154}

\subsection{物理层频带与短分组}

本课程讨论的经典低速率物理层包括868/915 MHz附近的频段与2.4 GHz频段。频带中信道的可用带宽影响能够承载的符号速率，调制和编码又决定有效比特率，因此带宽是速率约束的重要组成部分，但不能单凭频率更高就认定速率必然更高。课件把较窄信道与较低速率、较宽信道与较高速率并列，用于说明这一设计关系。

短报文适合少量测量与控制信息。经典帧结构中，物理层前导、定界与长度等字段合计6字节，随后承载最长127字节的物理层服务数据单元（PHY Service Data Unit, PSDU）。这127字节承载的是MAC帧，其中还包含MAC首部和校验等开销，所以应用有效数据小于这个上限。讨论“载荷”时必须说明所属层次，物理层载荷与应用载荷不能直接互换。

\subsection{信标使能模式与超帧}

信标使能模式（Beacon-enabled Mode）中，PAN协调器周期性发送信标（Beacon），为其他节点提供时间参考。由信标开始组织的超帧（Superframe）把活动区间划分成不同接入阶段。先出现竞争接入期（Contention Access Period, CAP），节点采用时隙化竞争；随后可安排无竞争期（Contention-Free Period, CFP），为特定通信提供分配的时隙。

\begin{figure}[htbp]
\centering
\begin{tikzpicture}[every node/.style={font=\small},x=1cm,y=1cm]
\draw (0,0) rectangle (0.7,1); \node[rotate=90] at (0.35,0.5) {信标};
\draw (0.7,0) rectangle (4.5,1); \node[align=center] at (2.6,0.5) {CAP\\时隙CSMA/CA};
\draw (4.5,0) rectangle (7,1); \node[align=center] at (5.75,0.5) {CFP\\分配时隙};
\draw[dashed] (7,0) rectangle (10,1); \node[align=center] at (8.5,0.5) {非活动区间\\可休眠};
\draw (10,0) rectangle (10.7,1); \node[rotate=90] at (10.35,0.5) {信标};
\draw[<->] (0,-0.5)--node[below]{信标间隔}(10,-0.5);
\end{tikzpicture}
\caption{超帧的原创结构示意。图中长度仅表达先后关系；CFP与非活动区间是否存在由配置决定。}
\label{fig:p4-superframe}
\end{figure}

同步意味着节点共享时隙边界，并不意味着CAP由协调器逐个批准发送。CAP内部仍由节点依据算法分布式竞争。CFP则通过分配时隙减少竞争，并让节点知道何时必须清醒。

超帧可留出全网非活动区间，此时包括协调器在内的设备都可进入低功耗状态。单个节点也可以在自己不参与通信的时段休眠。这会把一部分负担转移到缓存和时延：如果目的节点睡着，发送方或协调器必须保存待交付数据，直到合适的唤醒机会。缓存容量有限，等待过久可能导致丢包；因此，休眠节能、存储需求与交付时延需要共同配置。

\subsection{先理解外层：接入许可与停止等待}

时隙载波侦听多路访问与碰撞避免（Slotted Carrier Sense Multiple Access with Collision Avoidance, Slotted CSMA/CA）可以按内外两个逻辑模块理解。外层负责“已经有一帧待发时，发送及确认是否成功”；内层负责“这一次能否取得信道”。先把内层当作只有两种输出的模块，就能把整体流程看清。

外层从初始化发送失败计数开始，然后调用内层。若内层返回允许发送，节点发送数据，并按停止等待（Stop-and-Wait）方式等待ACK。收到正确ACK，本帧成功；若ACK超时，增加失败计数，在尚有尝试预算时重新调用内层，否则报告失败。若内层直接返回接入失败，则本次不能进入数据发送，外层也结束为失败。

\begin{figure}[htbp]
\centering
\begin{tikzpicture}[every node/.style={font=\small},box/.style={draw,rounded corners,align=center,text width=38mm,minimum height=9mm}]
\node[box] (init) at (0,0) {待发帧；初始化发送计数};
\node[box] (inner) at (0,-1.4) {调用内层接入过程};
\node[box] (send) at (0,-2.8) {允许发送：发送并等待ACK};
\node[box] (check) at (0,-4.2) {未获ACK：累计失败次数};
\node[box] (retry) at (0,-5.6) {仍有发送尝试预算？};
\node[box] (fail) at (5,-1.4) {接入失败：结束};
\node[box] (ok) at (5,-2.8) {收到ACK：成功};
\node[box] (drop) at (5,-5.6) {预算耗尽：结束};
\draw[->] (init)--(inner); \draw[->] (inner)--(send); \draw[->] (inner)--(fail);
\draw[->] (send)--(check); \draw[->] (send)--(ok); \draw[->] (check)--(retry);
\draw[->] (retry)--node[above]{否}(drop);
\draw[->] (retry.west)--(-3,-5.6)--(-3,-1.4)--node[above]{是}(inner.west);
\end{tikzpicture}
\caption{按接入与确认两个层次重新绘制的外层流程。接入失败和发送后未获ACK是不同事件。}
\label{fig:p4-outer}
\end{figure}

这一划分防止把所有失败都叫作碰撞。内层可能连续发现信道忙，完全没有发出数据，就用尽了接入尝试；外层也可能已发出数据，但没有获得确认。后者可能来自碰撞或其他接收失败。仅从“没有ACK”通常不能唯一判定物理原因。

\subsection{内层：随机退避、两次CCA与重试}

\subsubsection{三个变量的不同职责}

内层使用退避次数计数$NB$、竞争窗口$CW$和退避指数（Backoff Exponent）$BE$。在此算法中，$CW$表示还需要成功完成多少次空闲信道评估（Clear Channel Assessment, CCA），初始化为2；它不是Wi-Fi中随机抽取退避整数的窗口上限。$BE$决定随机退避区间，$NB$记录本轮接入过程中因信道忙而进行的重试。

每次进入内层，初始化为
\begin{equation}
 NB=0,\qquad CW=2,\qquad BE=BE_{\min}.
\end{equation}
这里$BE_{\min}$为最小退避指数。时隙化系统先定位下一个退避周期（Backoff Period）边界。此处“退避周期”是一个基本时间单位，记其长度为$T_{\mathrm{BP}}$；它与一次完整随机退避持续时间不是同一概念。

\subsubsection{先随机等待，再判断信道}

节点等概率抽取整数
\begin{equation}
 U\sim\mathrm{Unif}\{0,1,\ldots,2^{BE}-1\},
 \qquad T_{\mathrm{bo}}=U T_{\mathrm{BP}}.
\end{equation}
例如$BE=3$时，可抽到0至7，共8个等可能值；若抽到4，就等待4个基本退避周期。取0只表示随机等待为零，不代表绕过后面的CCA直接发送。

本算法在随机等待过程中不采用Wi-Fi式的“忙则冻结、闲则递减”规则。节点先经过所抽出的退避时间，再在规定边界执行CCA。若一次CCA空闲，则$CW$减1；$CW$尚未到0时，继续下一次规定的CCA。只有两次所需的空闲评估均通过，内层才返回发送许可。

若CCA发现信道忙，无论此前已经通过了多少次空闲评估，都重新设置
\begin{equation}
 CW\leftarrow2,\qquad NB\leftarrow NB+1,\qquad
 BE\leftarrow\min(BE+1,BE_{\max}).
\end{equation}
$CW$恢复为2，避免把一次较早的空闲评估与后来跨越忙时段的评估拼接成发送许可；$NB$增加，记录本轮又失败了一次；$BE$增加，但不超过上限，使下一次随机等待的候选范围更宽。若仍有接入重试预算，就重新抽取退避；预算耗尽，则返回接入失败。

\subsubsection{一个包含忙信道的变量演化例子}

设$BE_{\min}=3$，且退避指数上限至少为4，开始时$(NB,CW,BE)=(0,2,3)$。第一次抽到$U=2$，经过两个基本周期后进行CCA。如果该次评估发现信道忙，变量变为$(1,2,4)$。第二次抽样范围于是扩大为0至15；假设抽到$U=3$，节点再等待三个周期。

若随后第一次CCA空闲，$CW$由2变为1；下一次CCA也空闲，$CW$由1变为0，内层返回允许发送。若第二次CCA反而忙，则不能发送，必须把$CW$恢复到2，并再次更新$NB$与$BE$。这一例子只是变量规则的演示，不是对任何节点将抽到哪些随机数的预测。

\begin{figure}[htbp]
\centering
\begin{tikzpicture}[every node/.style={font=\small},x=0.95cm]
\draw[->] (0,0)--(11.3,0) node[right]{时间};
\foreach \x in {0,...,10} {\draw (\x,-0.12)--(\x,0.12);}
\draw (0,0.5) rectangle (2,1.2); \node at (1,0.85) {随机等待2};
\draw (2,0.5) rectangle (3,1.2); \node at (2.5,1.55) {CCA忙};
\draw (3,0.5) rectangle (6,1.2); \node at (4.5,0.85) {新随机等待3};
\draw (6,0.5) rectangle (7,1.2); \node at (6.5,1.55) {CCA闲};
\draw (7,0.5) rectangle (8,1.2); \node at (7.5,2.05) {CCA闲};
\draw (8,0.5) rectangle (11,1.2); \node at (9.5,0.85) {数据发送};
\node at (2.5,-0.6) {$NB=1,\ BE=4$};
\node at (6.5,-0.6) {$CW=1$}; \node at (8,-1.1) {$CW=0$};
\end{tikzpicture}
\caption{随机退避与CCA交替的原创逻辑时序。方框表达过程顺序，CCA方框宽度不作为实际持续时间参数。}
\label{fig:p4-cca}
\end{figure}

\subsection{两类重试预算与半双工时序}

内层的接入重试预算与外层的发送重试预算是两套独立约束。节点在内层遇到忙信道，会增加$NB$，却没有因此增加一次已经发出的数据帧。节点真正发出数据而未收到ACK，才改变外层的发送失败状态。只要外层还允许再次尝试，就重新进入内层，并重新把$NB$初始化为0；新一轮接入重新拥有其规定的退避尝试预算。

因此，一个待发帧可能经历“若干次CCA忙---获准发送---没有ACK---重新进行若干次CCA”这样的过程。分析时必须分别记录“接入了几次”和“真正发送了几次”。最大重试参数还应明确采用“初次尝试之外的重试次数”还是“包含初次尝试的总次数”，不能把两种计数约定直接混用。

\subsubsection{DATA、SIFS、ACK及回复后的恢复}

发送获准之后，数据接收端正确收到帧，经过短帧间间隔（Short Interframe Space, SIFS）再回送ACK。单收发器需要由接收状态转换为发送状态，不能瞬时同时承担两种工作。ACK完成之后，接收端还需从发送切回接收，才能处理下一帧。课程的简化交换时序把两个转换保护阶段都表示为SIFS：
\begin{equation}
 \mathrm{DATA}\;\longrightarrow\;\mathrm{SIFS}\;
 \longrightarrow\;\mathrm{ACK}\;\longrightarrow\;
 \mathrm{SIFS}\;\longrightarrow\;\text{下一次可接收阶段}.
\end{equation}
其核心含义是“没有发射能量”与“接收机已准备好接收”不完全等价。若另一节点恰在回复后的切换阶段开始发送，即使侦听到的无线能量很低，接收端也可能来不及接收新帧开头。

如果只核算这一个简化交换中的占用与保护时间，可以写为
\begin{equation}
 T_{\mathrm{exchange}}=T_{\mathrm{DATA}}+
 2T_{\mathrm{SIFS}}+T_{\mathrm{ACK}}.
\end{equation}
完整交付时间还应加入随机退避、CCA、对齐等待、排队及重传，不能把上式当作从上层产生数据到成功交付的全部时延。实际实现中的帧间隔、收发转换和CCA时序须按对应PHY与标准版本配置；该式服务于本课程的简化半双工分析。

\subsection{无信标模式与非时隙CSMA/CA}

无信标模式（Non-beacon-enabled Mode）不依靠周期信标组织上述同步超帧，采用非时隙CSMA/CA（Unslotted CSMA/CA）。外层仍然是接入许可、数据发送、ACK与失败后的重试；同样存在“没有取得信道”与“已经发送但未获确认”两类失败。

内层的主要简化是不必等候全网共享的退避边界，也不使用时隙方案中初始化为2的$CW$。节点初始化$NB$和$BE$，抽取随机退避并等待；随机等待结束后执行CCA。若评估空闲，返回允许发送；若评估忙，则更新$NB$和$BE$，在仍有预算时重新退避，否则返回接入失败。

“非时隙”不表示没有定时单位，也不表示可以忽略接收端切换时间。退避仍以规定的时间尺度形成随机等待，数据与ACK交换仍需满足相应间隔。与时隙方案一样，退避计数归零本身只是进入侦听步骤，是否发送仍取决于CCA结果。

\subsubsection{非时隙接入的补充数值例：从忙信道到重传成功}

非时隙接入的计算需要同时维护绝对时刻、内层状态与外层状态。以下为依据本节算法构造的补充例题，参数专为完整演示而设，并非原课堂数值题，也不代表标准默认参数。设一帧在$t=0$进入队列，采用
\begin{align}
 T_{\mathrm{BP}}&=100\,\mu\mathrm{s}, &
 T_{\mathrm{CCA}}&=20\,\mu\mathrm{s},\\
 T_{\mathrm{DATA}}&=200\,\mu\mathrm{s}, &
 T_{\mathrm{ACK}}&=40\,\mu\mathrm{s},\\
 BE_{\min}&=2, & BE_{\max}&=3.
\end{align}
取课堂简化模型的$T_{\mathrm{SIFS}}=192\,\mu\mathrm{s}$，分别保护DATA后与ACK后的收发转换。ACK超时从数据发送结束时开始计时，期限取$500\,\mu\mathrm{s}$。该期限长于正常回复所需的$192+40=232\,\mu\mathrm{s}$。

为消除“最大次数”不同命名带来的歧义，本例直接定义阈值：内层每次CCA忙后使$NB$加1，若$NB<3$则重试，达到3则接入失败；外层每次发送后ACK超时使$NF$加1，若$NF<2$则重试，达到2则本帧失败。初始$NB=NF=0$。这一判定方式按课件流程图的严格小于号组织，不据此推定其他标准版本同名参数的含义。

已知发送端第一次CCA判忙，之后列出的CCA均判空闲；第一次数据没有被目的端成功接收，因此不回ACK，第二次数据正确接收。第一次失败原因在本例中作为给定条件，不从“发送端CCA空闲”反推接收端一定没有干扰。

第一步，初始化$NB=0,BE=2$，从$\{0,1,2,3\}$抽到$U=2$。随机等待为$200\,\mu\mathrm{s}$，故在区间$[200,220)$执行CCA。即使退避期间外部信号已经出现，随机等待也继续流逝，不能采用Wi-Fi冻结规则延长这200微秒。到$t=220$得到忙结果，更新为$NB=1,BE=3$。因为$1<3$，继续内层；此时尚未发送数据，故$NF$仍为0。

第二步，从扩大后的$\{0,\ldots,7\}$抽到$U=3$，等待300微秒，故下一次CCA位于$[520,540)$。CCA空闲，数据在$[540,740)$发送。数据没有被正确接收，节点等待至
\begin{equation}
 t_{\mathrm{timeout}}=740+500=1240\,\mu\mathrm{s}
\end{equation}
仍未获得ACK，于是$NF=1$。由于$1<2$，允许再次发送；这时要重新进入内层，设置$NB=0,BE=2$。第一次内层积累的$NB=1$与$BE=3$不能延续到这次重新初始化之后。

第三步，新一轮抽到$U=1$，等待$100\,\mu\mathrm{s}$，在$[1340,1360)$执行CCA。空闲评估通过后，第二次数据在$[1360,1560)$发送。目的端正确收到数据，先切换$192\,\mu\mathrm{s}$，因此ACK开始于
\begin{equation}
 t_{\mathrm{ACK,start}}=1560+192=1752\,\mu\mathrm{s},
\end{equation}
并于$1752+40=1792\,\mu\mathrm{s}$结束。ACK及时到达，外层以成功结束；发送端从$t=0$到收到完整确认共经过1792微秒。

目的端的状态还需要继续向后画：在ACK结束后，它经过192微秒切回接收，至
\begin{equation}
 t_{\mathrm{ready}}=1792+192=1984\,\mu\mathrm{s}
\end{equation}
才按本例模型准备好接收下一帧。因此，“本帧已获确认”的1792微秒与“汇聚节点可接收下一帧”的1984微秒是两个不同事件。

\begin{figure}[htbp]
\centering
\begin{tikzpicture}[every node/.style={font=\small}]
\draw[->] (0,0)--(11.4,0) node[right]{$t$};
\draw (0,0.4) rectangle (2.1,1.2); \node[align=center] at (1.05,0.8) {退避、忙CCA\\再退避};
\draw (2.1,0.4) rectangle (4,1.2); \node[align=center] at (3.05,0.8) {CCA闲\\DATA失败};
\draw (4,0.4) rectangle (6.1,1.2); \node[align=center] at (5.05,0.8) {等待ACK\\超时};
\draw (6.1,0.4) rectangle (8.5,1.2); \node[align=center] at (7.3,0.8) {内层重置\\再接入与DATA};
\draw (8.5,0.4) rectangle (10,1.2); \node[align=center] at (9.25,0.8) {SIFS\\ACK};
\draw[dashed] (10,0.4) rectangle (11.2,1.2); \node at (10.6,0.8) {SIFS};
\foreach \x/\v in {0/0,4/740,6.1/1240,8.5/1560,10/1792,11.2/1984}{\draw (\x,-0.1)--(\x,0.2); \node[rotate=45,anchor=east] at (\x,-0.2) {$\v$};}
\node at (5.6,-1.25) {时间标签单位：微秒；图形宽度不按比例};
\end{tikzpicture}
\caption{补充非时隙例题的原创过程图。接入忙、数据失败和接收端恢复被分开表示。}
\label{fig:p4-unslotted-example}
\end{figure}

这个结果还可以按时间分量核对。第一次发送前耗时540微秒，第一次数据及超时耗时$200+500=700$微秒；第二轮退避、CCA、DATA、SIFS与ACK耗时
\begin{equation}
 100+20+200+192+40=552\,\mu\mathrm{s}.
\end{equation}
合计$540+700+552=1792\,\mu\mathrm{s}$，与绝对时间轴一致。成功之后的192微秒恢复时间没有计入发送端获得确认的时延，却必须计入目的端再次可接收的时刻。

\subsubsection{同一规则下的两种丢包分支}

如果内层连续三次CCA忙，在本例阈值下，$NB$依次为1、2、3；第三次忙后退出为接入失败。假设初始退避抽样依次为$U=2,3,0$，则三次CCA分别位于$[200,220)$、$[520,540)$、$[540,560)$，在$t=560$结束为接入失败。第三次抽到0不免除CCA；$BE$已达到上限3，不再继续增大。这个分支中数据发送次数为零，$NF$仍为0，外层却仍然结束为失败。

另一种情况是每次内层都成功取得信道，但连续两次实际发送后都ACK超时。第一次超时后$NF=1$，允许新一轮接入；第二次超时后$NF=2$，达到本例阈值，停止重传。此时最后一轮的$NB$即使为0，也不能抵消外层已经耗尽的发送预算。

若另一节点在成功例子的$t=1800\,\mu\mathrm{s}$开始发送给同一个汇聚节点，按本例简化模型它落在ACK后的恢复区间$[1792,1984)$中。接收端尚未恢复，因此新帧可能丢失。这个分支说明，数值题必须同时绘制发送端的接入状态和接收端的收发状态，不能只看一条代表无线能量的忙闲线。

\subsection{返回DCF：统一原则与不同实现}

\subsubsection{CSMA/CA不是某一套固定流程图}

载波侦听多路访问与碰撞避免（CSMA/CA）是一类接入思想：发送前评估介质状态，在不能安全接入时推迟，并通过随机化等机制减少同时发送。DCF是Wi-Fi中这一思想的一种具体实现，IEEE 802.15.4又给出另一套具体规则。说“DCF加退避”容易造成误解，因为退避本来就是DCF的一部分。

Wi-Fi的退避通常在允许递减的空闲时隙下降，信道忙时冻结，待满足恢复条件后继续剩余值。这里讨论的802.15.4基本流程则先经过随机退避，再执行CCA，遇忙重新抽取；两者不能只因都出现“backoff”就照搬计数规则。Wi-Fi的$CW$描述抽样窗口，时隙802.15.4的$CW$描述剩余空闲评估次数，也不能互相代入公式。

\subsubsection{NAV、冻结和推迟分别表达什么}

NAV是一种虚拟载波侦听（Virtual Carrier Sensing）状态，根据正确收到的Duration信息保留预期交换时间。物理载波侦听（Physical Carrier Sensing）则观察无线介质是否存在被判定为忙的活动。对DCF的接入判断，可以把两者关系概括为
\begin{equation}
 \mathrm{Busy}=\mathrm{PHYBusy}\ \lor\ (NAV>0).
\end{equation}
因此，NAV为零不保证可以递减退避；只要物理层判忙，仍然需要冻结。反之，即使当下没有检测到发射能量，NAV尚未到期，也不能把这段时间当成自由竞争机会。

冻结（Freeze）指保存退避计数器的当前值，暂不递减；推迟（Defer）是发送没有立即发生的行为结果，不是另一个与NAV并列的硬件计数器。设置NAV可能导致冻结，但冻结并不要求必定先设置NAV。例如隐藏站点听到AP发送ACK时，即便没有建立新的正NAV，也能因物理忙而冻结退避。

恢复退避一般仍需满足相应帧间空闲条件，不能在忙信号刚结束的瞬间就任意恢复时隙递减。是否需要退避，又由待发帧遇忙、失败后重传、成功后的后退避等状态决定。单独出现一段等待时间，不足以判断节点必然在其后新抽一个退避数。

\subsubsection{EIFS保护未知交换，RTS门限依据帧长}

EIFS适用于不能正确解码已接收活动、因而无法可靠获取交换持续时间的保护情形。两个站点同时发送，旁观站点可能只看到物理忙而无法读取首部；忙结束后，它用EIFS保护潜在的后续回复。如果等待期间正确接收到新的帧，节点可依据新的有效信息回到正常的接入规则。

不同包长的碰撞尤其能显示发送端之间的差异：短帧发送端结束后只听到长帧尾部，可能无法正确解码，从而采用EIFS；长帧发送端则主要依据自身ACK超时判断失败。EIFS在适用时替代通常的帧间保护等待，不应机械地再叠加一个完整DIFS；是否还需执行退避，则取决于节点已有的DCF状态。

请求发送/清除发送（Request to Send/Clear to Send, RTS/CTS）通过短控制交换帮助预留介质，能够减轻某些隐藏节点造成的长帧碰撞代价。在基于门限的选择中，是否使用RTS/CTS由待发帧长度与RTS门限比较决定，不能简单写成“发现隐藏节点就一定使用RTS”。保护机制的收益还需与控制报文开销比较。这一结论把本部分的传感器接入过程与前面的Wi-Fi机制重新连成一条主线：相同的共享介质问题，可以因设备能力、能量目标和业务规模不同而采用不同具体规则。

\subsection*{本部分资料说明}
\addcontentsline{toc}{subsection}{本部分资料说明}
本部分按10月14日上午、10月14日后半段、10月15日两段课程和10月16日上午的授课顺序组织；其中10月15日开头的隐藏节点解答与前一天题设合并，以保持推导连续。图均为根据概念与过程重新设计的教学示意，不作为原始课堂板书复原。隐藏节点例题采用已明确给出的归一化时隙关系；具体标准的边界参数与版本差异不由课堂简化模型替代。电网监测仅用于说明感知与控制需求，不据课堂案例判定某次实际停电的原因。
\clearpage
\section*{课程材料与阅读顺序}
\addcontentsline{toc}{section}{课程材料与阅读顺序}
本讲义以课程记录与配套幻灯片为内容范围依据，不将未讲授的算法细节推定为考试要求。公式和流程图是对相应课程内容的教学整理；涉及课堂口述不一致之处，以明确的题设、协议条件和量纲一致性为准。逐章材料定位随讲义源码一并保存。
\begin{longtable}{p{27mm}p{124mm}}
\toprule 课程日期 & 内容顺序\\\midrule
10月8日 & 课程背景、无线链路、信号与调制、传播及信道基础\\
10月9日 & 检错与重传、停止等待、介质访问与Wi-Fi基础\\
10月10日 & DCF、载波侦听、随机退避与帧交换\\
10月12日 & 退避细节、EIFS、PCF、QoS与EDCA\\
10月13日 & 确认与公平性、HCCA、高速率与聚合、DCF时序题\\
10月14日 & 隐藏节点时序题、设备间通信、网络图与多跳路由\\
10月15日 & 能量与功率、传感器网络、节能MAC与IEEE 802.15.4\\
10月16日 & 信标、超帧、时隙及非时隙CSMA/CA、协议回顾\\\bottomrule
\end{longtable}
\noindent 主要课件包括课程导论、Basic Concepts、Wi-Fi，以及后续设备间通信与传感器网络讲授中引用的课件页面。课件缺页不以猜测补齐；后续章节以现有课程记录和已核验图片为范围依据。
\subsection*{课程记录定位}
以下链接指向对应课次的课程文档，可按文档中的时间标题回到相关讲授位置。
\begin{itemize}
\item \href{https://lcn6dqn3m0yr.feishu.cn/docx/SJM4dzPKHo4upfxK4yjcLTz9nBj}{10月8日08:00：导论与网络组织}；\href{https://lcn6dqn3m0yr.feishu.cn/docx/R8JJdjfT2oKLnVxpZTTcSJO6nth}{09:10：信号与调制}；\href{https://lcn6dqn3m0yr.feishu.cn/docx/BaBFdm7AzoHPSYxulLwcWRgenBf}{10:00：频谱与传播}；\href{https://lcn6dqn3m0yr.feishu.cn/docx/ENxEdQVrvoRZ6lxgUeAcu5yBnng}{10:49：信道与编码}
\item \href{https://lcn6dqn3m0yr.feishu.cn/docx/GSiDdfhZHoJ99Ux1XducLUkgn6b}{10月9日07:59：CRC与停止等待}；\href{https://lcn6dqn3m0yr.feishu.cn/docx/PnXMd5tSLoFiWDxmYvDcS6BJnfe}{08:45：Wi-Fi基础}；\href{https://lcn6dqn3m0yr.feishu.cn/docx/JT5ydOxALoxXVSxQTN0cM40enid}{10:00：网络组织与物理层}
\item \href{https://lcn6dqn3m0yr.feishu.cn/docx/FfgDdqktToKXpyxwlCxcv7GanSb}{10月10日08:00：DCF与帧交换}；\href{https://lcn6dqn3m0yr.feishu.cn/docx/LN0Xdi9daou8cFxTeT8cSCsbnnh}{10:00：隐藏节点与退避}
\item \href{https://lcn6dqn3m0yr.feishu.cn/docx/BSkQdCu5LoOPopxZWPGcbeJhnhf}{10月12日13:31：退避细节}；\href{https://lcn6dqn3m0yr.feishu.cn/docx/EsHBd3OOCo6ODAxhlMKc6JxNn9d}{15:30：EIFS与PCF}；\href{https://lcn6dqn3m0yr.feishu.cn/docx/BEhKd4SRwo8qxexUkBjcSpTMnfE}{16:20：EDCA}
\item \href{https://lcn6dqn3m0yr.feishu.cn/docx/QX16d9bfUoakAHxlKBCcnf2qnPg}{10月13日08:03：公平性与HCCA}；\href{https://lcn6dqn3m0yr.feishu.cn/docx/KZRBdsn5hoo9eDxMz1ucm6h0ned}{10:00：聚合与时序题}
\item \href{https://lcn6dqn3m0yr.feishu.cn/docx/Z6bvdAi7wo7stVxgwETcOt25nDf}{10月14日08:01：考试说明与时序题}；\href{https://lcn6dqn3m0yr.feishu.cn/docx/TqofdOKfFoCaoixV2cfccaBEnMh}{10:00：设备间通信与多跳}
\item \href{https://lcn6dqn3m0yr.feishu.cn/docx/SVZTdsr4ZoJsP3xd8GLcV7JjnAh}{10月15日08:03：功率、能量与传感器}；\href{https://lcn6dqn3m0yr.feishu.cn/docx/Onm4dp9S1oLG5hxAtzVcrCXknGf}{10:00：低功耗网络与MAC}
\item \href{https://lcn6dqn3m0yr.feishu.cn/docx/H8UxdBbx7owjVmxCeMbcnihengg}{10月16日08:05：超帧与传感器信道接入}
\end{itemize}

\end{document}
